% Options for packages loaded elsewhere
\PassOptionsToPackage{unicode}{hyperref}
\PassOptionsToPackage{hyphens}{url}
%
\documentclass[
  11pt,
]{article}
\usepackage{lmodern}
\usepackage{amssymb,amsmath}
\usepackage{ifxetex,ifluatex}
\ifnum 0\ifxetex 1\fi\ifluatex 1\fi=0 % if pdftex
  \usepackage[T1]{fontenc}
  \usepackage[utf8]{inputenc}
  \usepackage{textcomp} % provide euro and other symbols
\else % if luatex or xetex
  \usepackage{unicode-math}
  \defaultfontfeatures{Scale=MatchLowercase}
  \defaultfontfeatures[\rmfamily]{Ligatures=TeX,Scale=1}
\fi
% Use upquote if available, for straight quotes in verbatim environments
\IfFileExists{upquote.sty}{\usepackage{upquote}}{}
\IfFileExists{microtype.sty}{% use microtype if available
  \usepackage[]{microtype}
  \UseMicrotypeSet[protrusion]{basicmath} % disable protrusion for tt fonts
}{}
\makeatletter
\@ifundefined{KOMAClassName}{% if non-KOMA class
  \IfFileExists{parskip.sty}{%
    \usepackage{parskip}
  }{% else
    \setlength{\parindent}{0pt}
    \setlength{\parskip}{6pt plus 2pt minus 1pt}}
}{% if KOMA class
  \KOMAoptions{parskip=half}}
\makeatother
\usepackage{xcolor}
\IfFileExists{xurl.sty}{\usepackage{xurl}}{} % add URL line breaks if available
\IfFileExists{bookmark.sty}{\usepackage{bookmark}}{\usepackage{hyperref}}
\hypersetup{
  pdftitle={ASUNNS Theorem Compendium, v003},
  pdfauthor={UNNS Research Collective},
  hidelinks,
  pdfcreator={LaTeX via pandoc}}
\urlstyle{same} % disable monospaced font for URLs
\usepackage[margin=25mm]{geometry}
\usepackage{longtable,booktabs}
% Correct order of tables after \paragraph or \subparagraph
\usepackage{etoolbox}
\makeatletter
\patchcmd\longtable{\par}{\if@noskipsec\mbox{}\fi\par}{}{}
\makeatother
% Allow footnotes in longtable head/foot
\IfFileExists{footnotehyper.sty}{\usepackage{footnotehyper}}{\usepackage{footnote}}
\makesavenoteenv{longtable}
\setlength{\emergencystretch}{3em} % prevent overfull lines
\providecommand{\tightlist}{%
  \setlength{\itemsep}{0pt}\setlength{\parskip}{0pt}}
\setcounter{secnumdepth}{-\maxdimen} % remove section numbering

\title{ASUNNS Theorem Compendium, v003}
\author{UNNS Research Collective}
\date{27 November 2026}

\begin{document}
\maketitle

{
\setcounter{tocdepth}{1}
\tableofcontents
}
Results of Analytic Structures of the UNNS Substrate through ANS-I v015,
ANS-II v002 and ANS-II-ND v003. Statements are copied by program from
the proof files; the proofs stay in those files. In this edition, review
attribution inside quoted statements reads UNNS Research Collective. A
statement may refer to a numbered display or section of its source.
Three statuses are kept apart: proof status (PROVED means that a written
argument exists in the releases), mathematical review, and release or
computational acceptance. No literature novelty review of these results
has been completed. Entries with status CONJECTURE are not theorems.

\hypertarget{standing-definitions}{%
\section{Standing definitions}\label{standing-definitions}}

These definitions apply to regions A to I. They are given for
orientation; where a statement needs more, the entry quotes it, and the
setting of each release is in the proof sources.

\textbf{Nesting grammar.} Letters are integers \(N\ge1\). The step
\(x\mapsto r_Nx\) with \(r_N=(N+1)^2/N\) is admissible at a positive
integer \(x\) exactly when \(N\mid x\). A word \(w=(N_1,\dots,N_k)\) has
multiplier \(Q(w)=\prod r_{N_i}\), minimal source \(K(w)\) (the least
integer from which \(w\) is admissible) and least endpoint
\(E(w)=K(w)Q(w)\). Put \(\mathcal W(n)=\{w:E(w)\mid n\}\).

\textbf{Series.} \(G_k(s)=\sum_{|w|=k}E(w)^{-s}\) and
\(\mathcal G(s,t)=\sum_kt^kG_k(s)\); \(\mathcal Z_H(s,t)\) is the
history Dirichlet series; \(B(\sigma)=\sum_Nr_N^{-\sigma}\), and
\(\sigma_c=1.5463810\ldots\) is the root of \(B=1\). For \(a>1\),
\(\nu_a(N)=r_N^{-a}/B(a)\).

\textbf{Counting.} \(N_x(X)\) counts the histories from \(x\) with
endpoint at most \(X\); \(A(X)=\#\{w:E(w)\le X\}\); \(A_k(X)\) is its
depth-\(k\) part; \(C(X)\) counts from all sources.

\textbf{Prime coordinates.} \(X_p(N)=2v_p(N+1)-v_p(N)\) and
\(S_p(j)=\sum_{i\le j}X_p(N_i)\).

Regions J and K use the extended grammar: finitely many further
multipliers with a rational power, states in finitely many torsion
classes, and admissibility meaning that integrality is preserved. Its
setting is quoted at the head of those regions.

\hypertarget{a.-generation-recurrences-composition-and-factorization}{%
\section{A. Generation: recurrences, composition and
factorization}\label{a.-generation-recurrences-composition-and-factorization}}

\hypertarget{as-001-immediate-predecessor-series}{%
\subsection{AS-001: Immediate predecessor
series}\label{as-001-immediate-predecessor-series}}

\textbf{Proof status:} PROVED. \textbf{Review:} Baseline note, not
separately reviewed.

\textbf{Review record:} None. \textbf{Release and computation:} Baseline
note.

\textbf{Source:} Initial analytic baseline, lines 7-12. \textbf{Scope:}
Selected nesting grammar. \textbf{Literature priority:} Classical
content.

\textbf{Statement as recorded:}

For the selected integer nesting grammar, T\_N(Nk)=k(N+1)\^{}2 for N,k
\textgreater= 1. The labelled in-degree D(y) counts m \textgreater= 2
with m\^{}2 dividing y.

\[\sum_{y\ge1}D(y)y^{-s}=\zeta(s)(\zeta(2s)-1),\qquad \Re s>1.\]

Proof: sum over y=km\^{}2 and use absolute convergence in the stated
half-plane.

\textbf{Depends on:} no earlier result cited. \textbf{Classical inputs:}
divisor sums.

Elementary baseline deduction; classical in content.

\hypertarget{as-002-root-series}{%
\subsection{AS-002: Root series}\label{as-002-root-series}}

\textbf{Proof status:} PROVED. \textbf{Review:} Baseline note, not
separately reviewed.

\textbf{Review record:} None. \textbf{Release and computation:} Baseline
note.

\textbf{Source:} Initial analytic baseline, lines 14-17. \textbf{Scope:}
Selected nesting grammar. \textbf{Literature priority:} Classical
content.

\textbf{Statement as recorded:}

A(y)=1 if D(y)=0 and 0 otherwise. These are exactly the squarefree
integers, including 1, hence A(y)=mu(y)\^{}2.

\[\sum_{y\ge1}A(y)y^{-s}=\frac{\zeta(s)}{\zeta(2s)},\qquad \Re s>1.\]

\textbf{Depends on:} no earlier result cited. \textbf{Classical inputs:}
squarefree Euler product.

Elementary baseline deduction; classical in content.

\hypertarget{ans-t03-source-and-endpoint-recurrences}{%
\subsection{ANS-T03: Source and endpoint
recurrences}\label{ans-t03-source-and-endpoint-recurrences}}

\textbf{Proof status:} PROVED. \textbf{Review:} UNNS Research Collective
review: valid.

\textbf{Review record:} Foundational audit, 27 November 2026
(PROVENANCE/REVIEWS). \textbf{Release and computation:} Release not
frozen: Windows acceptance pending.

\textbf{Source:} ANS-I v001, lines 13-15, 21-23. \textbf{Scope:}
Integral-image grammars (see statement). \textbf{Literature priority:}
Unreviewed.

\textbf{Release table summary} (ANS-I/v001): Recurrences for E(wb),
K(wb); predecessor recursion

\textbf{Statement as recorded:}

\textbf{Theorem.} For every word \(w\) there are integers
\(K(w),E(w)=K(w)Q(w)\) such that \(w\) is admissible exactly from the
sources \(K(w)\mathbb Z_{>0}\), with endpoints exactly
\(E(w)\mathbb Z_{>0}\). Starting from \(K=E=1\) for the empty word,
\[E(wb)=\frac{p_b\,E(w)}{\gcd(E(w),q_b)},\qquad K(wb)=\frac{q_b\,K(w)}{\gcd(E(w),q_b)} .\]
For nesting, \(E(wN)=(N+1)^2E(w)/\gcd(E(w),N)\) and
\(K(wN)=N\,K(w)/\gcd(E(w),N)\).

\textbf{Corollary (histories ending at \(n\)).} Let
\(\mathcal W(n)=\{w:E(w)\mid n\}\). Then \(\mathcal W(n)\) is the set of
words of histories ending at \(n\),
\[\mathcal W(n)=\{\varepsilon\}\ \sqcup\bigsqcup_{b:\;p_b\mid n}\mathcal W\!\left(\tfrac{nq_b}{p_b}\right)b,\]
and under (F) it is finite.

\textbf{Depends on:} no earlier result cited. \textbf{Used by:} ANS-T04,
ANS-T05, ANS-T11, ANS-T12, ANS-T13, ANS-T14, ANS-T46m, ANS-T47.

\hypertarget{ans-t04-prime-coordinate-evolution}{%
\subsection{ANS-T04: Prime-coordinate
evolution}\label{ans-t04-prime-coordinate-evolution}}

\textbf{Proof status:} PROVED. \textbf{Review:} UNNS Research Collective
review: valid. Prime-coordinate recurrence and maximum formula.

\textbf{Review record:} Foundational audit, 27 November 2026
(PROVENANCE/REVIEWS). \textbf{Release and computation:} Release not
frozen: Windows acceptance pending.

\textbf{Source:} ANS-I v001, lines 29-33. \textbf{Scope:} Selected
nesting grammar. \textbf{Literature priority:} Unreviewed.

\textbf{Release table summary} (ANS-I/v001): Prime-coordinate law,
closed (maximum) form

\textbf{Statement as recorded:}

\textbf{Theorem.} A nesting transition acts by
\(v_p(T_N(x))=v_p(x)+2v_p(N+1)-v_p(N)\). For minimal endpoints,
\(e_p(w)=v_p(E(w))\) satisfies
\[e_p(wN)=2v_p(N+1)+\max\{e_p(w)-v_p(N),\,0\},\] and, with
\(u_i=2v_p(N_i+1)\) and \(d_i=v_p(N_i)\), for \(k\ge1\)
\[e_p(N_1\cdots N_k)=\max_{1\le j\le k}\Big(\sum_{i=j}^{k}u_i-\sum_{i=j+1}^{k}d_i\Big).\]
For each \(p\) at most one of \(v_p(N)\), \(v_p(N+1)\) is non-zero.

\textbf{Depends on:} ANS-T03.

\hypertarget{ans-t47-composition-of-words-block-factorization-supermultiplicativity}{%
\subsection{ANS-T47: Composition of words, block factorization,
supermultiplicativity}\label{ans-t47-composition-of-words-block-factorization-supermultiplicativity}}

\textbf{Proof status:} PROVED. \textbf{Review:} UNNS Research Collective
review: valid.

\textbf{Review record:} Release review records. \textbf{Release and
computation:} Archived release.

\textbf{Source:} ANS-I v014, lines 23-24, 32-33, 37-41. \textbf{Scope:}
Selected nesting grammar. \textbf{Literature priority:} Underlying
structure classical; otherwise unreviewed.

\textbf{Release table summary} (ANS-I/v014): Composition law (bicyclic
per prime), gcd identity, block factorization, supermultiplicativity

\textbf{Statement as recorded:}

\textbf{Theorem.} For words \(u,v\) let \(d=\gcd(E(u),K(v))\). Then
\[K(uv)=\frac{K(u)K(v)}d,\qquad E(uv)=\frac{E(u)E(v)}d .\]

\textbf{Corollary 1.}
\(\gcd(K(uv),E(uv))=\gcd\big(K(u)K(v),E(u)E(v)\big)/d\). If moreover
\(\gcd(K(u),E(u))=\gcd(K(v),E(v))=1\), then
\[\gcd(K(uv),E(uv))=\gcd(K(u),E(v)).\]

\textbf{Corollary 2 (block factorization).} Let \(J_s=\mu*{\rm id}^s\)
and, for \(\sigma>\frac12\),
\[A_m(e;s)=\sum_{|u|=m,\;e\mid E(u)}E(u)^{-s},\qquad B_n(e;s)=\sum_{|v|=n,\;e\mid K(v)}E(v)^{-s}.\]
Then \[G_{m+n}(s)=\sum_{e\ge1}J_s(e)\,A_m(e;s)\,B_n(e;s),\] and
consequently \(G_{m+n}(\sigma)\ge G_m(\sigma)\,G_n(\sigma)\) for real
\(\sigma\).

\textbf{Depends on:} ANS-T03, ANS-T12. \textbf{Used by:} ANS-T46r,
ANS-T46s. \textbf{Classical inputs:} bicyclic monoid (named only).

\hypertarget{ans-t01-history-dirichlet-factorization}{%
\subsection{ANS-T01: History Dirichlet
factorization}\label{ans-t01-history-dirichlet-factorization}}

\textbf{Proof status:} PROVED. \textbf{Review:} UNNS Research Collective
review: valid. Absolute-convergence region as stated.

\textbf{Review record:} Foundational audit, 27 November 2026
(PROVENANCE/REVIEWS). \textbf{Release and computation:} Release not
frozen: Windows acceptance pending.

\textbf{Source:} ANS-I v001, lines 39-40, 44-46. \textbf{Scope:}
Integral-image grammars (see statement). \textbf{Literature priority:}
Unreviewed.

\textbf{Release table summary} (ANS-I/v001): History Dirichlet
factorization Z\_H = zeta G, with exact region

\textbf{Statement as recorded:}

Let \(H_k(n)=\#\{w\in\mathcal W(n):|w|=k\}\) and
\(a_k(n)=\#\{w:|w|=k,\ E(w)=n\}\), and put
\[\mathcal G(s,t)=\sum_w\frac{t^{|w|}}{E(w)^s},\qquad \mathcal Z_H(s,t)=\sum_{n,k}\frac{H_k(n)t^k}{n^s}.\]

\textbf{Theorem.} For every integral-image grammar with (F),
\[H_k(n)=\sum_{d\mid n}a_k(d),\qquad a_k(n)=\sum_{d\mid n}\mu(d)H_k(n/d),\]
and, for \(t\ne0\), \(\mathcal Z_H(s,t)\) converges absolutely iff
\(\operatorname{Re}s>1\) and
\(\mathcal G(\operatorname{Re}s,|t|)<\infty\); there
\(\mathcal Z_H(s,t)=\zeta(s)\,\mathcal G(s,t)\).

\textbf{Depends on:} no earlier result cited. \textbf{Used by:} ANS-T14.

\hypertarget{ans-t02-depth-resolved-moebius-inversion}{%
\subsection{ANS-T02: Depth-resolved Moebius
inversion}\label{ans-t02-depth-resolved-moebius-inversion}}

\textbf{Proof status:} PROVED. \textbf{Review:} UNNS Research Collective
review: valid. Exact depth-resolved Moebius inversion.

\textbf{Review record:} Foundational audit, 27 November 2026
(PROVENANCE/REVIEWS). \textbf{Release and computation:} Release not
frozen: Windows acceptance pending.

\textbf{Source:} ANS-I v001, lines 44-46. \textbf{Scope:} Integral-image
grammars (see statement). \textbf{Literature priority:} Unreviewed.

\textbf{Release table summary} (ANS-I/v001): Depth-resolved Moebius
inversion

\textbf{Statement as recorded:}

\textbf{Theorem.} For every integral-image grammar with (F),
\[H_k(n)=\sum_{d\mid n}a_k(d),\qquad a_k(n)=\sum_{d\mid n}\mu(d)H_k(n/d),\]
and, for \(t\ne0\), \(\mathcal Z_H(s,t)\) converges absolutely iff
\(\operatorname{Re}s>1\) and
\(\mathcal G(\operatorname{Re}s,|t|)<\infty\); there
\(\mathcal Z_H(s,t)=\zeta(s)\,\mathcal G(s,t)\).

\textbf{Depends on:} no earlier result cited.

\hypertarget{ans-t14-source-weighted-history-series-and-its-region}{%
\subsection{ANS-T14: Source-weighted history series and its
region}\label{ans-t14-source-weighted-history-series-and-its-region}}

\textbf{Proof status:} PROVED. \textbf{Review:} UNNS Research Collective
review: valid, with a recorded qualification. The stated convergence
regions; not a complete classification for negative source weights.

\textbf{Review record:} Foundational audit, 27 November 2026
(PROVENANCE/REVIEWS). \textbf{Release and computation:} Release not
frozen: Windows acceptance pending.

\textbf{Source:} ANS-I v001, lines 56-61, 101-105. \textbf{Scope:}
Integral-image grammars (see statement). \textbf{Literature priority:}
Unreviewed.

\textbf{Release table summary} (ANS-I/v001): Source-weighted series
F(s,u,t): all-history factorization, anti-diagonal closed form,
convergence region

\textbf{Statement as recorded:}

Let
\(\mathcal F(s,u,t)=\sum_w t^{|w|}E(w)^{-s}K(w)^{-u}=\sum_wt^{|w|}Q(w)^{-s}K(w)^{-(s+u)}\).

\textbf{Theorem.} Summing over all actual histories, given by a source
\(x\) and an admissible word \(w\) with endpoint \(y=Q(w)x\),
\[\sum_{(x,w)}\frac{t^{|w|}}{x^{u}\,y^{s}}=\zeta(s+u)\,\mathcal F(s,u,t)\]
wherever the left side converges absolutely, which for \(t\ne0\) means
\(\operatorname{Re}(s+u)>1\) and absolute convergence of \(\mathcal F\).
On the anti-diagonal, for \(\operatorname{Re}s>1\) and
\(|t|B(\operatorname{Re}s)<1\),
\[\mathcal F(s,-s,t)=\frac1{1-t\sum_b r_b^{-s}},\qquad B(\sigma)=\sum_br_b^{-\sigma}.\]

\textbf{Theorem.} Let \(t\ne0\), \(a=\operatorname{Re}s\) and
\(b=\operatorname{Re}(s+u)\).

\begin{enumerate}
\def\labelenumi{\arabic{enumi}.}
\tightlist
\item
  If \(a\le1\), or if \(a>1\) and \(|t|B(a)>1\), then
  \(\mathcal F(s,u,t)\) does not converge absolutely, for any real
  \(b\).
\item
  If \(b\ge0\), \(a>1\) and \(|t|B(a)<1\), it converges absolutely.
\item
  The series over all actual histories in T14a converges absolutely if
  \(b>1\), \(a>1\), \(|t|B(a)<1\); it does not if \(b\le1\), or
  \(a\le1\), or \(a>1\) and \(|t|B(a)>1\).
\end{enumerate}

\textbf{Depends on:} ANS-T03, ANS-T01, ANS-T07, ANS-T09.

\hypertarget{ans-t06-no-ordinary-euler-product}{%
\subsection{ANS-T06: No ordinary Euler
product}\label{ans-t06-no-ordinary-euler-product}}

\textbf{Proof status:} PROVED. \textbf{Review:} UNNS Research Collective
review: valid. Ordinary Euler-product obstruction.

\textbf{Review record:} Foundational audit, 27 November 2026
(PROVENANCE/REVIEWS). \textbf{Release and computation:} Release not
frozen: Windows acceptance pending.

\textbf{Source:} ANS-I v001, lines 161-161. \textbf{Scope:} Selected
nesting grammar. \textbf{Literature priority:} Unreviewed.

\textbf{Release table summary} (ANS-I/v001): No ordinary Euler product
for Z\_H or G at t = 1

\textbf{Statement as recorded:}

\textbf{Theorem.} Neither \(H(n)=\#\mathcal W(n)\) nor
\(a(n)=\sum_ka_k(n)\) is multiplicative; hence neither
\(\mathcal Z_H(s,1)\) nor \(\mathcal G(s,1)\) is a product over primes
of series in \(p^{-s}\).

\textbf{Depends on:} no earlier result cited.

\hypertarget{ans-t08-closed-form-for-unit-denominators}{%
\subsection{ANS-T08: Closed form for unit
denominators}\label{ans-t08-closed-form-for-unit-denominators}}

\textbf{Proof status:} PROVED. \textbf{Review:} UNNS Research Collective
review: valid. Unit-denominator grammar.

\textbf{Review record:} Foundational audit, 27 November 2026
(PROVENANCE/REVIEWS). \textbf{Release and computation:} Release not
frozen: Windows acceptance pending.

\textbf{Source:} ANS-I v001, lines 169-169. \textbf{Scope:}
Integral-image grammars (see statement). \textbf{Literature priority:}
Unreviewed.

\textbf{Release table summary} (ANS-I/v001): Closed form when all
denominators are 1

\textbf{Statement as recorded:}

\textbf{Theorem.} If every \(q_b=1\) then \(E(w)=Q(w)\) and
\(\mathcal G(s,t)=1/(1-t\sum_bp_b^{-s})\) wherever
\(|t|\sum_bp_b^{-\operatorname{Re}s}<1\). For the labels \(\{2,3\}\),
\(\mathcal G(s,t)=1/(1-t(2^{-s}+3^{-s}))\), absolutely convergent iff
\(|t|(2^{-\sigma}+3^{-\sigma})<1\), with meromorphic continuation to the
whole plane.

\textbf{Depends on:} no earlier result cited.

\hypertarget{b.-convergence-and-growth-in-depth}{%
\section{B. Convergence and growth in
depth}\label{b.-convergence-and-growth-in-depth}}

\hypertarget{ans-t07-sufficient-convergence-criterion}{%
\subsection{ANS-T07: Sufficient convergence
criterion}\label{ans-t07-sufficient-convergence-criterion}}

\textbf{Proof status:} PROVED. \textbf{Review:} UNNS Research Collective
review: valid. A sufficient convergence criterion.

\textbf{Review record:} Foundational audit, 27 November 2026
(PROVENANCE/REVIEWS). \textbf{Release and computation:} Release not
frozen: Windows acceptance pending.

\textbf{Source:} ANS-I v001, lines 73-73. \textbf{Scope:} Selected
nesting grammar. \textbf{Literature priority:} Unreviewed.

\textbf{Release table summary} (ANS-I/v001): Sufficient convergence
criterion, B(2) \textless{} 1

\textbf{Statement as recorded:}

\textbf{Theorem.} If \(\operatorname{Re}s=a>1\),
\(\operatorname{Re}(s+u)=b\ge0\) and \(|t|B(a)<1\), then
\(\mathcal F(s,u,t)\) converges absolutely. Also
\(B(2)=\zeta(2)-2\zeta(3)+\zeta(4)<1\).

\textbf{Depends on:} no earlier result cited. \textbf{Used by:} ANS-T14.

\textbf{Relations:} ANS-T09 extends ANS-T07.

\hypertarget{ans-t15-source-one-generation-every-prime-reached-saturation}{%
\subsection{ANS-T15: Source-one generation: every prime reached;
saturation}\label{ans-t15-source-one-generation-every-prime-reached-saturation}}

\textbf{Proof status:} PROVED. \textbf{Review:} UNNS Research Collective
review: valid.

\textbf{Review record:} Foundational audit, 27 November 2026
(PROVENANCE/REVIEWS). \textbf{Release and computation:} Release not
frozen: Windows acceptance pending.

\textbf{Source:} ANS-I v001, lines 79-79, 83-83. \textbf{Scope:}
Selected nesting grammar. \textbf{Literature priority:} Unreviewed.

\textbf{Release table summary} (ANS-I/v001): Every source's history
count grows as X\^{}(sigma\_c + o(1)); all primes reachable from 1

\textbf{Statement as recorded:}

\textbf{Lemma (every prime is reached from source 1).} For each prime
\(p\) there is a word \(w_p\) admissible from \(1\) with
\(p\mid Q(w_p)\).

\textbf{Lemma (saturation).} Fix \(\sigma\ge0\) and \(M\), and let the
letters of \(v=(v_1,v_2,\dots)\) be independent with law
\(\nu(N)=r_N^{-\sigma}/B_M(\sigma)\), \(N\le M\). There is a word
\(w_0\), admissible from \(1\), such that with probability at least
\(1/2\) the whole infinite word \(w_0v\) is admissible from \(1\).

\textbf{Depends on:} no earlier result cited. \textbf{Used by:} ANS-T09,
ANS-T10, ANS-T11, ANS-T23, ANS-T31.

\hypertarget{ans-t09-growth-in-depth-and-the-sharp-boundary-sigma_c}{%
\subsection{ANS-T09: Growth in depth and the sharp boundary
sigma\_c}\label{ans-t09-growth-in-depth-and-the-sharp-boundary-sigma_c}}

\textbf{Proof status:} PROVED. \textbf{Review:} UNNS Research Collective
review: valid.

\textbf{Review record:} Foundational audit, 27 November 2026
(PROVENANCE/REVIEWS). \textbf{Release and computation:} Release not
frozen: Windows acceptance pending.

\textbf{Source:} ANS-I v001, lines 93-93, 111-113. \textbf{Scope:}
Selected nesting grammar. \textbf{Literature priority:} Unreviewed.

\textbf{Release table summary} (ANS-I/v001): G\_k\^{}(1/k)
-\textgreater{} B(sigma); sharp abscissa sigma\_c with B(sigma\_c) = 1

\textbf{Statement as recorded:}

\textbf{Theorem.} Let \(G_k(\sigma)=\sum_{|w|=k}E(w)^{-\sigma}\). Then
\(\lim_{k\to\infty}G_k(\sigma)^{1/k}=B(\sigma)\) for \(\sigma>1\), and
the limit is \(+\infty\) for \(0\le\sigma\le1\).

\textbf{Corollary (abscissa).} \(\mathcal G(s,1)\) and
\(\mathcal Z_H(s,1)\) have abscissa of absolute convergence
\(\sigma_c\), the unique root of
\[\sum_{N\ge1}\Big(\frac{N}{(N+1)^2}\Big)^{\sigma_c}=1,\qquad\sigma_c=1.5463810\ldots\]
Because the coefficients are non-negative, \(s=\sigma_c\) is a singular
point of both series (Landau).

\textbf{Depends on:} ANS-T15. \textbf{Used by:} ANS-T11, ANS-T14,
ANS-T19. \textbf{Classical inputs:} Landau.

\textbf{Relations:} ANS-T09 extends ANS-T07.

\hypertarget{ans-t10-growth-exponent-of-every-counting-function}{%
\subsection{ANS-T10: Growth exponent of every counting
function}\label{ans-t10-growth-exponent-of-every-counting-function}}

\textbf{Proof status:} PROVED. \textbf{Review:} UNNS Research Collective
review: valid.

\textbf{Review record:} Foundational audit, 27 November 2026
(PROVENANCE/REVIEWS). \textbf{Release and computation:} Release not
frozen: Windows acceptance pending.

\textbf{Source:} ANS-I v001, lines 119-122. \textbf{Scope:} Selected
nesting grammar. \textbf{Literature priority:} Unreviewed.

\textbf{Release table summary} (ANS-I/v001): log A(X)/log X, log
C(X)/log X -\textgreater{} sigma\_c

\textbf{Statement as recorded:}

For a source \(x\) let \(N_x(X)\) be the number of histories from \(x\)
with endpoint at most \(X\). Let \(A(X)=\#\{w:E(w)\le X\}\) and
\(C(X)=\sum_{n\le X}\#\mathcal W(n)\).

\textbf{Theorem.} For every source \(x\ge1\),
\[\lim_{X\to\infty}\frac{\log N_x(X)}{\log X}=\lim\frac{\log A(X)}{\log X}=\lim\frac{\log C(X)}{\log X}=\sigma_c .\]

\textbf{Depends on:} ANS-T15. \textbf{Used by:} ANS-T37, ANS-T40.

\hypertarget{ans-t12-convergence-region-at-fixed-depth}{%
\subsection{ANS-T12: Convergence region at fixed
depth}\label{ans-t12-convergence-region-at-fixed-depth}}

\textbf{Proof status:} PROVED. \textbf{Review:} UNNS Research Collective
review: valid.

\textbf{Review record:} Foundational audit, 27 November 2026
(PROVENANCE/REVIEWS). \textbf{Release and computation:} Release not
frozen: Windows acceptance pending.

\textbf{Source:} ANS-I v001, lines 134-134. \textbf{Scope:} Selected
nesting grammar. \textbf{Literature priority:} Unreviewed.

\textbf{Release table summary} (ANS-I/v001): Fixed-depth region a + b
\textgreater{} 1; abscissa of G\_k is 1/2

\textbf{Statement as recorded:}

\textbf{Theorem.} Let \(k\ge1\), \(a>0\), \(b\ge0\) and
\(F_k(a,b)=\sum_{|w|=k}Q(w)^{-a}K(w)^{-b}\). Then \(F_k(a,b)<\infty\)
iff \(a+b>1\) when \(b>0\), and iff \(a>1\) when \(b=0\). In particular
\(G_k(\sigma)=F_k(\sigma,\sigma)\) has abscissa exactly \(1/2\) for
every \(k\ge1\).

\textbf{Depends on:} ANS-T03. \textbf{Used by:} ANS-T34, ANS-T35,
ANS-T46m, ANS-T47.

\hypertarget{ans-t11-saturation-criterion-for-general-grammars-classification-open}{%
\subsection{ANS-T11: Saturation criterion for general grammars
(classification
open)}\label{ans-t11-saturation-criterion-for-general-grammars-classification-open}}

\textbf{Proof status:} PROVED. \textbf{Review:} UNNS Research Collective
review: valid, with a recorded qualification. The saturation criterion
only; the general classification remains open.

\textbf{Review record:} Foundational audit, 27 November 2026
(PROVENANCE/REVIEWS). \textbf{Release and computation:} Release not
frozen: Windows acceptance pending.

\textbf{Source:} ANS-I v001, lines 175-177. \textbf{Scope:}
Integral-image grammars (see statement). \textbf{Literature priority:}
Unreviewed.

\textbf{Release table summary} (ANS-I/v001): Analytic classification
across grammars

\textbf{Statement as recorded:}

\textbf{Theorem.} Let \(G\) be an integral-image grammar with (F), \(S\)
a finite set of labels, \(\sigma\ge0\) and
\(B_S(\sigma)=\sum_{b\in S}r_b^{-\sigma}\). Suppose that for every prime
\(p\) dividing some \(q_b\), \(b\in S\),
\[\sum_{b\in S}r_b^{-\sigma}\big(v_p(p_b)-v_p(q_b)\big)>0 .\] Then
\(\liminf_kG_k(\sigma)^{1/k}\ge B_S(\sigma)\). If such sets exist with
\(B_S(\sigma)\) arbitrarily close to \(B(\sigma)\) (or arbitrarily large
when \(B(\sigma)=\infty\)), then \(\lim G_k(\sigma)^{1/k}=B(\sigma)\).

\textbf{Depends on:} ANS-T03, ANS-T09, ANS-T15.

\hypertarget{ans-t26-finite-alphabet-contrast}{%
\subsection{ANS-T26: Finite-alphabet
contrast}\label{ans-t26-finite-alphabet-contrast}}

\textbf{Proof status:} PROVED. \textbf{Review:} UNNS Research Collective
review: valid.

\textbf{Review record:} Release review records. \textbf{Release and
computation:} Archived release.

\textbf{Source:} ANS-I v002, lines 101-101. \textbf{Scope:} Selected
nesting grammar. \textbf{Literature priority:} Unreviewed.

\textbf{Release table summary} (ANS-I/v002): Finite alphabets: G\_k
-\textgreater{} c\_M \textgreater{} 0, blow-up c\_M/(mu\_M delta), c\_M
-\textgreater{} 0

\textbf{Statement as recorded:}

\textbf{Theorem.} Let \(M\ge2\), \(\sigma_M\) the root of \(B_M=1\),
letters \(N\le M\) with law \(r_N^{-\sigma_M}\). Then \(K_k\) stabilizes
almost surely at a finite \(K_\infty\);
\(G^{(M)}_k(\sigma_M)\downarrow c_M=\mathbb EK_\infty^{-\sigma_M}>0\);
\(\delta\,\mathcal G^{(M)}(\sigma_M+\delta,1)\to c_M/\mu_M>0\) as
\(\delta\downarrow0\); and \(c_M\to0\) as \(M\to\infty\).

\textbf{Depends on:} ANS-T23.

\hypertarget{c.-the-critical-exponent}{%
\section{C. The critical exponent}\label{c.-the-critical-exponent}}

\hypertarget{ans-t16-critical-probability-representation}{%
\subsection{ANS-T16: Critical probability
representation}\label{ans-t16-critical-probability-representation}}

\textbf{Proof status:} PROVED. \textbf{Review:} UNNS Research Collective
review: valid.

\textbf{Review record:} Foundational audit, 27 November 2026
(PROVENANCE/REVIEWS). \textbf{Release and computation:} Archived
release.

\textbf{Source:} ANS-I v002, lines 11-12. \textbf{Scope:} Selected
nesting grammar. \textbf{Literature priority:} Unreviewed.

\textbf{Release table summary} (ANS-I/v002): G\_k(a)/B(a)\^{}k =
E{[}K\_k\^{}-a{]}

\textbf{Statement as recorded:}

Since \(\Pr(N_1\cdots N_k=w)=Q(w)^{-a}/B(a)^k\) and \(E=KQ\),
\[\frac{G_k(a)}{B(a)^k}=\mathbb E\,K_k^{-a};\qquad G_k(\sigma)=\mathbb E_{\nu_c}K_k^{-\sigma}.\]

\textbf{Depends on:} no earlier result cited. \textbf{Used by:} ANS-T19.

\hypertarget{ans-t17-infinite-alphabet-admissibility-obstruction}{%
\subsection{ANS-T17: Infinite-alphabet admissibility
obstruction}\label{ans-t17-infinite-alphabet-admissibility-obstruction}}

\textbf{Proof status:} PROVED. \textbf{Review:} UNNS Research Collective
review: valid.

\textbf{Review record:} Foundational audit, 27 November 2026
(PROVENANCE/REVIEWS). \textbf{Release and computation:} Archived
release.

\textbf{Source:} ANS-I v002, lines 19-19. \textbf{Scope:} Selected
nesting grammar. \textbf{Literature priority:} Unreviewed.

\textbf{Release table summary} (ANS-I/v002): Critical random words are
a.s. inadmissible from any integer

\textbf{Statement as recorded:}

\textbf{Theorem (T17).} An infinite \(\nu_a\)-random word is almost
surely not admissible from any given integer \(x\).

\textbf{Depends on:} no earlier result cited. \textbf{Used by:} ANS-T18,
ANS-T23.

\hypertarget{ans-t18-divergence-of-minimal-sources}{%
\subsection{ANS-T18: Divergence of minimal
sources}\label{ans-t18-divergence-of-minimal-sources}}

\textbf{Proof status:} PROVED. \textbf{Review:} UNNS Research Collective
review: valid.

\textbf{Review record:} Foundational audit, 27 November 2026
(PROVENANCE/REVIEWS). \textbf{Release and computation:} Archived
release.

\textbf{Source:} ANS-I v002, lines 23-23. \textbf{Scope:} Selected
nesting grammar. \textbf{Literature priority:} Unreviewed.

\textbf{Release table summary} (ANS-I/v002): K\_k -\textgreater{}
infinity a.s.

\textbf{Statement as recorded:}

\textbf{Theorem (T18).} \(K_k\to\infty\) almost surely.

\textbf{Depends on:} ANS-T17. \textbf{Used by:} ANS-T19, ANS-T21.

\hypertarget{ans-t19-decay-of-the-critical-depth-terms}{%
\subsection{ANS-T19: Decay of the critical depth
terms}\label{ans-t19-decay-of-the-critical-depth-terms}}

\textbf{Proof status:} PROVED. \textbf{Review:} UNNS Research Collective
review: valid.

\textbf{Review record:} Foundational audit, 27 November 2026
(PROVENANCE/REVIEWS). \textbf{Release and computation:} Archived
release.

\textbf{Source:} ANS-I v002, lines 28-28. \textbf{Scope:} Selected
nesting grammar. \textbf{Literature priority:} Unreviewed.

\textbf{Release table summary} (ANS-I/v002): G\_k(sigma\_c)
-\textgreater{} 0 with G\_k\^{}(1/k) -\textgreater{} 1

\textbf{Statement as recorded:}

\textbf{T19.} \(G_k(\sigma)\to0\) by T18, T16 and bounded convergence,
while \(G_k(\sigma)^{1/k}\to B(\sigma)=1\) (v001, Theorem 4.4).

\textbf{Depends on:} ANS-T18, ANS-T16, ANS-T09. \textbf{Used by:}
ANS-T20.

\textbf{Relations:} ANS-T23 quantifies ANS-T19.

\hypertarget{ans-t20-no-meromorphic-continuation-through-sigma_c}{%
\subsection{ANS-T20: No meromorphic continuation through
sigma\_c}\label{ans-t20-no-meromorphic-continuation-through-sigma_c}}

\textbf{Proof status:} PROVED. \textbf{Review:} UNNS Research Collective
review: valid.

\textbf{Review record:} Foundational audit, 27 November 2026
(PROVENANCE/REVIEWS). \textbf{Release and computation:} Archived
release.

\textbf{Source:} ANS-I v002, lines 30-30. \textbf{Scope:} Selected
nesting grammar. \textbf{Literature priority:} Unreviewed.

\textbf{Release table summary} (ANS-I/v002): No meromorphic continuation
through sigma\_c

\textbf{Statement as recorded:}

\textbf{T20.} By (0), \(E(wN)\ge r_NE(w)\ge4E(w)\), so
\(G_k(\sigma+\delta)\le4^{-k\delta}G_k(\sigma)\). Hence
\(\delta\,\mathcal G(\sigma+\delta,1)\le\delta\sum_k4^{-k\delta}G_k(\sigma)\to0\)
since \(G_k(\sigma)\to0\). A pole of order \(m\ge1\) would make
\(\delta\,\mathcal G\) tend to a non-zero limit or to infinity; a
removable point is excluded by Landau's theorem. So \(\mathcal G(s,1)\)
has no meromorphic continuation to a neighbourhood of \(\sigma\).
Theorem T22 below strengthens this.

\textbf{Depends on:} ANS-T19. \textbf{Used by:} ANS-T31, ANS-T36.
\textbf{Classical inputs:} Landau.

\textbf{Relations:} ANS-T22 strengthens ANS-T20.

\hypertarget{ans-t21-strict-growth-suppression}{%
\subsection{ANS-T21: Strict growth
suppression}\label{ans-t21-strict-growth-suppression}}

\textbf{Proof status:} PROVED. \textbf{Review:} UNNS Research Collective
review: valid.

\textbf{Review record:} Foundational audit, 27 November 2026
(PROVENANCE/REVIEWS). \textbf{Release and computation:} Archived
release.

\textbf{Source:} ANS-I v002, lines 32-34. \textbf{Scope:} Selected
nesting grammar. \textbf{Literature priority:} Unreviewed.

\textbf{Release table summary} (ANS-I/v002): N\_x, A, C are
o(X\^{}sigma\_c)

\textbf{Statement as recorded:}

\textbf{T21.} Every \(w\) with \(E(w)\le X\) gives
\[\frac{A(X)}{X^{\sigma}}=\mathbb E_{\nu_c}\sum_{k\ge0}K_k^{-\sigma}\Big(\frac{E_k}{X}\Big)^{\sigma}\mathbf 1_{E_k\le X}.\qquad(1)\]
Since \(E_{k+1}\ge4E_k\), the sum is at most
\(C_4=(1-4^{-\sigma})^{-1}\). Fix a word with \(K_k\to\infty\). Terms
with \(k<k_0\) tend to \(0\) as \(X\to\infty\), and the others sum to at
most \(C_4K_{k_0}^{-\sigma}\). Dominated convergence gives
\(A(X)=o(X^{\sigma})\). A history from \(x\) with endpoint \(\le X\) has
\(E(w)\le X\), so \(N_x(X)\le A(X)\); and
\(C(X)=\sum_{m\le X}A(X/m)=o(X^\sigma)\) because \(\sigma>1\).
\(\square\)

\textbf{Depends on:} ANS-T18. \textbf{Used by:} ANS-T24.

\textbf{Relations:} ANS-T24 quantifies ANS-T21.

\hypertarget{ans-t22-convergence-on-the-critical-surface-smooth-singularity}{%
\subsection{ANS-T22: Convergence on the critical surface; smooth
singularity}\label{ans-t22-convergence-on-the-critical-surface-smooth-singularity}}

\textbf{Proof status:} PROVED. \textbf{Review:} UNNS Research Collective
review: valid.

\textbf{Review record:} Release review records. \textbf{Release and
computation:} Archived release.

\textbf{Source:} ANS-I v002, lines 54-54, 58-58. \textbf{Scope:}
Selected nesting grammar. \textbf{Literature priority:} Unreviewed.

\textbf{Release table summary} (ANS-I/v002): On \textbar t\textbar B(a)
= 1, a \textgreater{} 1: F converges absolutely iff Re(s+u)
\textgreater{} 0; G(sigma\_c,1) finite; G and Z\_H are C-infinity up to
Re s = sigma\_c

\textbf{Statement as recorded:}

\textbf{Corollary (T22, the whole critical surface).} Let
\(a=\operatorname{Re}s>1\), \(b=\operatorname{Re}(s+u)\) and
\(|t|B(a)=1\). Then \(\mathcal F(s,u,t)\) converges absolutely iff
\(b>0\). In particular \(\mathcal G(s,t)\) converges absolutely on the
surface; \(\mathcal G(\sigma,1)<\infty\).

\textbf{Corollary (smooth singularity).} All derivatives of
\(\mathcal G(s,1)\) and of \(\mathcal Z_H(s,1)\) converge absolutely and
uniformly on the closed half-plane \(\operatorname{Re}s\ge\sigma\). Thus
both functions extend to \(C^\infty\) functions there, although
\(s=\sigma\) is a singular point. No pole, and no algebraic or
logarithmic branch point that would make some derivative unbounded, can
occur at \(\sigma\).

\textbf{Depends on:} ANS-T23. \textbf{Used by:} ANS-T37a.
\textbf{Classical inputs:} Landau.

\textbf{Relations:} ANS-T22 strengthens ANS-T20.

\hypertarget{ans-t23-rate-of-critical-decay-up-to-logarithms}{%
\subsection{ANS-T23: Rate of critical decay up to
logarithms}\label{ans-t23-rate-of-critical-decay-up-to-logarithms}}

\textbf{Proof status:} PROVED. \textbf{Review:} UNNS Research Collective
review: valid.

\textbf{Review record:} Release review records. \textbf{Release and
computation:} Archived release.

\textbf{Source:} ANS-I v002, lines 37-37, 81-83. \textbf{Scope:}
Selected nesting grammar. \textbf{Literature priority:} Unreviewed.

\textbf{Release table summary} (ANS-I/v002): -log(G\_k(a)/B(a)\^{}k) =
k\^{}(1/a + o(1))

\textbf{Statement as recorded:}

\textbf{Theorem.} For every \(a>1\) and \(b>0\) and all \(k\ge2\),
\(\mathbb E\,K_k^{-b}\le C_b\exp(-c_b\,k^{1/a}/\log k)\).

\textbf{Theorem (T23).} For every \(a>1\) and large \(k\),
\[\frac{c\,k^{1/a}}{\log k}\ \le\ -\log\frac{G_k(a)}{B(a)^k}\ \le\ C\,k^{1/a}\log^2k .\]
Hence \(-\log\big(G_k(a)/B(a)^k\big)=k^{1/a+o(1)}\), and
\(G_k(\sigma_c)=\exp(-k^{1/\sigma_c+o(1)})\).

\textbf{Depends on:} ANS-T17, ANS-T15. \textbf{Used by:} ANS-T22,
ANS-T23a, ANS-T24, ANS-T26, ANS-T28a, ANS-T48, ANS2-T02, ANS2-T09a.
\textbf{Classical inputs:} Ville, Chebyshev.

\textbf{Relations:} ANS-T23a sharpens ANS-T23; ANS-T23 quantifies
ANS-T19.

\hypertarget{ans-t23a-strengthened-critical-depth-bounds}{%
\subsection{ANS-T23a: Strengthened critical depth
bounds}\label{ans-t23a-strengthened-critical-depth-bounds}}

\textbf{Proof status:} PROVED. \textbf{Review:} UNNS Research Collective
review: valid.

\textbf{Review record:} Foundational audit, 27 November 2026
(PROVENANCE/REVIEWS). \textbf{Release and computation:} Archived
release.

\textbf{Source:} ANS-I v003, lines 18-20. \textbf{Scope:} Selected
nesting grammar. \textbf{Literature priority:} Unreviewed.

\textbf{Release table summary} (ANS-I/v003): -log E{[}K\_k\^{}-b{]}
\textless= C k\^{}beta (with v002: \textgreater= c k\^{}beta / log k)

\textbf{Statement as recorded:}

\textbf{Theorem.} For \(a>1\), \(b>0\) and large \(k\),
\[\frac{c\,k^{\beta}}{\log k}\le-\log\mathbb E K_k^{-b}\le C\,k^{\beta}.\]
In particular
\(e^{-Ck^{\beta_c}}\le G_k(\sigma_c)\le e^{-ck^{\beta_c}/\log k}\).

\textbf{Depends on:} ANS-T23. \textbf{Used by:} ANS-T24a, ANS-T28a,
ANS-T37b, ANS-T48. \textbf{Classical inputs:} Jensen.

\textbf{Relations:} ANS-T23a sharpens ANS-T23.

\hypertarget{ans-t24-growth-suppression-law-up-to-logarithms}{%
\subsection{ANS-T24: Growth suppression law up to
logarithms}\label{ans-t24-growth-suppression-law-up-to-logarithms}}

\textbf{Proof status:} PROVED. \textbf{Review:} UNNS Research Collective
review: valid.

\textbf{Review record:} Release review records. \textbf{Release and
computation:} Archived release.

\textbf{Source:} ANS-I v002, lines 90-92. \textbf{Scope:} Selected
nesting grammar. \textbf{Literature priority:} Unreviewed.

\textbf{Release table summary} (ANS-I/v002): N\_x, A, C = X\^{}sigma\_c
exp(-(log X)\^{}(1/sigma\_c + o(1)))

\textbf{Statement as recorded:}

\textbf{Theorem.} With \(\sigma=\sigma_c\), for each fixed source \(x\)
and \(Z(X)\) any of \(N_x(X)\), \(A(X)\), \(C(X)\), for large \(X\),
\[\frac{c\,(\log X)^{1/\sigma}}{\log\log X}\ \le\ -\log\frac{Z(X)}{X^{\sigma}}\ \le\ C\,(\log X)^{1/\sigma}(\log\log X)^2 .\]
(\(C\) may depend on \(x\).) Hence
\(Z(X)=X^{\sigma}\exp\big(-(\log X)^{1/\sigma+o(1)}\big)\).

\textbf{Depends on:} ANS-T21, ANS-T23. \textbf{Used by:} ANS-T24a.
\textbf{Classical inputs:} Chebyshev.

\textbf{Relations:} ANS-T24a sharpens ANS-T24; ANS-T24 quantifies
ANS-T21.

\hypertarget{ans-t27-critical-derivative-growth-infimal-gevrey-order-sigma_c}{%
\subsection{ANS-T27: Critical derivative growth; infimal Gevrey order
sigma\_c}\label{ans-t27-critical-derivative-growth-infimal-gevrey-order-sigma_c}}

\textbf{Proof status:} PROVED. \textbf{Review:} UNNS Research Collective
review: valid, with a recorded qualification. Valid for the infimal
Gevrey order; membership at exactly that order is open.

\textbf{Review record:} Release review records. \textbf{Release and
computation:} Archived release.

\textbf{Source:} ANS-I v003, lines 42-44. \textbf{Scope:} Selected
nesting grammar. \textbf{Literature priority:} Unreviewed.

\textbf{Release table summary} (ANS-I/v003): c\^{}j j\^{}(sigma j)
\textless= D\_j \textless= (C j log j)\^{}(sigma j); infimal Gevrey
order sigma\_c

\textbf{Statement as recorded:}

Let
\(D_j=(-1)^j\mathcal G^{(j)}(\sigma,1)=\sum_w(\log E(w))^jE(w)^{-\sigma}\).

\textbf{Theorem.} There are \(c,C>0\) with
\(c^j j^{\sigma j}\le D_j\le(Cj\log j)^{\sigma j}\) for large \(j\).
Hence \(\log D_j=\sigma j\log j+O(j\log\log j)\),
\(D_j=(j!)^{\sigma+o(1)}\), the Taylor series of \(\mathcal G(\cdot,1)\)
at \(\sigma\) has radius \(0\), and \(\sigma_c\) is the exact infimal
Gevrey order of \(\mathcal G(\cdot,1)\) on
\(\operatorname{Re}s\ge\sigma\) (since
\(|\mathcal G^{(j)}(s,1)|\le D_j\) there).

\textbf{Depends on:} ANS-T24a. \textbf{Used by:} ANS-T37a, ANS2-T04.

\hypertarget{ans-t28a-exact-order-of-the-first-introduction-event}{%
\subsection{ANS-T28a: Exact order of the first-introduction
event}\label{ans-t28a-exact-order-of-the-first-introduction-event}}

\textbf{Proof status:} PROVED. \textbf{Review:} UNNS Research Collective
review: valid. Preventing harmful first introductions does not prevent
depletion of primes met earlier, so this does not settle T28.

\textbf{Review record:} Foundational audit, 27 November 2026
(PROVENANCE/REVIEWS). \textbf{Release and computation:} Archived
release.

\textbf{Source:} ANS-I v003, lines 53-56. \textbf{Scope:} Selected
nesting grammar. \textbf{Literature priority:} Unreviewed.

\textbf{Release table summary} (ANS-I/v003): Probability that every new
prime arrives through N+1 is exp(-order k\^{}beta / log k)

\textbf{Statement as recorded:}

A prime \emph{occurs} in \(N\) if it divides \(N(N+1)\). Let \(S_i\) be
the set of primes not occurring in \(N_1,\dots,N_{i-1}\), and let
\(\mathcal I_k\) be the event that no \(N_i\) (\(i\le k\)) is divisible
by a prime of \(S_i\): every prime met in the first \(k\) letters is
introduced through \(N+1\) only.

\textbf{Theorem.} For \(a>1\),
\[e^{-Ck^{\beta}/\log k}\le\Pr(\mathcal I_k)\le e^{-ck^{\beta}/\log k}.\]

\textbf{Depends on:} ANS-T23, ANS-T23a. \textbf{Used by:} ANS2-T07.

\hypertarget{ans-t31-holomorphic-continuation-across-the-critical-line-away-from-sigma_c}{%
\subsection{ANS-T31: Holomorphic continuation across the critical line
away from
sigma\_c}\label{ans-t31-holomorphic-continuation-across-the-critical-line-away-from-sigma_c}}

\textbf{Proof status:} PROVED. \textbf{Review:} UNNS Research Collective
review: valid.

\textbf{Review record:} Release review records. \textbf{Release and
computation:} Archived release.

\textbf{Source:} ANS-I v005, lines 21-27. \textbf{Scope:} Selected
nesting grammar. \textbf{Literature priority:} Unreviewed.

\textbf{Release table summary} (ANS-I/v005): G(s,1) and Z\_H(s,1)
continue holomorphically to \{sigma \textgreater{} sigma\_c - a
min(tau\^{}2, \textbar tau\textbar\^{}-A), tau != 0\}; sigma\_c is the
only singular point on the critical line (resolves the line part of
ANS-T29)

\textbf{Statement as recorded:}

\textbf{Theorem.} There are constants \(a>0\), \(A>0\) such that
\(\sum_kG_k(s)\) converges locally uniformly in
\[\Omega^*=\{\sigma+i\tau:\ \tau\ne0,\ \sigma>\sigma_c-a\min(\tau^2,|\tau|^{-A})\}\]
(for \(\sigma>\sigma_c\) this is the known absolute convergence, which
holds for every \(\tau\)). Consequently \(\mathcal G(s,1)\) has a
holomorphic continuation to
\(\{\operatorname{Re}s>\sigma_c\}\cup\Omega^*\), and so has
\(\mathcal Z_H(s,1)=\zeta(s)\,\mathcal G(s,1)\). On \(\Omega^*\) near
the line, \(|\mathcal G(s,1)|\le C\max(\tau^{-2},|\tau|^{A})\).

In particular \textbf{every point \(\sigma_c+i\tau\) with \(\tau\ne0\)
is a regular point}. The line \(\operatorname{Re}s=\sigma_c\) is not a
natural boundary, and the real point \(\sigma_c\), which is singular
(Landau; ANS-T20), is its only singular point.

The word series itself does not converge absolutely for
\(\operatorname{Re}s<\sigma_c\). The continuation comes from
cancellation inside each \(G_k\).

\textbf{Depends on:} ANS-T15, ANS-T20. \textbf{Used by:} ANS-T35,
ANS-T36, ANS-T37. \textbf{Classical inputs:} Landau.

\hypertarget{ans-t35-twisted-continuation-for-unitary-characters}{%
\subsection{ANS-T35: Twisted continuation for unitary
characters}\label{ans-t35-twisted-continuation-for-unitary-characters}}

\textbf{Proof status:} PROVED. \textbf{Review:} UNNS Research Collective
review: valid.

\textbf{Review record:} Release review records. \textbf{Release and
computation:} Archived release.

\textbf{Source:} ANS-I v006, lines 34-35, 43-43. \textbf{Scope:}
Selected nesting grammar. \textbf{Literature priority:} Unreviewed.

\textbf{Release table summary} (ANS-I/v006): For unitary chi not equal
to q\^{}(i tau), G\_chi continues across sigma\_c + i tau; finite order:
uniform region over the whole line

\textbf{Statement as recorded:}

A \emph{unitary character} is a homomorphism
\(\chi:\mathbb Q_{>0}\to S^1\); it is determined by arbitrary values on
the primes. Put
\[G_{\chi,k}(s)=\sum_{|w|=k}\frac{\chi(Q(w))}{E(w)^s},\qquad\mathcal G_\chi(s,t)=\sum_kt^kG_{\chi,k}(s).\]

\textbf{Theorem.} Let \(\chi\) be a unitary character and
\(\tau\in\mathbb R\). If \(\chi\) is not the character
\(q\mapsto q^{i\tau}\), then \(\mathcal G_\chi(s,t)\) continues
holomorphically to a neighbourhood of \((\sigma_c+i\tau,t_0)\) for every
\(|t_0|=1\).

\textbf{Depends on:} ANS-T31, ANS-T12. \textbf{Used by:} ANS-T36,
ANS-T37, ANS-T37b.

\hypertarget{ans-t36-finite-quotient-sector-functions-share-the-singular-part}{%
\subsection{ANS-T36: Finite-quotient sector functions share the singular
part}\label{ans-t36-finite-quotient-sector-functions-share-the-singular-part}}

\textbf{Proof status:} PROVED. \textbf{Review:} UNNS Research Collective
review: valid.

\textbf{Review record:} Release review records. \textbf{Release and
computation:} Archived release.

\textbf{Source:} ANS-I v006, lines 72-72, 76-76. \textbf{Scope:}
Selected nesting grammar. \textbf{Literature priority:} Unreviewed.

\textbf{Release table summary} (ANS-I/v006): For every finite quotient,
sector function = G/\textbar Gamma\textbar{} + function holomorphic
across the line

\textbf{Statement as recorded:}

Let \(\varphi:\mathbb Q_{>0}\to\Gamma\) be a surjective homomorphism
onto a finite abelian group, and
\(\mathcal G^{[\gamma]}(s)=\sum_{\varphi(Q(w))=\gamma}E(w)^{-s}\).

\textbf{Theorem.}
\(\mathcal G^{[\gamma]}(s)=|\Gamma|^{-1}\,\mathcal G(s,1)+R_\gamma(s)\),
where \(R_\gamma\) is holomorphic on a region
\(\{\sigma>\sigma_c-a(1+|\tau|)^{-A}\}\). Consequently each
\(\mathcal G^{[\gamma]}\) is singular at \(\sigma_c\) and regular at
every other point of the critical line.

\textbf{Depends on:} ANS-T35, ANS-T20, ANS-T31. \textbf{Used by:}
ANS-T37a.

\hypertarget{ans-t28-exact-order-of-the-critical-decay}{%
\subsection{ANS-T28: Exact order of the critical
decay}\label{ans-t28-exact-order-of-the-critical-decay}}

\textbf{Proof status:} CONJECTURE. \textbf{Review:} Not a proved result.

\textbf{Review record:} None. \textbf{Release and computation:} Archived
release.

\textbf{Source:} ANS-I v003, lines 71-73. \textbf{Scope:} Selected
nesting grammar. \textbf{Literature priority:} Unreviewed.

\textbf{Release table summary} (ANS-I/v003): Exact order: -log
E{[}K\_k\^{}-b{]} comparable to k\^{}beta / log k

\textbf{Statement as recorded:}

\textbf{Conjecture.} \(-\log\mathbb E K_k^{-b}\asymp k^{\beta}/\log k\).

Proved: it lies between \(ck^{\beta}/\log k\) and \(Ck^{\beta}\). What
is missing is control of repeated demands: under a law that keeps
\(K_k\) small, one needs the exponent of each small prime in the state
to stay supplied. Under \(\mathsf Q\) the introductions are handled by
§5, but the drift of the individual prime coordinates under
\(\mathsf Q\) is not controlled.

\textbf{Depends on:} no earlier result cited.

\hypertarget{d.-counting-equidistribution-and-typical-histories}{%
\section{D. Counting, equidistribution and typical
histories}\label{d.-counting-equidistribution-and-typical-histories}}

\hypertarget{ans-t24a-growth-suppression-without-logarithms-in-the-lower-estimate}{%
\subsection{ANS-T24a: Growth suppression without logarithms in the lower
estimate}\label{ans-t24a-growth-suppression-without-logarithms-in-the-lower-estimate}}

\textbf{Proof status:} PROVED. \textbf{Review:} UNNS Research Collective
review: valid.

\textbf{Review record:} Release review records. \textbf{Release and
computation:} Archived release.

\textbf{Source:} ANS-I v003, lines 32-33. \textbf{Scope:} Selected
nesting grammar. \textbf{Literature priority:} Unreviewed.

\textbf{Release table summary} (ANS-I/v003): -log(Z(X)/X\^{}sigma\_c)
\textless= C (log X)\^{}beta\_c, Z = N\_x, A, C

\textbf{Statement as recorded:}

\textbf{Theorem.} With \(\sigma=\sigma_c\), \(\beta=\beta_c\), any fixed
source \(x\) and \(Z(X)\) any of \(N_x(X)\), \(A(X)\), \(C(X)\):
\[\frac{c(\log X)^{\beta}}{\log\log X}\le-\log\frac{Z(X)}{X^{\sigma}}\le C(\log X)^{\beta}.\]

\textbf{Depends on:} ANS-T24, ANS-T23a. \textbf{Used by:} ANS-T24b,
ANS-T27, ANS-T42a, ANS2-T03, ANS2N-T04b.

\textbf{Relations:} ANS-T24a sharpens ANS-T24; ANS-T24b gives a second
proof of one inequality of ANS-T24a.

\hypertarget{ans-t24b-first-passage-proof-of-the-counting-estimate}{%
\subsection{ANS-T24b: First-passage proof of the counting
estimate}\label{ans-t24b-first-passage-proof-of-the-counting-estimate}}

\textbf{Proof status:} PROVED. \textbf{Review:} UNNS Research Collective
review: valid.

\textbf{Review record:} Release review records. \textbf{Release and
computation:} Release has no archive of its own.

\textbf{Source:} ANS-I v004, lines 5-5. \textbf{Scope:} Selected nesting
grammar. \textbf{Literature priority:} Unreviewed.

\textbf{Release table summary} (ANS-I/v004): First-passage proof of the
right inequality of T24a, without the central-limit window

\textbf{Statement as recorded:}

\textbf{Statement} (right inequality of ANS-T24a). For any fixed source
\(x\) and \(Z(X)\) any of \(N_x(X)\), \(A(X)\), \(C(X)\):
\(-\log\big(Z(X)/X^\sigma\big)\le C(\log X)^\beta\).

\textbf{Depends on:} ANS-T24a.

\textbf{Relations:} ANS-T24b gives a second proof of one inequality of
ANS-T24a.

\hypertarget{ans-t37a-equidistribution-of-critical-log-moments}{%
\subsection{ANS-T37a: Equidistribution of critical
log-moments}\label{ans-t37a-equidistribution-of-critical-log-moments}}

\textbf{Proof status:} PROVED. \textbf{Review:} UNNS Research Collective
review: valid.

\textbf{Review record:} Release review records. \textbf{Release and
computation:} Archived release.

\textbf{Source:} ANS-I v006, lines 84-89. \textbf{Scope:} Selected
nesting grammar. \textbf{Literature priority:} Unreviewed.

\textbf{Release table summary} (ANS-I/v006): Critical log-moments
equidistribute; liminf A\_gamma/A \textless= 1/\textbar Gamma\textbar{}
\textless= limsup (uses T27)

\textbf{Statement as recorded:}

Let \(D_j=\sum_w(\log E)^jE^{-\sigma_c}\) and let \(D_j^\gamma\) be the
part with \(\varphi(Q(w))=\gamma\).

\textbf{Theorem.} (1) \(D_j^\gamma/D_j\to1/|\Gamma|\); the error is at
most \(C^j j!/D_j\). (2) With
\(A_\gamma(X)=\#\{w:E(w)\le X,\ \varphi(Q(w))=\gamma\}\) and
\(A=\sum_\gamma A_\gamma\),
\[\liminf_{X\to\infty}\frac{A_\gamma(X)}{A(X)}\le\frac1{|\Gamma|}\le\limsup_{X\to\infty}\frac{A_\gamma(X)}{A(X)} .\]

In particular, if a sector proportion has a limit, the limit is
\(1/|\Gamma|\).

\textbf{Depends on:} ANS-T36, ANS-T22, ANS-T27.

\textbf{Relations:} ANS-T37 strengthens ANS-T37a.

\hypertarget{ans-t37-counting-equidistribution-over-finite-quotients}{%
\subsection{ANS-T37: Counting equidistribution over finite
quotients}\label{ans-t37-counting-equidistribution-over-finite-quotients}}

\textbf{Proof status:} PROVED. \textbf{Review:} UNNS Research Collective
review: valid.

\textbf{Review record:} Release review records. \textbf{Release and
computation:} Archived release.

\textbf{Source:} ANS-I v007, lines 17-22, 58-59. \textbf{Scope:}
Selected nesting grammar. \textbf{Literature priority:} Unreviewed.

\textbf{Release table summary} (ANS-I/v007): A\_gamma(X)
\textasciitilde{} A(X)/\textbar Gamma\textbar, and the same for each
fixed source, for every finite quotient

\textbf{Statement as recorded:}

Letters \(N\ge1\), \(r_N=(N+1)^2/N\), \(\sigma=\sigma_c\),
\(\nu(N)=r_N^{-\sigma}\) (a probability law, since \(B(\sigma_c)=1\)). A
word \(w\) has source \(K(w)\), multiplier \(Q(w)\) and least endpoint
\(E=KQ\). Fix a homomorphism \(\varphi\) of \(\mathbb Q_{>0}\) onto a
finite abelian group \(\Gamma\), written additively; the \emph{sector}
of \(w\) is \(\varphi(Q(w))\). Two counting problems are treated:

\begin{itemize}
\tightlist
\item
  \(H=A\): the words with \(E(w)\le X\);
\item
  \(H=N_x\), for a fixed source \(x\): the words with \(K(w)\mid x\) and
  \(xQ(w)\le X\).
\end{itemize}

\(H_\gamma(X)\) counts those in sector \(\gamma\). By ANS-T10 (v001),
\(H(X)=X^{\sigma+o(1)}\).

\textbf{Theorem.} For every finite quotient \(\Gamma\) and every
\(\gamma\in\Gamma\),
\[A_\gamma(X)\sim\frac{A(X)}{|\Gamma|},\qquad N_{x,\gamma}(X)\sim\frac{N_x(X)}{|\Gamma|}\quad(x\text{ fixed}).\]

\textbf{Depends on:} ANS-T35, ANS-T31, ANS-T10. \textbf{Used by:}
ANS-T40, ANS-T41, ANS-T42, ANS-T44.

\textbf{Relations:} ANS-T37 strengthens ANS-T37a.

\hypertarget{ans-t40-regular-variation-with-index-sigma_c}{%
\subsection{ANS-T40: Regular variation with index
sigma\_c}\label{ans-t40-regular-variation-with-index-sigma_c}}

\textbf{Proof status:} PROVED. \textbf{Review:} UNNS Research Collective
review: valid.

\textbf{Review record:} Release review records. \textbf{Release and
computation:} Archived release.

\textbf{Source:} ANS-I v007, lines 69-71. \textbf{Scope:} Selected
nesting grammar. \textbf{Literature priority:} Unreviewed.

\textbf{Release table summary} (ANS-I/v007): A(lambda X)/A(X)
-\textgreater{} lambda\^{}sigma\_c, same for N\_x: A = X\^{}sigma\_c
times a slowly varying factor

\textbf{Statement as recorded:}

\textbf{Theorem.} For every \(\lambda>0\),
\[\frac{A(\lambda X)}{A(X)}\to\lambda^{\sigma_c},\qquad\frac{N_x(\lambda X)}{N_x(X)}\to\lambda^{\sigma_c}.\]
Thus \(A(X)=X^{\sigma_c}L(X)\) with \(L\) slowly varying, and likewise
\(N_x\).

\textbf{Depends on:} ANS-T37, ANS-T10. \textbf{Used by:} ANS-T41,
ANS-T43, ANS-T44.

\hypertarget{ans-t41-asymptotic-arithmetic-replication}{%
\subsection{ANS-T41: Asymptotic arithmetic
replication}\label{ans-t41-asymptotic-arithmetic-replication}}

\textbf{Proof status:} PROVED. \textbf{Review:} UNNS Research Collective
review: valid.

\textbf{Review record:} Release review records. \textbf{Release and
computation:} Archived release.

\textbf{Source:} ANS-I v007, lines 83-83. \textbf{Scope:} Selected
nesting grammar. \textbf{Literature priority:} Unreviewed.

\textbf{Release table summary} (ANS-I/v007): C(X) \textasciitilde{}
zeta(sigma\_c) A(X), with equal sector shares

\textbf{Statement as recorded:}

\textbf{Theorem.} \(C(X)\sim\zeta(\sigma_c)\,A(X)\), and
\(C_\gamma(X)\sim\zeta(\sigma_c)\,A(X)/|\Gamma|\).

\textbf{Depends on:} ANS-T40, ANS-T37. \textbf{Used by:} ANS-T43,
ANS-T44.

\hypertarget{ans-t37b-fixed-depth-sector-equidistribution-at-sigma_c}{%
\subsection{ANS-T37b: Fixed-depth sector equidistribution at
sigma\_c}\label{ans-t37b-fixed-depth-sector-equidistribution-at-sigma_c}}

\textbf{Proof status:} PROVED. \textbf{Review:} UNNS Research Collective
review: valid.

\textbf{Review record:} Release review records. \textbf{Release and
computation:} Archived release.

\textbf{Source:} ANS-I v007, lines 91-92. \textbf{Scope:} Selected
nesting grammar. \textbf{Literature priority:} Unreviewed.

\textbf{Release table summary} (ANS-I/v007): G\_k sector ratio at
sigma\_c = 1/\textbar Gamma\textbar{} + O(e\^{}(-c'k))

\textbf{Statement as recorded:}

\textbf{Theorem.} There is \(c'>0\) with
\[\frac{G^{[\gamma]}_k(\sigma_c)}{G_k(\sigma_c)}=\frac1{|\Gamma|}+O(e^{-c'k}).\]

\textbf{Depends on:} ANS-T35, ANS-T23a.

\hypertarget{ans-t42-typical-history-laws}{%
\subsection{ANS-T42: Typical-history
laws}\label{ans-t42-typical-history-laws}}

\textbf{Proof status:} PROVED. \textbf{Review:} UNNS Research Collective
review: valid.

\textbf{Review record:} Release review records. \textbf{Release and
computation:} Archived release.

\textbf{Source:} ANS-I v007, lines 100-100, 104-109. \textbf{Scope:}
Selected nesting grammar. \textbf{Literature priority:} Unreviewed.

\textbf{Release table summary} (ANS-I/v007): Typical histories: log Q
\textasciitilde{} log X, log K = o(log X), depth \textasciitilde{} log X
/ lambda\_c, block frequencies -\textgreater{} product of
r\_N\^{}(-sigma\_c)

\textbf{Statement as recorded:}

\textbf{Theorem.} Let \(\mathcal E\) be a set of words whose part of
depth \(k\) has \(\nu\)-probability at most \(Ce^{-bk}\). Then the
number of words of \(\mathcal E\) counted by \(H(X)\) is
\(O(X^{\sigma-\varepsilon})=o(H(X))\).

\textbf{Corollary.} Choose \(w\) uniformly among the words counted by
\(A(X)\) (or by \(N_x(X)\)) and put
\(\lambda_c=\sum_Nr_N^{-\sigma}\log r_N=-B'(\sigma_c)=2.96846\ldots\).
Then, in probability as \(X\to\infty\):

\begin{enumerate}
\def\labelenumi{\arabic{enumi}.}
\tightlist
\item
  \(\log Q(w)/\log X\to1\) and \(\log E(w)/\log X\to1\);
\item
  \(\log K(w)/\log X\to0\);
\item
  \(|w|/\log X\to1/\lambda_c=0.33687\ldots\);
\item
  the frequency of every finite block \(v=(v_1,\dots,v_\ell)\) in \(w\)
  tends to \(\prod_ir_{v_i}^{-\sigma_c}\); in particular the letter
  \(N\) has asymptotic share \(r_N^{-\sigma_c}\).
\end{enumerate}

\textbf{Depends on:} ANS-T37.

\hypertarget{ans-t43-limit-law-of-normalized-endpoints}{%
\subsection{ANS-T43: Limit law of normalized
endpoints}\label{ans-t43-limit-law-of-normalized-endpoints}}

\textbf{Proof status:} PROVED. \textbf{Review:} UNNS Research Collective
review: valid.

\textbf{Review record:} Release review records. \textbf{Release and
computation:} Archived release.

\textbf{Source:} ANS-I v008, lines 110-110. \textbf{Scope:} Selected
nesting grammar. \textbf{Literature priority:} Unreviewed.

\textbf{Release table summary} (ANS-I/v008): Endpoint / X tends to the
Beta(sigma\_c, 1) law in the populations A, N\_x, C

\textbf{Statement as recorded:}

\textbf{Theorem.} Choose a history uniformly in the population counted
by \(A(X)\), \(N_x(X)\) or \(C(X)\), and let \(y\) be its endpoint
(\(E(w)\), \(xQ(w)\), or \(n\)). Then \(y/X\) converges in distribution
to the law with density \(\sigma_cu^{\sigma_c-1}\) on \((0,1)\).
Equivalently \(\log(X/y)\) tends to an exponential variable of rate
\(\sigma_c\).

\textbf{Depends on:} ANS-T40, ANS-T41.

\hypertarget{ans-t44-joint-limit-law-of-sector-replication-factor-and-endpoint}{%
\subsection{ANS-T44: Joint limit law of sector, replication factor and
endpoint}\label{ans-t44-joint-limit-law-of-sector-replication-factor-and-endpoint}}

\textbf{Proof status:} PROVED. \textbf{Review:} UNNS Research Collective
review: valid.

\textbf{Review record:} Release review records. \textbf{Release and
computation:} Archived release.

\textbf{Source:} ANS-I v008, lines 116-117. \textbf{Scope:} Selected
nesting grammar. \textbf{Literature priority:} Unreviewed.

\textbf{Release table summary} (ANS-I/v008): Sector, replication factor
and endpoint become independent: uniform, zeta, beta

\textbf{Statement as recorded:}

\textbf{Theorem.} For a finite quotient \(\Gamma\) and the sector
\(\varphi(Q(w))\): in the populations \(A\) and \(N_x\) the pair
(sector, \(y/X\)) tends to independent variables, uniform on \(\Gamma\)
and of density \(\sigma_cu^{\sigma_c-1}\). In the population \(C\),
writing the history uniquely as \((m,w)\) with source \(mK(w)\) and
endpoint \(mE(w)\),
\[\Pr\Big(m=j,\ \varphi(Q)=\gamma,\ \frac yX\le u\Big)\to\frac{j^{-\sigma_c}}{\zeta(\sigma_c)}\cdot\frac1{|\Gamma|}\cdot u^{\sigma_c}.\]

\textbf{Depends on:} ANS-T37, ANS-T40, ANS-T41.

\hypertarget{ans-t42a-concentration-bounds-for-typical-histories}{%
\subsection{ANS-T42a: Concentration bounds for typical
histories}\label{ans-t42a-concentration-bounds-for-typical-histories}}

\textbf{Proof status:} PROVED. \textbf{Review:} UNNS Research Collective
review: valid.

\textbf{Review record:} Release review records. \textbf{Release and
computation:} Archived release.

\textbf{Source:} ANS-I v008, lines 125-127. \textbf{Scope:} Selected
nesting grammar. \textbf{Literature priority:} Unreviewed.

\textbf{Release table summary} (ANS-I/v008): Concentration of multiplier
deficit, depth and block frequencies, exceptional proportion exp(-c (log
X)\^{}(1/sigma\_c))

\textbf{Statement as recorded:}

\textbf{Theorem.} Let \(H=A\) or \(N_x\), \(L=\log X\),
\(\beta=1/\sigma_c\). For every finite block \(v\) there are \(D,c>0\)
such that, outside a proportion \(e^{-cL^\beta}\) of the counted words,
\[L-\log Q(w)\le DL^\beta,\qquad\Big|\,|w|-\frac L{\lambda_c}\Big|\le DL^{(1+\beta)/2},\qquad|f_v(w)-\nu_c(v)|\le DL^{-(1-\beta)/2}.\]
Consequently \(\log K(w)\le DL^\beta\) outside the same proportion for
\(H=A\).

\textbf{Depends on:} ANS-T24a. \textbf{Classical inputs:} Bernstein.

\hypertarget{e.-fixed-depth-continuation-and-counting}{%
\section{E. Fixed depth: continuation and
counting}\label{e.-fixed-depth-continuation-and-counting}}

\hypertarget{ans-t32-fixed-depth-counting-asymptotic-and-the-constants-kappa_k}{%
\subsection{ANS-T32: Fixed-depth counting asymptotic and the constants
kappa\_k}\label{ans-t32-fixed-depth-counting-asymptotic-and-the-constants-kappa_k}}

\textbf{Proof status:} PROVED. \textbf{Review:} UNNS Research Collective
review: valid, with a recorded qualification. Valid after a correction
to the limiting step, which has been applied.

\textbf{Review record:} Release review records. \textbf{Release and
computation:} Archived release.

\textbf{Source:} ANS-I v005, lines 90-94. \textbf{Scope:} Selected
nesting grammar. \textbf{Literature priority:} Unreviewed.

\textbf{Release table summary} (ANS-I/v005): A\_k(X) \textasciitilde{}
kappa\_k sqrt(X) log\^{}(k-1) X / (2\^{}(k-1) (k-1)!), kappa\_k finite
with Euler product

\textbf{Statement as recorded:}

Let \(A_k(X)=\#\{w:|w|=k,\ E(w)\le X\}\). Let \(N_1,\dots,N_k\) be
independent and Haar-uniform in \(\hat{\mathbb Z}\); define
\(e_{p,0}=0\), \(e_{p,i}=2v_p(N_i+1)+\max\{e_{p,i-1}-v_p(N_i),0\}\), and
\(g_i=\prod_pp^{\min(e_{p,i-1},\,v_p(N_i))}\). Put
\[\kappa_k=\mathbb E\,\sqrt{g_2g_3\cdots g_k}=\prod_pc_p(k),\qquad c_p(k)=\mathbb E\,p^{\frac12\sum_{i\le k}\min(e_{p,i-1},\,v_p(N_i))}.\]

\textbf{Theorem.} For every \(k\ge1\), \(\kappa_k\) is finite and
\[A_k(X)\sim\frac{\kappa_k}{2^{k-1}(k-1)!}\,\sqrt X\,(\log X)^{k-1}.\qquad(5)\]

\textbf{Depends on:} no earlier result cited. \textbf{Used by:} ANS-C1,
ANS-T32a, ANS-T34, ANS-T34a.

\textbf{Relations:} ANS-T34a refines ANS-T32.

\hypertarget{ans-c1-real-limit-at-one-half-v001-conjecture-now-corollary}{%
\subsection{ANS-C1: Real limit at one half (v001 conjecture, now
corollary)}\label{ans-c1-real-limit-at-one-half-v001-conjecture-now-corollary}}

\textbf{Proof status:} PROVED. \textbf{Review:} UNNS Research Collective
review: valid.

\textbf{Review record:} Release review records. \textbf{Release and
computation:} Archived release.

\textbf{Source:} ANS-I v005, lines 96-96. \textbf{Scope:} Selected
nesting grammar. \textbf{Literature priority:} Unreviewed.

\textbf{Release table summary} (ANS-I/v005): (2 sigma - 1)\^{}k
G\_k(sigma) -\textgreater{} kappa\_k (conjectured in v001; corollary of
T32)

\textbf{Statement as recorded:}

\textbf{Corollary (ANS-C1, now proved).}
\((2\sigma-1)^kG_k(\sigma)\to\kappa_k\) as \(\sigma\downarrow\tfrac12\).

\textbf{Depends on:} ANS-T32.

\hypertarget{ans-t32a-local-coupling-asymptotic-of-c_pk}{%
\subsection{ANS-T32a: Local coupling asymptotic of
c\_p(k)}\label{ans-t32a-local-coupling-asymptotic-of-c_pk}}

\textbf{Proof status:} PROVED. \textbf{Review:} UNNS Research Collective
review: valid.

\textbf{Review record:} Release review records. \textbf{Release and
computation:} Archived release.

\textbf{Source:} ANS-I v006, lines 15-18, 28-28. \textbf{Scope:}
Selected nesting grammar. \textbf{Literature priority:} Unreviewed.

\textbf{Release table summary} (ANS-I/v006): c\_p(k) = 1 + C(k,2)
p\^{}-3/2 + C(k,3) p\^{}-2 + O\_k(p\^{}-5/2)

\textbf{Statement as recorded:}

\textbf{Theorem.} For every fixed \(k\ge1\), as \(p\to\infty\),
\[c_p(k)=1+\binom k2p^{-3/2}+\binom k3p^{-2}+O_k(p^{-5/2}).\]

The first term and the method are those of the review note; the
\(p^{-2}\) term is added here.

\textbf{Corollary.}
\(\displaystyle\log\kappa_k=\sum_{p\le P}\log c_p(k)+\binom k2\sum_{p>P}p^{-3/2}+\binom k3\sum_{p>P}p^{-2}+O_k(P^{-3/2}).\)

\textbf{Depends on:} ANS-T32.

\hypertarget{ans-t34-every-g_k-is-meromorphic-on-re-s-13}{%
\subsection{ANS-T34: Every G\_k is meromorphic on Re s \textgreater{}
1/3}\label{ans-t34-every-g_k-is-meromorphic-on-re-s-13}}

\textbf{Proof status:} PROVED. \textbf{Review:} UNNS Research Collective
review: valid.

\textbf{Review record:} Release review records. \textbf{Release and
computation:} Archived release.

\textbf{Source:} ANS-I v008, lines 15-21, 68-68. \textbf{Scope:}
Selected nesting grammar. \textbf{Literature priority:} Unreviewed.

\textbf{Release table summary} (ANS-I/v008): Every G\_k (and every
source-restricted G\_k(s;D)) is meromorphic on Re s \textgreater{} 1/3;
only pole at 1/2, of order k (exactness for D \textgreater{} 1: v009,
T34b)

\textbf{Statement as recorded:}

For \(D\ge1\) let \[G_k(s;D)=\sum_{|w|=k,\ D\mid K(w)}E(w)^{-s},\] so
that \(G_k(s;1)=G_k(s)\) and \(G_0(s;D)=[D=1]\). The sum converges
absolutely for \(\operatorname{Re}s>1/2\) (T12). Put
\(\sigma=\operatorname{Re}s\), \(\mu=\min(2\sigma,1)\),
\(H_k(s;D)=(2s-1)^kG_k(s;D)\),
\(J_s(d)=\sum_{e\mid d}\mu_{\rm M}(e)(d/e)^s\) (here \(\mu_{\rm M}\) is
the Möbius function) and \(\rho(d)=\prod_pp^{\lceil v_p(d)/2\rceil}\),
so that \(d\mid a^2\) exactly when \(\rho(d)\mid a\).

\textbf{Theorem.} For every \(k\ge0\) and \(D\ge1\), \(H_k(\cdot\,;D)\)
is holomorphic on \(1/3<\sigma<3/4\). On every closed substrip
\(\frac13+\delta\le\sigma\le\frac34-\delta\) and for every
\(\varepsilon>0\),
\[|H_k(s;D)|\le M\,(1+|\operatorname{Im}s|)^{2k}\,D^{-\mu+\varepsilon},\]
with \(M=M(k,\delta,\varepsilon)\).

\textbf{Corollary 1.} For every \(k\ge1\) the function \(G_k\) is
meromorphic on \(\operatorname{Re}s>1/3\). Its only singularity is a
pole at \(s=1/2\) of order exactly \(k\), with
\((2s-1)^kG_k(s)\to\kappa_k\). In particular every other point of the
line \(\operatorname{Re}s=1/2\) is regular.

\textbf{Depends on:} ANS-T12, ANS-T32. \textbf{Used by:} ANS-T34a,
ANS-T34b, ANS-T45, ANS-T46k, ANS-T46m.

\textbf{Relations:} ANS-T34b completes the exactness of ANS-T34.

\hypertarget{ans-t34a-full-counting-polynomial-with-power-saving}{%
\subsection{ANS-T34a: Full counting polynomial with power
saving}\label{ans-t34a-full-counting-polynomial-with-power-saving}}

\textbf{Proof status:} PROVED. \textbf{Review:} UNNS Research Collective
review: valid.

\textbf{Review record:} Release review records. \textbf{Release and
computation:} Archived release.

\textbf{Source:} ANS-I v008, lines 74-75. \textbf{Scope:} Selected
nesting grammar. \textbf{Literature priority:} Unreviewed.

\textbf{Release table summary} (ANS-I/v008): A\_k(X) = sqrt(X)
P\_\{k-1\}(log X) + O(X\^{}(1/2 - 1/(6(k+2)) + eps))

\textbf{Statement as recorded:}

\textbf{Corollary 2 (ANS-T34a).} For every \(k\ge1\) there is a
polynomial \(P_{k-1}\) of degree \(k-1\), with leading coefficient
\(\kappa_k/(2^{k-1}(k-1)!)\), such that
\[A_k(X)=\sqrt X\,P_{k-1}(\log X)+O_{k,\varepsilon}\Big(X^{\frac12-\frac1{6(k+2)}+\varepsilon}\Big).\]

\textbf{Depends on:} ANS-T34, ANS-T32.

\textbf{Relations:} ANS-T34a refines ANS-T32.

\hypertarget{ans-t34b-exact-pole-order-k-at-one-half-for-every-restriction}{%
\subsection{ANS-T34b: Exact pole order k at one half for every
restriction}\label{ans-t34b-exact-pole-order-k-at-one-half-for-every-restriction}}

\textbf{Proof status:} PROVED. \textbf{Review:} UNNS Research Collective
review: valid.

\textbf{Review record:} Release review records. \textbf{Release and
computation:} Archived release.

\textbf{Source:} ANS-I v009, lines 15-15. \textbf{Scope:} Selected
nesting grammar. \textbf{Literature priority:} Unreviewed.

\textbf{Release table summary} (ANS-I/v009): Pole of G\_k(s;D) at 1/2
has order exactly k, for every D

\textbf{Statement as recorded:}

\textbf{Theorem.} For all \(k\ge1\), \(D\ge1\), the pole of \(G_k(s;D)\)
at \(1/2\) has order exactly \(k\), and
\(\lim_{s\to1/2}(2s-1)^kG_k(s;D)\ge1/D\).

\textbf{Depends on:} ANS-T34.

\textbf{Relations:} ANS-T34b completes the exactness of ANS-T34.

\hypertarget{f.-depth-two}{%
\section{F. Depth two}\label{f.-depth-two}}

\hypertarget{ans-t13-depth-two-double-pole-and-counting-asymptotic}{%
\subsection{ANS-T13: Depth-two double pole and counting
asymptotic}\label{ans-t13-depth-two-double-pole-and-counting-asymptotic}}

\textbf{Proof status:} PROVED. \textbf{Review:} UNNS Research Collective
review: valid. Uses the classical Delange theorem.

\textbf{Review record:} Foundational audit, 27 November 2026
(PROVENANCE/REVIEWS). \textbf{Release and computation:} Release not
frozen: Windows acceptance pending.

\textbf{Source:} ANS-I v001, lines 144-151. \textbf{Scope:} Selected
nesting grammar. \textbf{Literature priority:} Unreviewed.

\textbf{Release table summary} (ANS-I/v001): Depth-two double pole,
constant kappa, A\_2(X) \textasciitilde{} (kappa/2) sqrt(X) log X

\textbf{Statement as recorded:}

\textbf{Theorem.} \(G_1(s)=\zeta(2s)-1\) and
\[G_2(s)=\sum_{a,b\ge1}\left[\frac{\gcd((a+1)^2,b)}{(a+1)^2(b+1)^2}\right]^{s}.\]
With \(J_s(d)=\sum_{e\mid d}\mu(e)(d/e)^s\),
\(c(d)=\prod_pp^{\lceil v_p(d)/2\rceil}\) and the Hurwitz zeta function
\(\zeta(\cdot,\cdot)\),
\[G_2(s)=\sum_{d\ge1}\frac{J_s(d)}{d^{2s}}\,\zeta\!\left(2s,1+\tfrac1d\right)\Big(\frac{\zeta(2s)}{c(d)^{2s}}-[d=1]\Big).\]
The right side continues \(G_2\) meromorphically to
\(\operatorname{Re}s>1/3\) with a single pole, at \(s=1/2\), of order
two, and \(G_2(s)\sim\dfrac{\kappa}{4(s-1/2)^2}\), where
\[\kappa=\prod_p\Big(1+(1-p^{-1/2})\sum_{j\ge1}p^{-j/2-\lceil j/2\rceil}\Big)\approx1.8536 .\]
Consequently the number of depth-two words with minimal endpoint at most
\(X\) satisfies \[A_2(X)\sim\frac{\kappa}{2}\sqrt X\log X .\]

\textbf{Depends on:} ANS-T03. \textbf{Classical inputs:} Delange.

\hypertarget{ans-t05-depth-two-representation-and-continuation}{%
\subsection{ANS-T05: Depth-two representation and
continuation}\label{ans-t05-depth-two-representation-and-continuation}}

\textbf{Proof status:} PROVED. \textbf{Review:} UNNS Research Collective
review: valid.

\textbf{Review record:} Foundational audit, 27 November 2026
(PROVENANCE/REVIEWS). \textbf{Release and computation:} Release not
frozen: Windows acceptance pending.

\textbf{Source:} ANS-I v001, lines 144-151. \textbf{Scope:} Selected
nesting grammar. \textbf{Literature priority:} Unreviewed.

\textbf{Release table summary} (ANS-I/v001): Depth-two double sum,
Hurwitz representation, continuation to Re s \textgreater{} 1/3

\textbf{Statement as recorded:}

\textbf{Theorem.} \(G_1(s)=\zeta(2s)-1\) and
\[G_2(s)=\sum_{a,b\ge1}\left[\frac{\gcd((a+1)^2,b)}{(a+1)^2(b+1)^2}\right]^{s}.\]
With \(J_s(d)=\sum_{e\mid d}\mu(e)(d/e)^s\),
\(c(d)=\prod_pp^{\lceil v_p(d)/2\rceil}\) and the Hurwitz zeta function
\(\zeta(\cdot,\cdot)\),
\[G_2(s)=\sum_{d\ge1}\frac{J_s(d)}{d^{2s}}\,\zeta\!\left(2s,1+\tfrac1d\right)\Big(\frac{\zeta(2s)}{c(d)^{2s}}-[d=1]\Big).\]
The right side continues \(G_2\) meromorphically to
\(\operatorname{Re}s>1/3\) with a single pole, at \(s=1/2\), of order
two, and \(G_2(s)\sim\dfrac{\kappa}{4(s-1/2)^2}\), where
\[\kappa=\prod_p\Big(1+(1-p^{-1/2})\sum_{j\ge1}p^{-j/2-\lceil j/2\rceil}\Big)\approx1.8536 .\]
Consequently the number of depth-two words with minimal endpoint at most
\(X\) satisfies \[A_2(X)\sim\frac{\kappa}{2}\sqrt X\log X .\]

\textbf{Depends on:} ANS-T03. \textbf{Used by:} ANS-T45.

\textbf{Relations:} ANS-T45 continues further than ANS-T05.

\hypertarget{ans-t45-depth-two-to-re-s-15-with-a-simple-pole-at-13}{%
\subsection{ANS-T45: Depth two to Re s \textgreater{} 1/5 with a simple
pole at
1/3}\label{ans-t45-depth-two-to-re-s-15-with-a-simple-pole-at-13}}

\textbf{Proof status:} PROVED. \textbf{Review:} UNNS Research Collective
review: valid.

\textbf{Review record:} Release review records. \textbf{Release and
computation:} Archived release.

\textbf{Source:} ANS-I v008, lines 87-91. \textbf{Scope:} Selected
nesting grammar. \textbf{Literature priority:} Unreviewed.

\textbf{Release table summary} (ANS-I/v008): G\_2 meromorphic on Re s
\textgreater{} 1/5, poles only at 1/2 (double) and 1/3 (simple); second
term in smoothed form

\textbf{Statement as recorded:}

\textbf{Theorem.} On \(\sigma>0\) let
\[F_2(s)=\sum_{d\ge1}\frac{J_s(d)}{(d\,\rho(d))^{2s}}\Big(\zeta\big(2s,1+\tfrac1d\big)-\zeta(2s)\Big),\qquad H(s)=\prod_p\Big(1-\frac{p^{-5s}+p^{-6s}}{1+p^{-2s}}\Big)\quad(\sigma>\tfrac15).\]
Then
\[G_2(s)=\zeta(2s)\Big(\zeta(2s)\,\zeta(3s)\,H(s)+F_2(s)-1\Big)+1.\] So
\(G_2\) is meromorphic on \(\operatorname{Re}s>1/5\), with poles only at
\(s=1/2\) (order two) and \(s=1/3\) (simple, residue
\(\zeta(2/3)^2H(1/3)/3\ne0\)).

\textbf{Depends on:} ANS-T05, ANS-T34. \textbf{Used by:} ANS-T45a,
ANS-T45b, ANS-T46b, ANS-T46k, ANS-T46s.

\textbf{Relations:} ANS-T45 continues further than ANS-T05; ANS-T45a
extends ANS-T45.

\hypertarget{ans-t45a-depth-two-meromorphic-on-re-s-0-with-natural-boundary}{%
\subsection{ANS-T45a: Depth two: meromorphic on Re s \textgreater{} 0
with natural
boundary}\label{ans-t45a-depth-two-meromorphic-on-re-s-0-with-natural-boundary}}

\textbf{Proof status:} PROVED. \textbf{Review:} UNNS Research Collective
review: valid.

\textbf{Review record:} Release review records. \textbf{Release and
computation:} Archived release.

\textbf{Source:} ANS-I v009, lines 23-26. \textbf{Scope:} Selected
nesting grammar. \textbf{Literature priority:} Unreviewed.

\textbf{Release table summary} (ANS-I/v009): G\_2 meromorphic on Re s
\textgreater{} 0 with natural boundary Re s = 0 (uses Estermann)

\textbf{Statement as recorded:}

Put \(\mathcal P(s)=\prod_p(1+p^{-2s}+p^{-3s})\). Since
\(1+p^{-3s}/(1+p^{-2s})=(1+p^{-2s}+p^{-3s})/(1+p^{-2s})\) and
\(\prod_p(1+p^{-2s})=\zeta(2s)/\zeta(4s)\), T45 can be written
\[G_2(s)=\zeta(2s)\,\zeta(4s)\,\mathcal P(s)+\zeta(2s)\big(F_2(s)-1\big)+1.\qquad(1)\]

\textbf{Theorem.} \(G_2\) is meromorphic on \(\operatorname{Re}s>0\),
and the line \(\operatorname{Re}s=0\) is its natural boundary: it has no
meromorphic continuation through any point of that line.

\textbf{Depends on:} ANS-T45. \textbf{Used by:} ANS-T46h-i.
\textbf{Classical inputs:} Estermann.

\textbf{Relations:} ANS-T45a extends ANS-T45.

\hypertarget{ans-t45b-depth-two-zeta-factorisation-and-real-poles-to-112}{%
\subsection{ANS-T45b: Depth two: zeta factorisation and real poles to
1/12}\label{ans-t45b-depth-two-zeta-factorisation-and-real-poles-to-112}}

\textbf{Proof status:} PROVED. \textbf{Review:} UNNS Research Collective
review: valid.

\textbf{Review record:} Release review records. \textbf{Release and
computation:} Archived release.

\textbf{Source:} ANS-I v009, lines 58-62, 67-67. \textbf{Scope:}
Selected nesting grammar. \textbf{Literature priority:} Unreviewed.

\textbf{Release table summary} (ANS-I/v009): Zeta factorisation of H up
to Re s \textgreater{} 1/12; real poles of G\_2 there at 1/2, 1/3, 1/7,
1/8

\textbf{Statement as recorded:}

With \(z=p^{-s}\) the local factor of \(H\) is \(1-(z^5+z^6)/(1+z^2)\),
and as power series
\[1-\frac{z^5+z^6}{1+z^2}=(1-z^5)(1-z^6)(1-z^7)^{-1}(1-z^8)^{-1}(1-z^9)(1-z^{10})\,\big(1+O(z^{12})\big).\]
(Coefficients computed exactly.) Hence, with \(H_{12}\) the Euler
product of the last factor, absolutely convergent and holomorphic on
\(\operatorname{Re}s>1/12\),
\[H(s)=\frac{\zeta(7s)\,\zeta(8s)}{\zeta(5s)\,\zeta(6s)\,\zeta(9s)\,\zeta(10s)}\,H_{12}(s),\]
\[G_2(s)=\zeta(2s)^2\zeta(3s)\,H(s)+\zeta(2s)\big(F_2(s)-1\big)+1.\qquad(3)\]

\textbf{Theorem.} On the real segment \(s>1/12\) the poles of \(G_2\)
are at \(1/2\) (double) and at \(1/3\), \(1/7\), \(1/8\) (simple).

\textbf{Depends on:} ANS-T45. \textbf{Used by:} ANS-T38a, ANS-T45d.

\hypertarget{ans-t45d-depth-two-exact-pole-divisor-on-re-s-0}{%
\subsection{ANS-T45d: Depth two: exact pole divisor on Re s
\textgreater{}
0}\label{ans-t45d-depth-two-exact-pole-divisor-on-re-s-0}}

\textbf{Proof status:} PROVED. \textbf{Review:} UNNS Research Collective
review: valid.

\textbf{Review record:} Release review records. \textbf{Release and
computation:} Archived release.

\textbf{Source:} ANS-I v010, lines 27-33, 35-38, 52-52. \textbf{Scope:}
Selected nesting grammar. \textbf{Literature priority:} Unreviewed.

\textbf{Release table summary} (ANS-I/v010): Exact pole divisor of G\_2
on Re s \textgreater{} 0; all real poles; all non-real poles under RH

\textbf{Statement as recorded:}

Let \(h(z)=1-(z^5+z^6)/(1+z^2)=(1-z^3)(1+z^2+z^3)/(1+z^2)\) be the local
factor of \(H\), and define integers \(c_n\) by
\[h(z)=\prod_{n\ge1}(1-z^n)^{-c_n}.\qquad(1)\] Put \(e_2=c_2+2\),
\(e_3=c_3+1\) and \(e_n=c_n\) otherwise. Since \(c_1=\dots=c_4=0\) we
have \(e_1=e_4=0\), \(e_2=2\), \(e_3=1\). For \(M\ge4\) let
\(H_{M+1}=H\prod_{n\le M}\zeta(ns)^{-c_n}\). Its local factor
\(h(z)\prod_{n\le M}(1-z^n)^{c_n}\) is analytic in \(|z|<1\) and equals
\(1+O(z^{M+1})\), so the product converges absolutely and locally
uniformly on \(\operatorname{Re}s>1/(M+1)\). There, by v009 (3),
\[G_2(s)=\mathcal M(s)+\mathcal A(s),\qquad\mathcal M=\prod_{2\le n\le M}\zeta(ns)^{e_n}\,H_{M+1},\qquad\mathcal A=\zeta(2s)(F_2-1)+1,\qquad(2)\]
and \(\mathcal A\) is holomorphic on \(\operatorname{Re}s>0\) apart from
a simple pole at \(1/2\).

For a complex number \(w\) let \(\mu(w)\) be the multiplicity of \(w\)
as a zero of \(\zeta\) (\(0\) if \(\zeta(w)\ne0\)). For \(s_0\) with
\(\operatorname{Re}s_0>0\) let \(\delta(s_0)=1\) if
\(1+p^{-2s_0}+p^{-3s_0}=0\) for some prime \(p\), and \(0\) otherwise.

\textbf{Theorem.} (i) The poles of \(G_2\) on the half-line \(s>0\) are
exactly the points \(1/n\) with \(e_n>0\), and the order at \(1/n\) is
\(e_n\).

\begin{enumerate}
\def\labelenumi{(\roman{enumi})}
\setcounter{enumi}{1}
\tightlist
\item
  For non-real \(s_0\) with \(\sigma_0=\operatorname{Re}s_0>0\), the
  order of the pole of \(G_2\) at \(s_0\) is \(\max(0,-\nu(s_0))\),
  where \[\nu(s_0)=\sum_{2\le n<1/\sigma_0}e_n\,\mu(ns_0)+\delta(s_0).\]
\end{enumerate}

\textbf{Corollary 2.} Assume the Riemann hypothesis. The non-real poles
of \(G_2\) in \(\operatorname{Re}s>0\) are exactly the points
\(\rho/n\), with \(\rho\) a non-trivial zero and \(e_n<0\), and the
order is \(|e_n|\mu(\rho)\).

\textbf{Depends on:} ANS-T45b. \textbf{Used by:} ANS-T38d, ANS-T38e,
ANS-T45c. \textbf{Classical inputs:} Riemann hypothesis only where
stated.

Corollary 2 (the non-real poles) assumes the Riemann hypothesis; the
theorem itself does not.

\hypertarget{ans-t45c-depth-two-pole-of-order-two-at-112}{%
\subsection{ANS-T45c: Depth two: pole of order two at
1/12}\label{ans-t45c-depth-two-pole-of-order-two-at-112}}

\textbf{Proof status:} PROVED. \textbf{Review:} UNNS Research Collective
review: valid.

\textbf{Review record:} Release review records. \textbf{Release and
computation:} Archived release.

\textbf{Source:} ANS-I v010, lines 44-50. \textbf{Scope:} Selected
nesting grammar. \textbf{Literature priority:} Unreviewed.

\textbf{Release table summary} (ANS-I/v010): G\_2 has a pole of order
two at 1/12

\textbf{Statement as recorded:}

\textbf{Corollary 1 (ANS-T45c, UNNS Research Collective review).}
\(G_2\) has a pole of order exactly two at \(1/12\): \(c_{12}=2\).

The first real poles and their orders are:

\begin{longtable}[]{@{}lllllllllllllll@{}}
\toprule
\(n\) & 2 & 3 & 7 & 8 & 12 & 13 & 17 & 18 & 21 & 22 & 23 & 26 & 27 &
28\tabularnewline
\midrule
\endhead
order at \(1/n\) & 2 & 1 & 1 & 1 & 2 & 1 & 3 & 2 & 1 & 6 & 2 & 4 & 13 &
3\tabularnewline
\bottomrule
\end{longtable}

\textbf{Depends on:} ANS-T45d.

\hypertarget{ans-t45e-depth-two-asymptotics-and-sign-densities-of-the-euler-exponents}{%
\subsection{ANS-T45e: Depth two: asymptotics and sign densities of the
Euler
exponents}\label{ans-t45e-depth-two-asymptotics-and-sign-densities-of-the-euler-exponents}}

\textbf{Proof status:} PROVED. \textbf{Review:} UNNS Research Collective
review: valid.

\textbf{Review record:} Release review records. \textbf{Release and
computation:} Archived release.

\textbf{Source:} ANS-I v010, lines 70-76. \textbf{Scope:} Selected
nesting grammar. \textbf{Literature priority:} Unreviewed.

\textbf{Release table summary} (ANS-I/v010): Asymptotic formula for the
exponents c\_n; each sign has density 1/2; unbounded pole orders

\textbf{Statement as recorded:}

Let \(\alpha\) be a non-real root of \(x^3+x+1\), put
\(R=|\alpha|=\beta^{1/2}=1.2106\ldots\) and \(\theta=\arg\alpha\), and
let \(p_N\) be the sum of the \(N\)-th powers of the three roots.

\textbf{Theorem.} (i) \(n\,c_n=\sum_{d\mid n}\mu_{\rm M}(n/d)\,A_d\),
where \(\mu_{\rm M}\) is the Möbius function and
\[A_N=-p_N-3\,[3\mid N]+2(-1)^{N/2}[2\mid N].\] (ii)
\(n\,c_n=-2R^n\cos n\theta+O(\tau(n)R^{n/2})\).

\begin{enumerate}
\def\labelenumi{(\roman{enumi})}
\setcounter{enumi}{2}
\tightlist
\item
  The sets \(\{n:e_n<0\}\) and \(\{n:e_n>0\}\) both have natural density
  \(\tfrac12\).
\end{enumerate}

\textbf{Depends on:} no earlier result cited. \textbf{Used by:}
ANS-T45f.

\hypertarget{ans-t45f-depth-two-arcsine-law-of-the-normalised-exponents}{%
\subsection{ANS-T45f: Depth two: arcsine law of the normalised
exponents}\label{ans-t45f-depth-two-arcsine-law-of-the-normalised-exponents}}

\textbf{Proof status:} PROVED. \textbf{Review:} UNNS Research Collective
review: valid.

\textbf{Review record:} Release review records. \textbf{Release and
computation:} Archived release.

\textbf{Source:} ANS-I v011, lines 17-19. \textbf{Scope:} Selected
nesting grammar. \textbf{Literature priority:} Unreviewed.

\textbf{Release table summary} (ANS-I/v011): Arcsine limit law of the
normalised Euler exponents

\textbf{Statement as recorded:}

\textbf{Theorem.} Put \(Y_n=ne_n/(2R^n)\). For \(-1\le a<b\le1\),
\[\lim_{N\to\infty}\frac1N\#\{n\le N:\ a\le Y_n\le b\}=\frac{\arcsin b-\arcsin a}\pi.\]
In particular the set where \(|e_n|>R^n/n\) has density \(\tfrac23\),
split equally between the two signs.

\textbf{Depends on:} ANS-T45e. \textbf{Classical inputs:} Weyl
equidistribution.

\hypertarget{ans-t38d-depth-two-vertical-pole-family-for-every-negative-exponent}{%
\subsection{ANS-T38d: Depth two: vertical pole family for every negative
exponent}\label{ans-t38d-depth-two-vertical-pole-family-for-every-negative-exponent}}

\textbf{Proof status:} PROVED. \textbf{Review:} UNNS Research Collective
review: valid.

\textbf{Review record:} Release review records. \textbf{Release and
computation:} Archived release.

\textbf{Source:} ANS-I v010, lines 58-58. \textbf{Scope:} Selected
nesting grammar. \textbf{Literature priority:} Unreviewed.

\textbf{Release table summary} (ANS-I/v010): Vertical pole family at
rho/n for every n with e\_n \textless{} 0

\textbf{Statement as recorded:}

\textbf{Theorem.} Let \(n\ge2\) with \(e_n<0\). For all but \(O_n(U)\)
zeros \(\rho=\tfrac12+i\gamma\), \(0<\gamma\le U\), the point \(\rho/n\)
is a pole of \(G_2\) of order exactly \(|e_n|\,\mu(\rho)\).

\textbf{Depends on:} ANS-T45d. \textbf{Used by:} ANS-T38b, ANS-T38c,
ANS-T38e. \textbf{Classical inputs:} Bohr-Landau.

\textbf{Relations:} ANS-T38d generalizes ANS-T38a; ANS-T38d generalizes
ANS-T38b.

\hypertarget{ans-t38a-depth-two-poles-at-one-fifth-of-zeta-zeros}{%
\subsection{ANS-T38a: Depth two: poles at one fifth of zeta
zeros}\label{ans-t38a-depth-two-poles-at-one-fifth-of-zeta-zeros}}

\textbf{Proof status:} PROVED. \textbf{Review:} UNNS Research Collective
review: valid.

\textbf{Review record:} Release review records. \textbf{Release and
computation:} Archived release.

\textbf{Source:} ANS-I v009, lines 73-73. \textbf{Scope:} Selected
nesting grammar. \textbf{Literature priority:} Unreviewed.

\textbf{Release table summary} (ANS-I/v009): G\_2 has poles at rho/5 for
zeta zeros rho on the critical line outside a set of smaller order;
\textgreater\textgreater{} T log T poles on Re s = 1/10 (uses Selberg,
Bohr-Landau)

\textbf{Statement as recorded:}

\textbf{Theorem.} Let \(\rho=\tfrac12+i\gamma\) be a zero of \(\zeta\)
of multiplicity \(m\) with \(\zeta(n\rho/5)\ne0\) for \(n=2,3,6,7,8,9\).
Then \(G_2\) has a pole of order exactly \(m\) at \(\rho/5\).
Consequently the number of distinct poles of \(G_2\) on
\(\operatorname{Re}s=1/10\) with \(|\operatorname{Im}s|\le T\) is
\(\gg T\log T\).

\textbf{Depends on:} ANS-T45b. \textbf{Classical inputs:} Selberg,
Bohr-Landau.

\textbf{Relations:} ANS-T38d generalizes ANS-T38a.

\hypertarget{ans-t38b-depth-two-poles-at-one-sixth-of-zeta-zeros}{%
\subsection{ANS-T38b: Depth two: poles at one sixth of zeta
zeros}\label{ans-t38b-depth-two-poles-at-one-sixth-of-zeta-zeros}}

\textbf{Proof status:} PROVED. \textbf{Review:} UNNS Research Collective
review: valid.

\textbf{Review record:} Release review records. \textbf{Release and
computation:} Archived release.

\textbf{Source:} ANS-I v010, lines 62-62. \textbf{Scope:} Selected
nesting grammar. \textbf{Literature priority:} Unreviewed.

\textbf{Release table summary} (ANS-I/v010): Poles at rho/6 for all but
O(T) critical-line zeros

\textbf{Statement as recorded:}

The cases \(n=5\) (T38a) and \(n=6\) (T38b, UNNS Research Collective
review) have \(e_n=-1\).

\textbf{Depends on:} ANS-T38d.

\textbf{Relations:} ANS-T38d generalizes ANS-T38b.

\hypertarget{ans-t38c-depth-two-lower-bounds-for-the-number-of-poles}{%
\subsection{ANS-T38c: Depth two: lower bounds for the number of
poles}\label{ans-t38c-depth-two-lower-bounds-for-the-number-of-poles}}

\textbf{Proof status:} PROVED. \textbf{Review:} UNNS Research Collective
review: valid.

\textbf{Review record:} Release review records. \textbf{Release and
computation:} Archived release.

\textbf{Source:} ANS-I v010, lines 64-64. \textbf{Scope:} Selected
nesting grammar. \textbf{Literature priority:} Unreviewed.

\textbf{Release table summary} (ANS-I/v010): At least (n/(5 pi)) T log T
poles of order exactly abs(e\_n) on Re s = 1/(2n) when e\_n \textless{}
0 (uses Conrey)

\textbf{Statement as recorded:}

\textbf{Corollary (ANS-T38c; the method is the UNNS Research
Collective's).} The number of poles of \(G_2\) of order exactly
\(|e_n|\) on \(\operatorname{Re}s=1/(2n)\) with
\(0<\operatorname{Im}s\le T\) is at least
\(\big(\tfrac{n}{5\pi}+o(1)\big)T\log T\). In particular \(\rho/5\) and
\(\rho/6\) give at least \((1/\pi+o(1))T\log T\) and
\((6/(5\pi)+o(1))T\log T\) \textbf{simple} poles.

\textbf{Depends on:} ANS-T38d. \textbf{Classical inputs:} Conrey.

\hypertarget{ans-t38e-depth-two-pole-populations-on-the-lines-re-s-12n}{%
\subsection{ANS-T38e: Depth two: pole populations on the lines Re s =
1/(2n)}\label{ans-t38e-depth-two-pole-populations-on-the-lines-re-s-12n}}

\textbf{Proof status:} PROVED. \textbf{Review:} UNNS Research Collective
review: valid.

\textbf{Review record:} Release review records. \textbf{Release and
computation:} Archived release.

\textbf{Source:} ANS-I v011, lines 25-29. \textbf{Scope:} Selected
nesting grammar. \textbf{Literature priority:} Unreviewed.

\textbf{Release table summary} (ANS-I/v011): Number of distinct poles of
G\_2 on Re s = 1/(2n), for fixed n

\textbf{Statement as recorded:}

Let \(N_0^*(T)\) count the distinct zeros \(\tfrac12+i\gamma\) of
\(\zeta\) with \(0<\gamma\le T\), and \(P_n(T)\) the distinct poles of
\(G_2\) on \(\operatorname{Re}s=1/(2n)\) with
\(0<\operatorname{Im}s\le T\).

\textbf{Theorem.} For each fixed \(n\ge1\),
\[P_n(T)=\begin{cases}N_0^*(nT)+O_n(T),&e_n<0,\\O_n(T),&e_n\ge0.\end{cases}\]
Hence \(P_n(T)\asymp_nT\log T\) when \(e_n<0\).

\textbf{Depends on:} ANS-T45d, ANS-T38d. \textbf{Classical inputs:}
Selberg, Bohr-Landau.

\hypertarget{g.-depth-three}{%
\section{G. Depth three}\label{g.-depth-three}}

\hypertarget{ans-t46c-linked-letter-model-only-a-simple-pole-at-13}{%
\subsection{ANS-T46c: Linked-letter model: only a simple pole at
1/3}\label{ans-t46c-linked-letter-model-only-a-simple-pole-at-13}}

\textbf{Proof status:} PROVED (model). \textbf{Review:} UNNS Research
Collective review: valid.

\textbf{Review record:} Release review records. \textbf{Release and
computation:} Archived release.

\textbf{Source:} ANS-I v010, lines 90-95. \textbf{Scope:} Selected
nesting grammar. \textbf{Literature priority:} Classical shifted divisor
sum; no novelty claimed.

\textbf{Release table summary} (ANS-I/v010): Model series with a letter
linked on both sides: only a simple pole at 1/3 (classical shifted
divisor sum)

\textbf{Statement as recorded:}

In the mechanism of v009 §3, a middle letter \(M\) may receive a prime
\(q\) from the first letter (\(q\mid M\)) and send a prime \(r\) to the
last (\(r\mid M+1\)). The weights are about \(q^{-s}\) and \(r^{-s}\).
This leads to the series
\[U(s)=\sum_{M\ge1}M^{-2s}\sigma_{-s}(M)\,\sigma_{-s}(M+1),\qquad\sigma_{-s}(M)=\sum_{d\mid M}d^{-s},\]
absolutely convergent for \(\operatorname{Re}s>1/2\). It is a
\textbf{model}: \(G_3\) also contains Möbius corrections, the link from
the first letter to the last, and the interaction of the three gcds.
Counting each link separately suggests two factors \(\zeta(3s)\).

\textbf{Theorem.} \(U\) is meromorphic on \(\operatorname{Re}s>3/10\).
Its poles there are simple poles at \(1/2\) and \(1/3\), and
\[\operatorname{Res}_{s=1/3}U=\frac{2\,\zeta(2/3)\,\zeta(4/3)}{3\,\zeta(2)}\ne0.\]

\textbf{Depends on:} no earlier result cited. \textbf{Used by:}
ANS-T46f, ANS-T46j. \textbf{Classical inputs:} Ingham-Estermann, Weil
bound, Shparlinski.

A model series, not the depth-three function.

\hypertarget{ans-t46f-linked-letter-model-residue-at-one-half}{%
\subsection{ANS-T46f: Linked-letter model: residue at one
half}\label{ans-t46f-linked-letter-model-residue-at-one-half}}

\textbf{Proof status:} PROVED (model). \textbf{Review:} UNNS Research
Collective review: valid.

\textbf{Review record:} Release review records. \textbf{Release and
computation:} Archived release.

\textbf{Source:} ANS-I v011, lines 37-37. \textbf{Scope:} Selected
nesting grammar. \textbf{Literature priority:} Unreviewed.

\textbf{Release table summary} (ANS-I/v011): Residue of the
linked-letter model at 1/2

\textbf{Statement as recorded:}

\textbf{Theorem.}
\(\operatorname{Res}_{s=1/2}U=\zeta(3/2)^2/(2\zeta(3))\).

\textbf{Depends on:} ANS-T46c.

A model series, not the depth-three function.

\hypertarget{ans-t46d-depth-three-exact-decomposition-for-re-s-12}{%
\subsection{ANS-T46d: Depth three: exact decomposition for Re s
\textgreater{}
1/2}\label{ans-t46d-depth-three-exact-decomposition-for-re-s-12}}

\textbf{Proof status:} PROVED. \textbf{Review:} UNNS Research Collective
review: valid.

\textbf{Review record:} Release review records. \textbf{Release and
computation:} Archived release.

\textbf{Source:} ANS-I v011, lines 45-46, 49-56, 60-65. \textbf{Scope:}
Selected nesting grammar. \textbf{Literature priority:} Unreviewed.

\textbf{Release table summary} (ANS-I/v011): Exact decomposition of
G\_3: depth-two factor, a shifted series V, a remainder R, minus G\_2
(Re s \textgreater{} 1/2)

\textbf{Statement as recorded:}

Write a word of length three as \((a-1,\,M,\,N)\) and put \(b=M+1\),
\(c=N+1\), with \(a,b,c\ge2\). Then
\[E=\frac{(abc)^2}{g_1g_2},\qquad g_1=\gcd(N,b^2),\qquad g_2=\gcd\Big(\frac{MN}{g_1},\,a^2\Big).\qquad(1)\]

Let \(z=p^{-s}\). Define
\(\Phi(K)=\sum_{d\mid K}J_s(d)\,\rho(d)^{-2s}\), a multiplicative
function with
\[\Phi(p^v)=1+(1-z)\Big(\Big\lfloor\frac v2\Big\rfloor+\Big\lceil\frac v2\Big\rceil z\Big).\]
Define \(w_1,w_2\) multiplicative, with \(w_i(1)=1\) and for exponents
\(\ge1\)
\[w_1(p^\mu)=\frac{\Phi(p^\mu)+c_\mu}{L_0},\qquad w_2(p^\beta)=\frac{1+z-z^{2\beta+1}/(1+z^2)}{L_0},\qquad L_0=\frac{1+z^2+z^3}{1+z^2},\qquad(2)\]
where \(c_\mu=z^3/(1+z^2)\) for even \(\mu\) and
\(c_\mu=z^2(1-z+z^2)/(1+z^2)\) for odd \(\mu\). On primes the two
weights agree: \[w_1(p)=w_2(p)=\frac{1+z+z^2}{1+z^2+z^3}.\] Put
\[V(s)=\sum_{M\ge1}(M+1)^{-2s}\,w_1(M)\,w_2(M+1).\]

For fixed \(M\) let \(f(N)=g_1^{\,s}\,\Phi(MN/g_1)\), let
\(\Lambda_M=\mu_{\rm M}*f\) be its Möbius transform, and let
\[R(s)=\sum_{M\ge1}(M+1)^{-2s}\sum_{k\ge1}\Lambda_M(k)\,k^{-2s}\,\Delta(s,k),\qquad\Delta(s,k)=\zeta\big(2s,1+\tfrac1k\big)-\zeta(2s).\]

\textbf{Theorem.} For \(\operatorname{Re}s>1/2\),
\[\boxed{\,G_3(s)=\zeta(2s)^2\,\zeta(3s)\,H(s)\,V(s)+\zeta(2s)\,R(s)-G_2(s)\,}\]
where \(\zeta(3s)H(s)=\prod_pL_0\) is the Euler product of depth two.

\textbf{Depends on:} no earlier result cited. \textbf{Used by:}
ANS-T46e, ANS-T46e-i, ANS-T46g, ANS-T46k.

\hypertarget{ans-t46e-i-depth-three-positivity-of-e1-e2-and-the-implication-for-the-pole}{%
\subsection{ANS-T46e-i: Depth three: positivity of E1, E2 and the
implication for the
pole}\label{ans-t46e-i-depth-three-positivity-of-e1-e2-and-the-implication-for-the-pole}}

\textbf{Proof status:} PROVED. \textbf{Review:} UNNS Research Collective
review: valid.

\textbf{Review record:} Release review records. \textbf{Release and
computation:} Archived release.

\textbf{Source:} ANS-I v011, lines 98-103, 111-113. \textbf{Scope:}
Selected nesting grammar. \textbf{Literature priority:} Unreviewed.

\textbf{Release table summary} (ANS-I/v011): E1 and E2 are positive;
T46e implies a pole of order exactly two of G\_3 at 1/3

\textbf{Statement as recorded:}

Take \(s=\tfrac13\), write \(\kappa_i=\mu_{\rm M}*w_i\), and define for
each prime, with \(x=p^{-2/3}\) and \(y=p^{-1}\),
\[W_i(t)=\sum_{j\ge0}w_i(p^j)\,t^j,\qquad E_{1,p}=(1-y)(1-x)W_1(x)+(1-z^2)\big((1-y)W_2(y)-1\big),\]
\(E_{2,p}\) with the indices exchanged, and \(E_i=\prod_pE_{i,p}\).

\textbf{Positivity (proved).} \(E_1,E_2\) converge to positive numbers.
From (2), \(w_1(p^\mu)\ge1\), since \(\Phi(p^\mu)\ge1\) and
\(c_\mu-(L_0-1)\) is \(0\) or \(z^2(1-z)^2/(1+z^2)\); and
\(w_2(p^\beta)\ge1\), since the numerator minus \(L_0\) is
\((z-z^{2\beta+1})/(1+z^2)\). So \((1-y)W_i(y)\ge1\) and each
\(E_{i,p}>0\). Expanding the closed forms in \(z=p^{-1/3}\) gives
\(E_{i,p}=1+O(z^4)\).

\textbf{Proposition (proved implication).} If ANS-T46e holds, \(G_3\)
has a pole of order exactly two at \(1/3\) and
\[\lim_{s\to1/3}\big(s-\tfrac13\big)^2G_3(s)=\frac{\zeta(2/3)^2\,H(1/3)\,r_V}3\ne0.\]
So T46b holds for \(k=3\).

\textbf{Depends on:} ANS-T46d. \textbf{Used by:} ANS-T46e.

\textbf{Relations:} ANS-T46e discharges the hypothesis of ANS-T46e-i.

\hypertarget{ans-t46g-depth-three-square-divisor-convolution-identities}{%
\subsection{ANS-T46g: Depth three: square-divisor convolution
identities}\label{ans-t46g-depth-three-square-divisor-convolution-identities}}

\textbf{Proof status:} PROVED. \textbf{Review:} UNNS Research Collective
review: valid.

\textbf{Review record:} Release review records. \textbf{Release and
computation:} Archived release.

\textbf{Source:} ANS-I v012, lines 17-24. \textbf{Scope:} Selected
nesting grammar. \textbf{Literature priority:} Unreviewed.

\textbf{Release table summary} (ANS-I/v012): Exact convolution
identities for w1, w2 (square-divisor extraction)

\textbf{Statement as recorded:}

Let \(z=p^{-s}\), \(P(z)=1+z^2+z^3\), \(A=z(1-z^2)/P\),
\(B=(1-z)(1+z^3)/P\). The UNNS Research Collective's identities
\[w_1(p^{2j+1})-w_1(p^{2j})=A,\quad w_1(p^{2j+2})-w_1(p^{2j+1})=B,\quad w_2(p^j)-w_2(p^{j-1})=Az^{2j-2}\]
\((j\ge0\), resp. \(j\ge1)\) are \textbf{valid}: they follow from (2) of
v011, and were checked symbolically. Hence
\(w_1=\sigma_{-s}*I_\square*r_{1}\) and \(w_2=\sigma_{-s}*r_2\) with the
stated local factors. Put \(q_1=I_\square*r_1\), so that
\(w_1=\sigma_{-s}*q_1\), with local series
\[Q(t)=\frac{(1-zt)\big(1+At+(B-1)t^2\big)}{1-t^2},\qquad R_2(t)=(1-zt)\Big(1+\frac{At}{1-z^2t}\Big).\qquad(1)\]
For \(j\ge1\) the coefficients are: \(q_1(p)=A-z\); \(q_1(p^j)=B-zA\)
for even \(j\) and \(A-zB\) for odd \(j\ge3\); \(r_2(p)=A-z\);
\(r_2(p^j)=-A(1-z)z^{2j-3}\) for \(j\ge2\). Consequently, uniformly in
\(\operatorname{Im}s\) for \(\sigma\) in compact subsets of
\((0.3,\infty)\) (where \(P\) is bounded away from \(0\)),
\[|q_1(p)|,|r_2(p)|\le Cp^{-3\sigma},\qquad|q_1(p^j)|\le C,\qquad|r_2(p^j)|\le Cp^{-(2j-2)\sigma}\quad(j\ge2).\qquad(\mathrm H)\]

One qualification: the extra decay of the corrections is at primes.
\(q_1(p^2)\to1\), and \(r_2(p^2)\) is only \(O(p^{-2\sigma})\).

\textbf{Depends on:} ANS-T46d. \textbf{Used by:} ANS-T46e, ANS-T46h,
ANS-T46i.

\hypertarget{ans-t46h-depth-three-first-remainder-component-continues-past-13}{%
\subsection{ANS-T46h: Depth three: first remainder component continues
past
1/3}\label{ans-t46h-depth-three-first-remainder-component-continues-past-13}}

\textbf{Proof status:} PROVED. \textbf{Review:} UNNS Research Collective
review: valid.

\textbf{Review record:} Release review records. \textbf{Release and
computation:} Archived release.

\textbf{Source:} ANS-I v012, lines 28-28. \textbf{Scope:} Selected
nesting grammar. \textbf{Literature priority:} Unreviewed.

\textbf{Release table summary} (ANS-I/v012): R\_1 continues to Re s
\textgreater{} 1/4; poles and residue at 1/3

\textbf{Statement as recorded:}

The UNNS Research Collective's factorisation
\(\sum\Phi(M)M^{-2s}=\zeta(2s)\zeta(3s)\zeta(4s)B_\Phi(s)\), the
continuation of \(R_1\) to \(\sigma>\frac14\) and the residue at \(1/3\)
are \textbf{valid}.

\textbf{Depends on:} ANS-T46g. \textbf{Used by:} ANS-T46i.

Consolidation note: the releases record the factorisation behind this
result as the UNNS Research Collective's and declare it valid without
writing out a derivation. A derivation is in
\texttt{PROVENANCE/AUDIT\_v001.md}, section 5; the foundational audit of
27 November 2026 accepted it.

\hypertarget{ans-t46h-i-depth-three-r_1-has-a-natural-boundary}{%
\subsection{ANS-T46h-i: Depth three: R\_1 has a natural
boundary}\label{ans-t46h-i-depth-three-r_1-has-a-natural-boundary}}

\textbf{Proof status:} PROVED. \textbf{Review:} UNNS Research Collective
review: valid, with a recorded qualification. A normal-convergence
clarification was requested and added in ANS-I v013.

\textbf{Review record:} Release review records. \textbf{Release and
computation:} Archived release.

\textbf{Source:} ANS-I v012, lines 30-30. \textbf{Scope:} Selected
nesting grammar. \textbf{Literature priority:} Unreviewed.

\textbf{Release table summary} (ANS-I/v012): R\_1 meromorphic on Re s
\textgreater{} 0 with natural boundary Re s = 0

\textbf{Statement as recorded:}

\textbf{Strengthening (ANS-T46h-i).} \(R_1\) is meromorphic on
\(\sigma>0\), and the line \(\sigma=0\) is a natural boundary.

\textbf{Depends on:} ANS-T45a. \textbf{Classical inputs:} Estermann.

\hypertarget{ans-t46i-depth-three-remainder-meromorphic-on-re-s-27}{%
\subsection{ANS-T46i: Depth three: remainder meromorphic on Re s
\textgreater{}
2/7}\label{ans-t46i-depth-three-remainder-meromorphic-on-re-s-27}}

\textbf{Proof status:} PROVED. \textbf{Review:} UNNS Research Collective
review: valid.

\textbf{Review record:} Release review records. \textbf{Release and
computation:} Archived release.

\textbf{Source:} ANS-I v012, lines 36-36. \textbf{Scope:} Selected
nesting grammar. \textbf{Literature priority:} Unreviewed.

\textbf{Release table summary} (ANS-I/v012): R meromorphic on Re s
\textgreater{} 2/7; at most simple poles at 1/2, 1/3

\textbf{Statement as recorded:}

\textbf{Theorem.} \(R\) is meromorphic on \(\sigma>\frac27\). Its only
possible singularities there are simple poles at \(\frac12\) and
\(\frac13\).

\textbf{Depends on:} ANS-T46h, ANS-T46g. \textbf{Used by:} ANS-T46e.
\textbf{Classical inputs:} convexity bounds for L-functions.

\hypertarget{ans-t46j-shifted-series-with-square-supported-weights}{%
\subsection{ANS-T46j: Shifted series with square-supported
weights}\label{ans-t46j-shifted-series-with-square-supported-weights}}

\textbf{Proof status:} PROVED. \textbf{Review:} UNNS Research Collective
review: valid, with a recorded qualification. Two scope clarifications
were requested and made in ANS-I v013.

\textbf{Review record:} Release review records. \textbf{Release and
computation:} Archived release.

\textbf{Source:} ANS-I v012, lines 65-72. \textbf{Scope:} Selected
nesting grammar. \textbf{Literature priority:} Unreviewed.

\textbf{Release table summary} (ANS-I/v012): Shifted series with weights
satisfying (H): continuation to Re s \textgreater{} 1/3 - 1/1500 and
residue at 1/3

\textbf{Statement as recorded:}

Let \(q,r\) be multiplicative, with values holomorphic in \(s\),
satisfying (H) with \(q_1,r_2\) replaced by \(q,r\) for
\(\sigma>\frac13-\delta_0\), where \(\delta_0=1/1500\); constants are
uniform in \(\operatorname{Im}s\) up to a power of \(|s|\). Put
\(w_q=\sigma_{-s}*q\) and
\[V(s)=\sum_{M\ge1}(M+1)^{-2s}\,w_q(M)\,w_r(M+1).\]

\textbf{Theorem.} \(V\) is meromorphic on \(\sigma>\frac13-\delta_0\).
Its only possible singularities there are simple poles at \(\frac12\)
and \(\frac13\). With \(\kappa_q=\mu_{\rm M}*w_q\) (so
\(\kappa_q(d)=\sum_{am=d}q(a)m^{-s}\)), all quantities taken at
\(s=\frac13\),
\[\operatorname{Res}_{1/3}V=\frac{\zeta(2/3)}3\,(E_1+E_2),\qquad E_1=\sum_{(e,a)=1}\frac{\kappa_r(e)}{e}\prod_{p\mid e}(1-p^{-2/3})\cdot\frac{q(a)}{a^{2/3}},\]
and \(E_2\) is the same with \(q\) and \(r\) exchanged.

The hypotheses hold for \(q=q_1\), \(r=r_2\) (the series of T46d), for
\(q=I_\square\), \(r=\delta_1\) (the UNNS Research Collective's model
\(U_\square\)), and for \(q=r=\delta_1\) (the model of T46c).

\textbf{Depends on:} ANS-T46c. \textbf{Used by:} ANS-T46e, ANS-T46k-i.
\textbf{Classical inputs:} Weil-Estermann.

\hypertarget{ans-t46e-depth-three-continuation-of-v-and-r-and-the-residue-of-v}{%
\subsection{ANS-T46e: Depth three: continuation of V and R and the
residue of
V}\label{ans-t46e-depth-three-continuation-of-v-and-r-and-the-residue-of-v}}

\textbf{Proof status:} PROVED. \textbf{Review:} UNNS Research Collective
review: valid.

\textbf{Review record:} Release review records. \textbf{Release and
computation:} Archived release.

\textbf{Source:} ANS-I v012, lines 160-162. \textbf{Scope:} Selected
nesting grammar. \textbf{Literature priority:} Unreviewed.

\textbf{Release table summary} (ANS-I/v012): V and R continue past 1/3;
residue of V is zeta(2/3)(E1+E2)/3 \textless{} 0

\textbf{Statement as recorded:}

\textbf{Theorem (ANS-T46e).} \(V\) and \(R\) of T46d continue
meromorphically past \(\frac13\), and
\[\operatorname{Res}_{1/3}V=r_V=\frac{\zeta(2/3)(E_1+E_2)}3<0,\] with
\(E_1,E_2\) the Euler products of v011 §3.1. This is the residue
conjectured there.

\textbf{Depends on:} ANS-T46d, ANS-T46g, ANS-T46i, ANS-T46j, ANS-T46e-i.
\textbf{Used by:} ANS-T46k.

\textbf{Relations:} ANS-T46e discharges the hypothesis of ANS-T46e-i.

\hypertarget{ans-t46k-depth-three-pole-of-order-exactly-two-at-13}{%
\subsection{ANS-T46k: Depth three: pole of order exactly two at
1/3}\label{ans-t46k-depth-three-pole-of-order-exactly-two-at-13}}

\textbf{Proof status:} PROVED. \textbf{Review:} UNNS Research Collective
review: valid.

\textbf{Review record:} Release review records. \textbf{Release and
computation:} Archived release.

\textbf{Source:} ANS-I v012, lines 168-170. \textbf{Scope:} Selected
nesting grammar. \textbf{Literature priority:} Unreviewed.

\textbf{Release table summary} (ANS-I/v012): G\_3 meromorphic on Re s
\textgreater{} 1/3 - 1/1500; pole of order exactly two at 1/3

\textbf{Statement as recorded:}

\textbf{Corollary (ANS-T46k).} \(G_3\) is meromorphic on
\(\sigma>\frac13-\frac1{1500}\). Its only singularities there are the
pole of order three at \(\frac12\) and a pole of order \textbf{exactly
two} at \(\frac13\), with
\[\lim_{s\to1/3}\big(s-\tfrac13\big)^2G_3(s)=\frac{\zeta(2/3)^2\,H(1/3)\,r_V}3\approx-0.75916.\]
In particular ANS-T46b holds for \(k=3\).

\textbf{Depends on:} ANS-T46d, ANS-T46e, ANS-T45, ANS-T34. \textbf{Used
by:} ANS-T46b, ANS-T46l, ANS-T46s.

\textbf{Relations:} ANS-T46k-i widens the margin of ANS-T46k.

\hypertarget{ans-t46l-depth-three-second-term-for-smoothed-counts}{%
\subsection{ANS-T46l: Depth three: second term for smoothed
counts}\label{ans-t46l-depth-three-second-term-for-smoothed-counts}}

\textbf{Proof status:} PROVED. \textbf{Review:} UNNS Research Collective
review: valid.

\textbf{Review record:} Release review records. \textbf{Release and
computation:} Archived release.

\textbf{Source:} ANS-I v012, lines 174-176. \textbf{Scope:} Selected
nesting grammar. \textbf{Literature priority:} Unreviewed.

\textbf{Release table summary} (ANS-I/v012): Depth three: second term
X\^{}(1/3)(a log X + b) for smoothed counts

\textbf{Statement as recorded:}

\textbf{Corollary (smoothed second term).} For all large \(m\) there are
constants \(b_m\) and \(\eta>0\) with
\[\sum_{E\le X}\Big(1-\frac EX\Big)^m=\mathrm{Res}_{1/2}+X^{1/3}\big(a_m\log X+b_m\big)+O\big(X^{1/3-\eta}\big),\]
where the sum is over words of length three, \(\mathrm{Res}_{1/2}\) is
the residue of \(G_3(s)X^s\,m!/(s(s+1)\cdots(s+m))\) at \(\frac12\), and
\(a_m=\dfrac{\zeta(2/3)^2H(1/3)r_V}3\cdot\dfrac{m!}{\prod_{j=0}^m(\frac13+j)}\ne0\).

\textbf{Depends on:} ANS-T46k. \textbf{Classical inputs:} Perron-type
smoothing.

\hypertarget{ans-t46k-i-depth-three-continuation-margin-1200}{%
\subsection{ANS-T46k-i: Depth three: continuation margin
1/200}\label{ans-t46k-i-depth-three-continuation-margin-1200}}

\textbf{Proof status:} PROVED. \textbf{Review:} UNNS Research Collective
review: valid.

\textbf{Review record:} Release review records. \textbf{Release and
computation:} Archived release.

\textbf{Source:} ANS-I v013, lines 24-24. \textbf{Scope:} Selected
nesting grammar. \textbf{Literature priority:} Unreviewed.

\textbf{Release table summary} (ANS-I/v013): Continuation margin of
G\_3: Re s \textgreater{} 1/3 - 1/200

\textbf{Statement as recorded:}

\textbf{Proposition.} T46j, and therefore T46k and T46l, hold with
\(\delta_0=1/200\). So \(G_3\) is meromorphic on
\(\sigma>\frac13-\frac1{200}\).

\textbf{Depends on:} ANS-T46j.

\textbf{Relations:} ANS-T46k-i widens the margin of ANS-T46k.

\hypertarget{h.-arbitrary-depth-at-one-third}{%
\section{H. Arbitrary depth at one
third}\label{h.-arbitrary-depth-at-one-third}}

\hypertarget{ans-t46m-every-depth-decomposition-by-end-letters}{%
\subsection{ANS-T46m: Every depth: decomposition by end
letters}\label{ans-t46m-every-depth-decomposition-by-end-letters}}

\textbf{Proof status:} PROVED. \textbf{Review:} UNNS Research Collective
review: valid.

\textbf{Review record:} Release review records. \textbf{Release and
computation:} Archived release.

\textbf{Source:} ANS-I v013, lines 42-57. \textbf{Scope:} Selected
nesting grammar. \textbf{Literature priority:} Unreviewed.

\textbf{Release table summary} (ANS-I/v013): Decomposition of every G\_k
(k \textgreater= 2, Re s \textgreater{} 1/2) by its end letters; weight
depends on source and endpoint of the middle word

\textbf{Statement as recorded:}

For a word \(\mu\) let \(E(\mu)\) be its endpoint and \(K(\mu)\) its
minimal source; \(E=K=1\) for the empty word. By T03 and the
first-letter identity of v008,
\[E(x\,w)=\frac{(x+1)^2E(w)}{\gcd(K(w),(x+1)^2)},\] and, with
\(g=\gcd(P,E(\mu))\),
\[E(\mu P)=\frac{(P+1)^2E(\mu)}g,\qquad K(\mu P)=\frac{PK(\mu)}g.\] Here
\(E=\prod(x+1)^2/G\) and \(K=\prod x/G\) with the same total gcd \(G\),
which gives the formula for \(K(\mu P)\). Hence a word with first letter
\(a-1\), middle word \(\mu\) and last letter \(P\) has
\[E=\frac{a^2(P+1)^2E(\mu)}{g\,h},\qquad h=\gcd\Big(\frac{PK(\mu)}g,\;a^2\Big).\qquad(1)\]

Let \(\Phi\), \(L_0=(1+z^2+z^3)/(1+z^2)\) and
\(c_u=\sum_{i\ge1}\big(\Phi(p^{u+i})-\Phi(p^{u+i-1})\big)z^{2i}\) be as
in v011, so that \(c_u=z^3/(1+z^2)\) for even \(u\) and
\(z^2(1-z+z^2)/(1+z^2)\) for odd \(u\). For \(u,v\ge0\) put
\[L(u,v)=\Phi(p^u)\,\big(1+z-z^{v+1}\big)+z^v\,c_u,\qquad\Omega(K,E)=\prod_p\frac{L(v_pK,\,v_pE)}{L_0}.\qquad(2)\]
Note \(L(0,0)=L_0\). For \(k\ge2\) define
\[M_k(s)=\sum_{|\mu|=k-2}E(\mu)^{-s}\,\Omega\big(K(\mu),E(\mu)\big),\]
and, with
\(f_\mu(P)=\gcd(P,E(\mu))^s\,\Phi\big(PK(\mu)/\gcd(P,E(\mu))\big)\) and
\(\Lambda_\mu=\mu_{\rm M}*f_\mu\),
\[R_k(s)=\sum_{|\mu|=k-2}E(\mu)^{-s}\sum_{q\ge1}\Lambda_\mu(q)\,q^{-2s}\,\Delta(s,q),\qquad\Delta(s,q)=\zeta\big(2s,1+\tfrac1q\big)-\zeta(2s).\]

\textbf{Theorem.} For every \(k\ge2\) and \(\sigma>\frac12\), with all
series absolutely convergent,
\[\boxed{\,G_k(s)=\zeta(2s)^2\,\zeta(3s)\,H(s)\,M_k(s)\;+\;\zeta(2s)\,R_k(s)\;-\;G_{k-1}(s).\,}\]

\textbf{Depends on:} ANS-T03, ANS-T34, ANS-T12. \textbf{Used by:}
ANS-T46n, ANS-T46p, ANS-T46q.

\hypertarget{ans-t46n-links-active-at-one-prime-maximum-floork24}{%
\subsection{ANS-T46n: Links active at one prime: maximum
floor(k\^{}2/4)}\label{ans-t46n-links-active-at-one-prime-maximum-floork24}}

\textbf{Proof status:} PROVED. \textbf{Review:} UNNS Research Collective
review: valid.

\textbf{Review record:} Release review records. \textbf{Release and
computation:} Archived release.

\textbf{Source:} ANS-I v013, lines 75-77. \textbf{Scope:} Selected
nesting grammar. \textbf{Literature priority:} Unreviewed.

\textbf{Release table summary} (ANS-I/v013): Links active at one prime:
at most floor(k\^{}2/4); maximal role assignments

\textbf{Statement as recorded:}

Write a word as \(x_1\cdots x_k\) and \(a_i=x_i+1\). A prime \(p\)
activates the link \((i,j)\), \(i<j\), when \(p\mid a_i\) and
\(p\mid x_j\). Since \(\gcd(x_i,a_i)=1\), at one prime a letter can
supply (\(p\mid a_i\), role S) or receive (\(p\mid x_i\), role R), not
both. So the sets of links active at one prime are contained in the sets
\(\{(i,j):i\in S,\ j\in R,\ i<j\}\) for assignments of roles.

\textbf{Proposition.} The largest such set has \(\lfloor k^2/4\rfloor\)
links. It requires all S letters to precede all R letters, with
\(\lfloor k/2\rfloor\) or \(\lceil k/2\rceil\) letters of role S: one
assignment for even \(k\), two for odd \(k\ge3\) (the first letter taken
as S and the last as R).

\textbf{Depends on:} ANS-T46m.

\hypertarget{ans-t46o-model-f_rt-pole-of-order-exactly-rt-at-13}{%
\subsection{ANS-T46o: Model F\_\{r,t\}: pole of order exactly rt at
1/3}\label{ans-t46o-model-f_rt-pole-of-order-exactly-rt-at-13}}

\textbf{Proof status:} PROVED (model). \textbf{Review:} UNNS Research
Collective review: valid.

\textbf{Review record:} Release review records. \textbf{Release and
computation:} Archived release.

\textbf{Source:} ANS-I v013, lines 103-108. \textbf{Scope:} Selected
nesting grammar. \textbf{Literature priority:} Unreviewed.

\textbf{Release table summary} (ANS-I/v013): Model F\_\{r,t\}:
meromorphic on Re s \textgreater{} 1/4, pole of order exactly rt at 1/3

\textbf{Statement as recorded:}

Keep only the links from \(r\) supplying letters to \(t\) later
receiving letters:
\[F_{r,t}(s)=\sum_{a\in\mathbb N^r}\ \sum_{x\in\mathbb N^t}\ \prod_ia_i^{-2s}\prod_j(x_j+1)^{-2s}\,\gcd\Big(\prod_jx_j,\;\big(\prod_ia_i\big)^2\Big)^s.\]

\textbf{Theorem.} For \(\sigma>\frac12\),
\[F_{r,t}(s)=\sum_{i=0}^t\zeta(2s)^{r+t-i}\,\zeta(3s)^{r(t-i)}\,Q_i(s),\]
with every \(Q_i\) holomorphic on \(\sigma>\frac14\). Hence \(F_{r,t}\)
is meromorphic there, with poles only at \(\frac12\), of order exactly
\(r+t\), and at \(\frac13\), of order \textbf{exactly \(rt\)}. Moreover
\(Q_0(\frac13)>0\), and
\(\lim_{s\to1/3}(s-\frac13)^{rt}F_{r,t}=\zeta(\frac23)^{r+t}\,Q_0(\frac13)\,3^{-rt}\).

\textbf{Depends on:} no earlier result cited. \textbf{Used by:}
ANS-T46s.

A model series, not G\_k.

\hypertarget{ans-t46p-conditions-implying-pole-order-floork24}{%
\subsection{ANS-T46p: Conditions implying pole order
floor(k\^{}2/4)}\label{ans-t46p-conditions-implying-pole-order-floork24}}

\textbf{Proof status:} PROVED (implication). \textbf{Review:} UNNS
Research Collective review: valid.

\textbf{Review record:} Release review records. \textbf{Release and
computation:} Archived release.

\textbf{Source:} ANS-I v013, lines 131-131. \textbf{Scope:} Selected
nesting grammar. \textbf{Literature priority:} Unreviewed.

\textbf{Release table summary} (ANS-I/v013): Conditions on M\_k, R\_k,
G\_\{k-1\} that imply T46b at depth k

\textbf{Statement as recorded:}

\textbf{Proposition (ANS-T46p, proved implication).} Let \(k\ge3\).
Suppose \(M_k\) and \(R_k\) continue meromorphically to a neighbourhood
of \(\frac13\), that \(M_k\) has there a pole of order exactly
\(\lfloor k^2/4\rfloor-1\), and that the pole order of \(R_k\) is less
than \(\lfloor k^2/4\rfloor\). Suppose also that \(G_{k-1}\) continues
there with pole order less than \(\lfloor k^2/4\rfloor\). Then \(G_k\)
has a pole of order exactly \(\lfloor k^2/4\rfloor\) at \(\frac13\).

\textbf{Depends on:} ANS-T46m.

An implication; its hypotheses are not established for any k
\textgreater= 4.

\hypertarget{ans-t46r-middle-word-gcd-identity-and-overlap-decomposition}{%
\subsection{ANS-T46r: Middle-word gcd identity and overlap
decomposition}\label{ans-t46r-middle-word-gcd-identity-and-overlap-decomposition}}

\textbf{Proof status:} PROVED. \textbf{Review:} UNNS Research Collective
review: valid.

\textbf{Review record:} Release review records. \textbf{Release and
computation:} Archived release.

\textbf{Source:} ANS-I v014, lines 47-52. \textbf{Scope:} Selected
nesting grammar. \textbf{Literature priority:} Unreviewed.

\textbf{Release table summary} (ANS-I/v014): Middle-word gcd identity
and overlap decomposition of L(u,v)/L0

\textbf{Statement as recorded:}

\textbf{Theorem.} For a middle word \(xy\) with \(g=\gcd(y,(x+1)^2)\),
\(K=xy/g\) and \(E=(x+1)^2(y+1)^2/g\):

\begin{enumerate}
\def\labelenumi{(\roman{enumi})}
\item
  \(\gcd(K,E)=\gcd\big(x,(y+1)^2\big)\).
\item
  With \(L_0=L(0,0)=1+c_0\), for all \(u,v\ge0\),
  \[\frac{L(u,v)}{L_0}=\frac{L(u,0)}{L_0}\,\frac{L(0,v)}{L_0}+\frac{(1-z^v)(1+z)\,\big[\Phi(p^u)c_0-c_u\big]}{L_0^2}.\]
\end{enumerate}

\textbf{Depends on:} ANS-T47.

\hypertarget{ans-t46s-role-patterns-local-expectation-and-agreement-with-proved-coefficients}{%
\subsection{ANS-T46s: Role patterns: local expectation and agreement
with proved
coefficients}\label{ans-t46s-role-patterns-local-expectation-and-agreement-with-proved-coefficients}}

\textbf{Proof status:} PROVED. \textbf{Review:} UNNS Research Collective
review: valid.

\textbf{Review record:} Release review records. \textbf{Release and
computation:} Archived release.

\textbf{Source:} ANS-I v014, lines 62-69, 73-75, 87-87, 97-98, 106-107.
\textbf{Scope:} Selected nesting grammar. \textbf{Literature priority:}
Unreviewed.

\textbf{Release table summary} (ANS-I/v014): (i)
E\_R(p)=1+L(R)/p+O(p\^{}-4/3); (ii) formula equals proved coefficients
k=2,3; (iii) and the model F\_\{r,t\}

\textbf{Statement as recorded:}

Fix a prime \(p\), put \(x=p^{-2/3}\), \(y=1/p\), \(z=p^{-1/3}\). A
letter \(N\) has the two coprime coordinates \(a=N+1\) and \(N\). A
\textbf{role pattern} \(R\in\{{\sf S},{\sf R}\}^k\) gives each letter a
dominant coordinate: \(a\) for a supplying letter \({\sf S}\), \(N\) for
a receiving letter \({\sf R}\). Under the measure \(\mathbb P_R\) the
letters are independent, and for each letter

\begin{itemize}
\tightlist
\item
  the dominant coordinate has \(p\)-valuation \(j\ge1\) with probability
  \((1-x)x^j\);
\item
  the other coordinate has valuation \(j\ge1\) with probability
  \((1-x)y^j\);
\item
  otherwise both valuations are \(0\).
\end{itemize}

(The three cases have total probability \(1\) for every \(p\).) Let
\(T\) be the \(p\)-valuation of the total gcd \(G\) of the word and put
\[\mathcal E_R(p)=\mathbb E_R\big[p^{T/3}\big],\qquad L(R)=\#\{i<j:R_i={\sf S},\,R_j={\sf R}\}.\]

\textbf{Theorem (T46s-i).} \(\mathcal E_R(p)=1+L(R)/p+O_k(p^{-4/3})\).
Consequently \[C_R=\prod_p\Big(1-\frac1p\Big)^{L(R)}\mathcal E_R(p)\]
converges absolutely and \(C_R>0\).

\textbf{Lemma.} Let \(R\) begin with \({\sf S}\) and end with
\({\sf R}\), and let \((u,v)=(v_pK,v_pE)\) for the middle word. The
expectation over the first and last letters of \(p\) raised to one third
of the two gcd exponents that involve them equals \(L(u,v)\).

\textbf{Theorem (T46s-ii).} For \(k=2,3\), with
\(\ell_k=\lfloor k^2/4\rfloor\),
\[\lim_{s\to1/3}\big(s-\tfrac13\big)^{\ell_k}G_k(s)=\zeta(\tfrac23)^k\,3^{-\ell_k}\sum_{L(R)=\ell_k}C_R.\qquad(\ast)\]

\textbf{Theorem (T46s-iii).} In the model \(F_{r,t}\) of T46o, let the
letters be independent with only their dominant coordinate, of law
\((1-x)x^j\), and let \(T=\min\big(2\sum\alpha_i,\sum\beta_j\big)\).
Then
\[\lim_{s\to1/3}\big(s-\tfrac13\big)^{rt}F_{r,t}(s)=\zeta(\tfrac23)^{r+t}\,3^{-rt}\prod_p\Big(1-\frac1p\Big)^{rt}\mathbb E\big[p^{T/3}\big].\]

\textbf{Depends on:} ANS-T47, ANS-T45, ANS-T46k, ANS-T46o. \textbf{Used
by:} ANS-T46t.

\hypertarget{ans-t46b-pole-of-order-floork24-at-13}{%
\subsection{ANS-T46b: Pole of order floor(k\^{}2/4) at
1/3}\label{ans-t46b-pole-of-order-floork24-at-13}}

\textbf{Proof status:} PROVED for k \textless= 3; CONJECTURE for k
\textgreater= 4. \textbf{Review:} UNNS Research Collective review:
valid. k = 2, 3 proved and reviewed (T45, T46k); k \textgreater= 4
conjecture.

\textbf{Review record:} Release review records. \textbf{Release and
computation:} Archived release.

\textbf{Source:} ANS-I v009, lines 87-89. \textbf{Scope:} Selected
nesting grammar. \textbf{Literature priority:} Unreviewed.

\textbf{Release table summary} (ANS-I/v014): Pole order floor(k\^{}2/4)

\textbf{Statement as recorded:}

\textbf{Conjecture (ANS-T46b).} For \(k\ge2\), \(G_k\) has a pole at
\(s=1/3\) of order \(\lfloor k^2/4\rfloor\).

This gives order \(1\) at \(k=2\) (proved, T45), order \(2\) at \(k=3\)
and order \(4\) at \(k=4\). A proof needs (i) continuation past \(1/3\),
(ii) equidistribution, with a power saving, of the least elements above,
which involves inverses modulo \(P_{\rm out}\) and so estimates of
Kloosterman type, and (iii) a non-vanishing check on the leading
coefficient. The depth-three counts of v008 are compatible with order
\(2\) but do not test it.

\textbf{Depends on:} ANS-T45, ANS-T46k.

\hypertarget{ans-t46q-weighted-middle-word-sum-pole-order-floork24-1}{%
\subsection{ANS-T46q: Weighted middle-word sum: pole order
floor(k\^{}2/4)-1}\label{ans-t46q-weighted-middle-word-sum-pole-order-floork24-1}}

\textbf{Proof status:} CONJECTURE (k = 2, 3 proved). \textbf{Review:}
Not a proved result.

\textbf{Review record:} None. \textbf{Release and computation:} Archived
release.

\textbf{Source:} ANS-I v013, lines 135-137. \textbf{Scope:} Selected
nesting grammar. \textbf{Literature priority:} Unreviewed.

\textbf{Release table summary} (ANS-I/v014): M\_k has order
floor(k\^{}2/4)-1

\textbf{Statement as recorded:}

\textbf{Conjecture (ANS-T46q).} \(M_k\) has a pole of order exactly
\(\lfloor k^2/4\rfloor-1\) at \(\frac13\).

Since \(\lfloor k^2/4\rfloor-1=\lfloor(k-2)^2/4\rfloor+(k-2)\), this
says that the weight \(\Omega\) raises the order of \(G_{k-2}\) by one
for each middle letter. It holds for \(k=2\) (\(M_2=1\)) and \(k=3\)
(v012). This is the form in which depth four can now be tested, and it
is a candidate for an induction on depth; no induction step is proved.

\textbf{Depends on:} ANS-T46m.

\hypertarget{ans-t46t-general-continuation-pole-order-floork24-and-leading-coefficient-at-13}{%
\subsection{ANS-T46t: General continuation, pole order floor(k\^{}2/4)
and leading coefficient at
1/3}\label{ans-t46t-general-continuation-pole-order-floork24-and-leading-coefficient-at-13}}

\textbf{Proof status:} CONJECTURE (k = 2, 3 proved). \textbf{Review:}
Not a proved result.

\textbf{Review record:} None. \textbf{Release and computation:} Archived
release.

\textbf{Source:} ANS-I v015, lines 149-151. \textbf{Scope:} Selected
nesting grammar. \textbf{Literature priority:} Unreviewed.

\textbf{Release table summary} (ANS-I/v015): G\_k: continuation near
1/3, order floor(k\^{}2/4), leading coefficient

\textbf{Statement as recorded:}

Following the review: \textbf{for every \(k\ge2\), \(G_k\) continues
meromorphically to a neighbourhood of \(\frac13\), has there a pole of
order exactly \(\ell_k=\lfloor k^2/4\rfloor\), and its leading Laurent
coefficient is}
\[\zeta(\tfrac23)^k\,3^{-\ell_k}\sum_{L(R)=\ell_k}C_R .\] Status:
CONJECTURE for \(k\ge4\); proved for \(k=2,3\) (T45, T46k, T46s-ii).

\textbf{Depends on:} ANS-T46s.

\hypertarget{i.-collective-growth-at-and-below-s-1}{%
\section{I. Collective growth at and below s =
1}\label{i.-collective-growth-at-and-below-s-1}}

\hypertarget{ans-t48-growth-of-g_k-at-s-1-of-exact-order-log-k}{%
\subsection{ANS-T48: Growth of G\_k at s = 1 of exact order log
k}\label{ans-t48-growth-of-g_k-at-s-1-of-exact-order-log-k}}

\textbf{Proof status:} PROVED. \textbf{Review:} UNNS Research Collective
review: valid.

\textbf{Review record:} Release review records. \textbf{Release and
computation:} Archived release.

\textbf{Source:} ANS-I v014, lines 147-149. \textbf{Scope:} Selected
nesting grammar. \textbf{Literature priority:} Unreviewed.

\textbf{Release table summary} (ANS-I/v014): G\_k(1)\^{}(1/k) of exact
order log k; lower bound k\^{}(1-s) for 1/2\textless s\textless1

\textbf{Statement as recorded:}

\textbf{Theorem.} (i) As \(k\to\infty\),
\[\log k-\log\log k+O(1)\ \le\ G_k(1)^{1/k}\ \le\ \big(e^{\gamma}\zeta(2)+o(1)\big)\log k .\]
(ii) For \(\frac12<s<1\) and large \(k\),
\(G_k(s)^{1/k}\ge c_s\,k^{1-s}\).

\textbf{Depends on:} ANS-T23, ANS-T23a. \textbf{Used by:} ANS-T48a,
ANS-T48b, ANS-T49, ANS-T50. \textbf{Classical inputs:} Ville, Mertens,
Chebyshev.

\textbf{Relations:} ANS-T48a sharpens the upper constant of ANS-T48;
ANS-T49 raises the lower constant of ANS-T48; ANS-T50 extends to a
window the setting of ANS-T48.

\hypertarget{ans-t48a-upper-constant-egamma}{%
\subsection{ANS-T48a: Upper constant
e\^{}gamma}\label{ans-t48a-upper-constant-egamma}}

\textbf{Proof status:} PROVED. \textbf{Review:} UNNS Research Collective
review: valid.

\textbf{Review record:} UNNS Research Collective review, 27 November
2026 (PROVENANCE/REVIEWS). \textbf{Release and computation:} Archived
release.

\textbf{Source:} ANS-I v015, lines 28-31. \textbf{Scope:} Selected
nesting grammar. \textbf{Literature priority:} Unreviewed.

\textbf{Release table summary} (ANS-I/v015): limsup G\_k(1)\^{}(1/k)/log
k \textless= e\^{}gamma

\textbf{Statement as recorded:}

For \(a,b>0\) put
\(D_{a,b}(E)=\sum_{d\mid E}d^{-a}\prod_{p\mid d}(1-p^{-b})\). The step
estimate of v014 §5.2 holds in the sharper form
\[\sum_Nr_N^{-a}\Big(\frac{\gcd(E,N)}N\Big)^b\le\zeta(a+b)\,D_{a,b}(E),\]
because \(J_b(d)=d^b\prod_{p\mid d}(1-p^{-b})\) was previously bounded
by \(d^b\). The function \(D_{a,b}\) is multiplicative and
\[D_{a,b}(p^{\ell})=1+(1-p^{-b})\sum_{j=1}^{\ell}p^{-aj}\le F_p=\frac{1-p^{-a-b}}{1-p^{-a}} .\]

\textbf{Depends on:} ANS-T48. \textbf{Used by:} ANS-T50.
\textbf{Classical inputs:} Mertens.

\textbf{Relations:} ANS-T48a sharpens the upper constant of ANS-T48.

\hypertarget{ans-t48b-exponential-history-function-order-one-infinite-type}{%
\subsection{ANS-T48b: Exponential history function: order one, infinite
type}\label{ans-t48b-exponential-history-function-order-one-infinite-type}}

\textbf{Proof status:} PROVED. \textbf{Review:} UNNS Research Collective
review: valid, with a recorded qualification. The bounding constants
satisfy 0 \textless{} c\_- \textless= c\_+ \textless{} infinity.

\textbf{Review record:} UNNS Research Collective review, 27 November
2026 (PROVENANCE/REVIEWS). \textbf{Release and computation:} Archived
release.

\textbf{Source:} ANS-I v015, lines 36-39. \textbf{Scope:} Selected
nesting grammar. \textbf{Literature priority:} Unreviewed.

\textbf{Release table summary} (ANS-I/v015): Exponential history
function: order 1, infinite type, maximum-modulus bounds

\textbf{Statement as recorded:}

Let \(\mathcal B_1(t)=\sum_kG_k(1)t^k/k!\). Suppose
\[c_-\le\liminf\frac{G_k(1)^{1/k}}{\log k},\qquad\limsup\frac{G_k(1)^{1/k}}{\log k}\le c_+ .\]
Then \(\mathcal B_1\) is entire of order \(1\) and infinite type, and
\[c_-\le\liminf_{r\to\infty}\frac{\log\mathcal B_1(r)}{r\log r}\le\limsup_{r\to\infty}\frac{\log\mathcal B_1(r)}{r\log r}\le c_+ .\]

\textbf{Depends on:} ANS-T48. \textbf{Classical inputs:} Stirling.

\hypertarget{ans-t49-lower-constant-1.1079}{%
\subsection{ANS-T49: Lower constant
1.1079}\label{ans-t49-lower-constant-1.1079}}

\textbf{Proof status:} PROVED. \textbf{Review:} UNNS Research Collective
review: valid.

\textbf{Review record:} UNNS Research Collective review, 27 November
2026 (PROVENANCE/REVIEWS). \textbf{Release and computation:} Archived
release.

\textbf{Source:} ANS-I v015, lines 47-52, 56-59. \textbf{Scope:}
Selected nesting grammar. \textbf{Literature priority:} Unreviewed.

\textbf{Release table summary} (ANS-I/v015): liminf G\_k(1)\^{}(1/k)/log
k \textgreater= 1 + beta* I(1/beta*) = 1.1079\ldots{}

\textbf{Statement as recorded:}

Let \(\rho\) be the Dickman function and
\(I(v)=\int_v^{\infty}\rho(u)\,du\). We use two classical facts:
\(\int_0^{\infty}\rho=e^{\gamma}\), and for fixed \(u\ge0\) the sum of
\(1/m\) over \(z\)-smooth \(m\le z^u\) is
\(\big(\int_0^u\rho+o(1)\big)\log z\). With Mertens' theorem this gives,
uniformly for \(t=z^v\) with \(v\) in a bounded range,
\[T_z(t):=\sum_{m>t,\;P^+(m)\le z}\frac1m=\big(I(v)+o(1)\big)\log z.\qquad(1)\]

\textbf{Theorem.} Put \(\beta_*=2/(2+e^{\gamma})=0.52895\ldots\) Then
\[\liminf_{k\to\infty}\frac{G_k(1)^{1/k}}{\log k}\ \ge\ 1+\beta_*\,I(1/\beta_*)=1.1079\ldots,\]
where, since \(1<1/\beta_*<2\), \(I(v)=e^{\gamma}+1-v(2-\log v)\) at
\(v=1/\beta_*\).

Fix \(\beta<\beta_*\) and put \(Y=\lfloor k/\log k\rfloor\),
\(z=Y^{\beta}\). Let
\[\mathcal A=\{N\le Y\}\ \cup\ \{N>Y:\ P^+(N)\le z\},\qquad\mu(N)=r_N^{-1}=\frac N{(N+1)^2}.\]
Call the two parts \emph{base} and \emph{extra}. Only primes \(p\le Y\)
divide letters of \(\mathcal A\). The extra letters are unbounded. By
(1) with \(v=1/\beta\),
\[B_{\mathcal A}=\sum_{N\in\mathcal A}\mu(N)=\big(1+\beta I(1/\beta)+o(1)\big)\log Y.\]

\textbf{Depends on:} ANS-T48. \textbf{Classical inputs:} Dickman
function, Mertens.

\textbf{Relations:} ANS-T49 raises the lower constant of ANS-T48.

\hypertarget{ans-t50-growth-bounds-in-the-window-around-s-1}{%
\subsection{ANS-T50: Growth bounds in the window around s =
1}\label{ans-t50-growth-bounds-in-the-window-around-s-1}}

\textbf{Proof status:} PROVED. \textbf{Review:} UNNS Research Collective
review: valid, with a recorded qualification. Fixed real tau as k tends
to infinity; no uniform statement for tau = tau(k).

\textbf{Review record:} UNNS Research Collective review, 27 November
2026 (PROVENANCE/REVIEWS). \textbf{Release and computation:} Archived
release.

\textbf{Source:} ANS-I v015, lines 94-99. \textbf{Scope:} Selected
nesting grammar. \textbf{Literature priority:} Unreviewed.

\textbf{Release table summary} (ANS-I/v015): Growth bounds in the window
s = 1 + tau/log k

\textbf{Statement as recorded:}

Let \(\mathrm{Ein}(\tau)=\int_0^1(1-e^{-\tau u})\,du/u\).

\textbf{Theorem.} For every real \(\tau\), with \(s_k=1+\tau/\log k\),
\[\int_0^1e^{-\tau u}du\ \le\ \liminf_k\frac{G_k(s_k)^{1/k}}{\log k}\ \le\ \limsup_k\frac{G_k(s_k)^{1/k}}{\log k}\ \le\ e^{\gamma-\mathrm{Ein}(\tau)} .\]

At \(\tau=0\) these are \(1\) and \(e^{\gamma}\). As \(\tau\to+\infty\)
the left side is \((1-e^{-\tau})/\tau\) and the right side is
\(e^{-E_1(\tau)}/\tau\), with \(E_1\) the exponential integral. Both
approach \(1/\tau\).

\textbf{Depends on:} ANS-T48, ANS-T48a. \textbf{Classical inputs:}
Mertens.

\textbf{Relations:} ANS-T50 extends to a window the setting of ANS-T48.

\hypertarget{j.-extended-generation-torsion-multipliers}{%
\section{J. Extended generation: torsion
multipliers}\label{j.-extended-generation-torsion-multipliers}}

\textbf{Setting (ANS-II v001).}

\textbf{Extension.} \(A=(\alpha_1,\dots,\alpha_m)\) is a finite list of
multipliers \(\alpha_j>1\) with a rational power,
\(\alpha_j=\prod_pp^{w_{jp}}\), \(w_{jp}\in\mathbb Q\). Letters are the
nesting labels \(N\ge1\) (multiplier \(r_N\)) and the symbols
\(\alpha_j\). Write \(\mu(\ell)\) for the multiplier of a letter and
\(X_p(\ell)\) for its exponent at \(p\): \(X_p(N)=2v_p(N+1)-v_p(N)\),
\(X_p(\alpha_j)=w_{jp}\). Let \(P_0\) be the finite set of primes with
some \(w_{jp}\ne0\), and \(R=\prod_{p\in P_0}p\).

\textbf{States and admissibility {[}TC{]}.} States are
\(x=a_\varphi n\), \(n\in\mathbb N\), where
\(a_\varphi=\prod p^{\varphi_p}\), \(0\le\varphi_p<1\), and \(\varphi\)
runs over the finite group \(\Phi\) generated by the vectors \(w_j\)
modulo \(\mathbb Z\). A letter is admissible at \(x\) when
\(\mu(\ell)x\) again has all prime exponents non-negative (integrality
preservation). For a nesting letter this is \(N\mid n\).

\textbf{Histories.} For a word \(w\) and a class \(\varphi\),
\(K_\varphi(w)\) is the least \(n\) such that \(w\) is admissible from
\(a_\varphi n\); \(Q(w)=\prod\mu\);
\(E_\varphi(w)=a_\varphi K_\varphi(w)Q(w)\). Put
\(\mathcal G_\varphi(s,t)=\sum_wt^{|w|}E_\varphi(w)^{-s}\),
\(A_\varphi(X)=\#\{w:E_\varphi(w)\le X\}\), and let
\(\mathcal Z_\varphi(s,t)\) be the sum of
\(t^{|w|}(\text{endpoint})^{-s}\) over all pairs (source
\(a_\varphi n\), admissible \(w\)).

\textbf{Characteristic function.} For \(a>1\),
\(B_A(a)=\sum_N r_N^{-a}+\sum_j\alpha_j^{-a}\). It decreases from
\(\infty\) to \(0\); \(\sigma_A\) is the root of \(B_A=1\). The letter
law is \(\nu_a(\ell)=\mu(\ell)^{-a}/B_A(a)\) and
\(\mu_p(a)=\mathbb E_{\nu_a}X_p\) is the \textbf{drift} at \(p\).

\textbf{Hypothesis (D).} \(\mu_p(\sigma_A)>0\) for every \(p\in P_0\).

\hypertarget{ans2-t00-classwise-invariance-on-valuation-type-classes}{%
\subsection{ANS2-T00: Classwise invariance on valuation-type
classes}\label{ans2-t00-classwise-invariance-on-valuation-type-classes}}

\textbf{Proof status:} PROVED. \textbf{Review:} UNNS Research Collective
review: valid.

\textbf{Review record:} Release review records. \textbf{Release and
computation:} Archived release.

\textbf{Source:} ANS-II v001, lines 17-17. \textbf{Scope:} Extended
grammar with torsion multipliers. \textbf{Literature priority:}
Unreviewed.

\textbf{Release table summary} (ANS-II/v001): Z\_C = alpha\_C\^{}-s Z\_H
on a valuation-type class

\textbf{Statement as recorded:}

\textbf{Theorem.} Take \(A\) empty and any base point \(\alpha_C>0\),
with admissibility \(N\mid n\) on \(\alpha_Cn\) ({[}AT{]}, classwise
decomposition). Then
\(\mathcal Z_C(s,t)=\alpha_C^{-s}\mathcal Z_H(s,t)\) and
\(N_{\alpha_Cx}(\alpha_CX)=N_x(X)\).

\textbf{Depends on:} no earlier result cited. \textbf{Used by:}
ANS2-T05.

\hypertarget{ans2-l1-source-formula-in-extended-grammars}{%
\subsection{ANS2-L1: Source formula in extended
grammars}\label{ans2-l1-source-formula-in-extended-grammars}}

\textbf{Proof status:} PROVED. \textbf{Review:} UNNS Research Collective
review: valid.

\textbf{Review record:} Release review records. \textbf{Release and
computation:} Archived release.

\textbf{Source:} ANS-II v001, lines 24-24. \textbf{Scope:} Extended
grammar with torsion multipliers. \textbf{Literature priority:}
Unreviewed.

\textbf{Release table summary} (ANS-II/v001): v\_p(K) = max\_j
-floor(phi\_p + S\_p(j))

\textbf{Statement as recorded:}

\textbf{Lemma 1 (ANS2-L1).} With \(S_p(j)\) the sum of \(X_p\) over the
first \(j\) letters,
\(v_p(K_\varphi(w))=\max_{0\le j\le|w|}\big(-\lfloor\varphi_p+S_p(j)\rfloor\big)\),
and \(w\) is admissible from \(a_\varphi n\) exactly when
\(K_\varphi(w)\mid n\).

\textbf{Depends on:} no earlier result cited. \textbf{Used by:} ANS2-L2,
ANS2-T02, ANS2-T05, ANS2-T16, ANS2N-T04b.

\hypertarget{ans2-l2-comparison-of-torsion-classes}{%
\subsection{ANS2-L2: Comparison of torsion
classes}\label{ans2-l2-comparison-of-torsion-classes}}

\textbf{Proof status:} PROVED. \textbf{Review:} UNNS Research Collective
review: valid.

\textbf{Review record:} Release review records. \textbf{Release and
computation:} Archived release.

\textbf{Source:} ANS-II v001, lines 28-30. \textbf{Scope:} Extended
grammar with torsion multipliers. \textbf{Literature priority:}
Unreviewed.

\textbf{Release table summary} (ANS-II/v001): All torsion classes are
comparable, K\_phi / K\_0 / R K\_phi

\textbf{Statement as recorded:}

\textbf{Lemma 2 (ANS2-L2).}
\(K_\varphi(w)\mid K_0(w)\mid R\,K_\varphi(w)\). Hence for real \(a>0\)
and each depth \(k\),
\[a_\varphi^{-a}\,\mathcal G_{0,k}(a)\le\mathcal G_{\varphi,k}(a)\le(R/a_\varphi)^{a}\,\mathcal G_{0,k}(a).\]

\textbf{Depends on:} ANS2-L1. \textbf{Used by:} ANS2-T16.

\hypertarget{ans2-t01-factorization-in-torsion-classes}{%
\subsection{ANS2-T01: Factorization in torsion
classes}\label{ans2-t01-factorization-in-torsion-classes}}

\textbf{Proof status:} PROVED. \textbf{Review:} UNNS Research Collective
review: valid.

\textbf{Review record:} Release review records. \textbf{Release and
computation:} Archived release.

\textbf{Source:} ANS-II v001, lines 35-35. \textbf{Scope:} Extended
grammar with torsion multipliers. \textbf{Literature priority:}
Unreviewed.

\textbf{Release table summary} (ANS-II/v001): Z\_phi = zeta G\_phi

\textbf{Statement as recorded:}

\textbf{Factorization (ANS2-T01).}
\(\mathcal Z_\varphi(s,t)=\zeta(s)\,\mathcal G_\varphi(s,t)\) where
convergent: the admissible sources are \(n=K_\varphi(w)m\), with
endpoint \(E_\varphi(w)m\).

\textbf{Depends on:} no earlier result cited.

\hypertarget{ans2-t02-positive-drift-depth-bounds-and-abscissa}{%
\subsection{ANS2-T02: Positive drift: depth bounds and
abscissa}\label{ans2-t02-positive-drift-depth-bounds-and-abscissa}}

\textbf{Proof status:} PROVED. \textbf{Review:} UNNS Research Collective
review: valid.

\textbf{Review record:} Release review records. \textbf{Release and
computation:} Archived release.

\textbf{Source:} ANS-II v001, lines 43-43, 50-51, 57-57. \textbf{Scope:}
Extended grammar with torsion multipliers. \textbf{Literature priority:}
Unreviewed.

\textbf{Release table summary} (ANS-II/v001): Depth bounds k\^{}beta/log
k to k\^{}beta; abscissa sigma\_A; critical convergence

\textbf{Statement as recorded:}

\textbf{Lemma 3.} Under (D) there are a neighbourhood \(I\) of
\(\sigma_A\), \(\eta>0\) and \(M_0\) such that
\(\mathbb E\,e^{-\eta X_p}\le1\) for all primes \(p\), all \(a\in I\),
and letters restricted to \(N\le M\) (all \(\alpha_j\) kept),
\(M_0\le M\le\infty\). If all \(w_{jp}\ge0\), then (D) holds and every
compact interval \(I\subset(1,\infty)\) may be used.

\textbf{Theorem.} Under (D), for \(a\in I\), \(\beta=1/a\), \(b>0\) and
large \(k\),
\[\frac{c\,k^{\beta}}{\log k}\le-\log\mathbb E\,K_{\varphi,k}^{-b}\le C\,k^{\beta}.\]

\textbf{Corollaries} (under (D)). (i)
\(\mathcal G_{\varphi,k}(a)^{1/k}\to B_A(a)\), and the abscissa of
\(\mathcal G_\varphi(\cdot,1)\) is exactly \(\sigma_A\); for divergence
below \(\sigma_A\) use \(a\in I\) and monotonicity (\(E_\varphi\ge1\)).
The counting statement \(A_\varphi(X)=X^{\sigma_A+o(1)}\) does not
follow from the abscissa alone; it is a consequence of the counting
theorem of §5 (whose proof uses only §§2--4). (ii)
\(K_{\varphi,k}\to\infty\) almost surely at \(a\in I\). (iii) On
\(|t|B_A(a)=1\), \(a=\operatorname{Re}s\in I\), the series
\(\sum_wt^{|w|}E_\varphi^{-s}K_\varphi^{-u}\) converges absolutely iff
\(\operatorname{Re}(s+u)>0\). (iv)
\(\mathcal G_\varphi(\sigma_A,1)<\infty\).

\textbf{Depends on:} ANS2-L1, ANS-T23. \textbf{Used by:} ANS2-T03,
ANS2-T08, ANS2-T09a, ANS2-T16. \textbf{Classical inputs:} Jensen, Ville.

\textbf{Relations:} ANS2-T16 treats the abscissa for any drift signs,
beyond ANS2-T02.

\hypertarget{ans2-t03-positive-drift-counting-functions}{%
\subsection{ANS2-T03: Positive drift: counting
functions}\label{ans2-t03-positive-drift-counting-functions}}

\textbf{Proof status:} PROVED. \textbf{Review:} UNNS Research Collective
review: valid.

\textbf{Review record:} Release review records. \textbf{Release and
computation:} Archived release.

\textbf{Source:} ANS-II v001, lines 60-61. \textbf{Scope:} Extended
grammar with torsion multipliers. \textbf{Literature priority:}
Unreviewed.

\textbf{Release table summary} (ANS-II/v001): X\^{}sigma\_A exp(-(log
X)\^{}beta) up to log log, all counting functions

\textbf{Statement as recorded:}

\textbf{Theorem.} Under (D), with \(\sigma=\sigma_A\),
\(\beta=1/\sigma\), and \(Z(X)\) equal to \(A_\varphi(X)\), to the
number \(N_x(X)\) of histories from a fixed state \(x=a_\varphi n\) with
endpoint at most \(X\), or to the number
\(C_\varphi(X)=\sum_{m\le X}A_\varphi(X/m)\) from all sources,
\[\frac{c(\log X)^{\beta}}{\log\log X}\le-\log\frac{Z(X)}{X^{\sigma}}\le C(\log X)^{\beta}.\]

\textbf{Depends on:} ANS2-T02, ANS-T24a. \textbf{Used by:} ANS2-T04.

\hypertarget{ans2-t04-positive-drift-derivatives-and-infimal-gevrey-order}{%
\subsection{ANS2-T04: Positive drift: derivatives and infimal Gevrey
order}\label{ans2-t04-positive-drift-derivatives-and-infimal-gevrey-order}}

\textbf{Proof status:} PROVED. \textbf{Review:} UNNS Research Collective
review: valid, with a recorded qualification. Valid for the infimal
Gevrey order.

\textbf{Review record:} Release review records. \textbf{Release and
computation:} Archived release.

\textbf{Source:} ANS-II v001, lines 80-80. \textbf{Scope:} Extended
grammar with torsion multipliers. \textbf{Literature priority:}
Unreviewed.

\textbf{Release table summary} (ANS-II/v001): Derivative bounds; smooth
singular point; infimal Gevrey order sigma\_A (membership at order
sigma\_A open)

\textbf{Statement as recorded:}

\textbf{Theorem.} Under (D),
\(D_j=\sum_w(\log E_\varphi)^jE_\varphi^{-\sigma_A}\) satisfies
\(c^jj^{\sigma_Aj}\le D_j\le(Cj\log j)^{\sigma_Aj}\). Consequently
\(\mathcal G_\varphi(\cdot,1)\) and \(\mathcal Z_\varphi(\cdot,1)\)
extend infinitely differentiably to \(\operatorname{Re}s\ge\sigma_A\);
\(\sigma_A\) is a singular point at which neither has a meromorphic
continuation; and the infimal Gevrey order of
\(\mathcal G_\varphi(\cdot,1)\) on the closed half-plane is
\(\sigma_A\).

\textbf{Depends on:} ANS2-T03, ANS-T27. \textbf{Classical inputs:}
Landau.

\hypertarget{ans2-t05-negative-drift-abscissa-below-the-characteristic-root}{%
\subsection{ANS2-T05: Negative drift: abscissa below the characteristic
root}\label{ans2-t05-negative-drift-abscissa-below-the-characteristic-root}}

\textbf{Proof status:} PROVED. \textbf{Review:} UNNS Research Collective
review: valid.

\textbf{Review record:} Release review records. \textbf{Release and
computation:} Archived release.

\textbf{Source:} ANS-II v001, lines 91-93. \textbf{Scope:} Extended
grammar with torsion multipliers. \textbf{Literature priority:}
Unreviewed.

\textbf{Release table summary} (ANS-II/v001): Abscissa strictly below
the characteristic root

\textbf{Statement as recorded:}

\textbf{Theorem.} Suppose \(\mu_p(\sigma_A)<0\) for some \(p\in P_0\).
Then the abscissa \(\sigma^\ast\) of \(\mathcal G_\varphi(\cdot,1)\)
satisfies \(\sigma_c\le\sigma^\ast<\sigma_A\), where \(\sigma_c\) is the
nesting exponent. Moreover \(\sigma^\ast\le a\) whenever
\(\Phi_p(a,\theta)<1\) for some \(\theta\in[0,1]\), where
\[\Phi_p(a,\theta)=\sum_\ell\mu(\ell)^{-a}p^{\theta aX_p(\ell)}.\]

\textbf{Depends on:} ANS2-L1, ANS2-T00.

\textbf{Relations:} ANS2-T16 gives the exact value bounded in ANS2-T05.

\hypertarget{ans2-t06-first-letter-identity-between-classes}{%
\subsection{ANS2-T06: First-letter identity between
classes}\label{ans2-t06-first-letter-identity-between-classes}}

\textbf{Proof status:} PROVED. \textbf{Review:} UNNS Research Collective
review: valid.

\textbf{Review record:} Release review records. \textbf{Release and
computation:} Archived release.

\textbf{Source:} ANS-II v001, lines 99-101. \textbf{Scope:} Extended
grammar with torsion multipliers. \textbf{Literature priority:}
Unreviewed.

\textbf{Release table summary} (ANS-II/v001): First-letter identity
between classes

\textbf{Statement as recorded:}

\textbf{Proposition.} Let all \(w_{jp}\ge0\) for \(\alpha=\alpha_j\),
applied in class \(\varphi\), with image class \(\psi\) and integer
carry \(c\) (\(\alpha a_\varphi=a_\psi c\)). Then, with \(K'=K_\psi(w)\)
and \(g=\gcd(K',c)\),
\[K_\varphi(\alpha w)=K'/g,\qquad E_\varphi(\alpha w)=E_\psi(w)\cdot c/g .\]

\textbf{Depends on:} no earlier result cited.

\hypertarget{ans2-t07-first-introductions-in-extended-grammars}{%
\subsection{ANS2-T07: First introductions in extended
grammars}\label{ans2-t07-first-introductions-in-extended-grammars}}

\textbf{Proof status:} PROVED. \textbf{Review:} UNNS Research Collective
review: valid, with a recorded qualification. Valid with the precise
event definition, now in the text.

\textbf{Review record:} Release review records. \textbf{Release and
computation:} Archived release.

\textbf{Source:} ANS-II v001, lines 84-86. \textbf{Scope:} Extended
grammar with torsion multipliers. \textbf{Literature priority:}
Unreviewed.

\textbf{Release table summary} (ANS-II/v001): Event I\_k (every prime
dividing a nesting letter occurred earlier through N or N+1):
probability exp(-order k\^{}beta/log k)

\textbf{Statement as recorded:}

\textbf{First introductions (ANS2-T07).} A prime \emph{occurs} in a
nesting letter \(N\) if it divides \(N(N+1)\); it has been
\emph{encountered before step \(i\)} if it occurs in some nesting letter
at a step \(j<i\), whether through \(N_j\) or through \(N_j+1\).
Adjoined letters are ignored by this observable. Let \(\mathcal I_k\) be
the event that, for every \(i\le k\) at which the letter is a nesting
letter \(N_i\), every prime dividing \(N_i\) has been encountered before
step \(i\).

\textbf{Theorem.} For every \(a>1\),
\(\Pr(\mathcal I_k)=\exp(-\Theta(k^\beta/\log k))\).

\textbf{Depends on:} ANS-T28a.

\hypertarget{ans2-t08-dependence-on-the-multiplier}{%
\subsection{ANS2-T08: Dependence on the
multiplier}\label{ans2-t08-dependence-on-the-multiplier}}

\textbf{Proof status:} PROVED. \textbf{Review:} UNNS Research Collective
review: valid.

\textbf{Review record:} Release review records. \textbf{Release and
computation:} Archived release.

\textbf{Source:} ANS-II v001, lines 108-110. \textbf{Scope:} Extended
grammar with torsion multipliers. \textbf{Literature priority:}
Unreviewed.

\textbf{Release table summary} (ANS-II/v001): Dependence on multiplier
size; 2\^{}(1/d) asymptotic

\textbf{Statement as recorded:}

\begin{enumerate}
\def\labelenumi{(\roman{enumi})}
\tightlist
\item
  Under (D), \(\sigma_A\) is determined by the real values of the
  multipliers, and adjoining a multiplier increases the root of
  \(B_A=1\). (ii) The torsion order does not determine it: \(\sqrt2\)
  and \(2^{3/2}\) both have order 2 and different roots. (iii) For
  \(\alpha=2^{1/d}\),
  \[\sigma_d=\log_4d-\log_4\log d+O(1),\qquad\beta_d=1/\sigma_d\to0 .\]
\end{enumerate}

\textbf{Depends on:} ANS2-T02.

\hypertarget{ans2-t14-one-letter-formula-for-any-signs}{%
\subsection{ANS2-T14: One-letter formula for any
signs}\label{ans2-t14-one-letter-formula-for-any-signs}}

\textbf{Proof status:} PROVED. \textbf{Review:} UNNS Research Collective
review: valid.

\textbf{Review record:} Release review records. \textbf{Release and
computation:} Archived release.

\textbf{Source:} ANS-II v002, lines 9-11. \textbf{Scope:} Extended
grammar with torsion multipliers. \textbf{Literature priority:}
Unreviewed.

\textbf{Release table summary} (ANS-II/v002): One letter formula and
operator for multipliers of any signs

\textbf{Statement as recorded:}

A letter \(\ell\) takes class \(\psi\) to a class \(\psi'\) with
\(\mu(\ell)a_\psi=a_{\psi'}\,u/v\), where \(u/v\) is reduced. Here
\(u,v\) depend on \((\psi,\ell)\). For a nesting letter \(\psi'=\psi\),
\(u=(N+1)^2\), \(v=N\). For fixed \(\ell\) the map \(\psi\mapsto\psi'\)
is a bijection of \(\Phi_A\). No sign condition on the exponents is
used.

\textbf{Lemma 1.}
\(\displaystyle n(w\ell)=\frac{u\,n(w)}{\gcd(n(w),v)}\) and
\(E_\varphi(w\ell)=E_\varphi(w)\,\mu(\ell)\,\dfrac{v}{\gcd(n(w),v)}\).

\textbf{Depends on:} no earlier result cited. \textbf{Used by:}
ANS2-T10, ANS2-T11, ANS2-T15.

\hypertarget{ans2-t10-divisibility-pullback}{%
\subsection{ANS2-T10: Divisibility
pullback}\label{ans2-t10-divisibility-pullback}}

\textbf{Proof status:} PROVED. \textbf{Review:} UNNS Research Collective
review: valid.

\textbf{Review record:} Release review records. \textbf{Release and
computation:} Archived release.

\textbf{Source:} ANS-II v002, lines 15-17. \textbf{Scope:} Extended
grammar with torsion multipliers. \textbf{Literature priority:}
Unreviewed.

\textbf{Release table summary} (ANS-II/v002): Divisibility pullback d
divides n' iff D(d) divides n

\textbf{Statement as recorded:}

\textbf{Lemma 2 (divisibility pullback).} For \(d\ge1\) put
\[D(d)=\prod_{p:\,v_p(d)>v_p(u)}p^{\,v_p(v)+v_p(d)-v_p(u)}.\] Then
\(d\mid n(w\ell)\) if and only if \(D(d)\mid n(w)\). For a nesting
letter this is
\(D_N(d)=\prod_{v_p d>2v_p(N+1)}p^{\,v_pN+v_pd-2v_p(N+1)}\).

\textbf{Depends on:} ANS2-T14. \textbf{Used by:} ANS2-T11.

\hypertarget{ans2-t11-nesting-transfer-operator}{%
\subsection{ANS2-T11: Nesting transfer
operator}\label{ans2-t11-nesting-transfer-operator}}

\textbf{Proof status:} PROVED. \textbf{Review:} UNNS Research Collective
review: valid.

\textbf{Review record:} Release review records. \textbf{Release and
computation:} Archived release.

\textbf{Source:} ANS-II v002, lines 22-28. \textbf{Scope:} Extended
grammar with torsion multipliers. \textbf{Literature priority:}
Unreviewed.

\textbf{Release table summary} (ANS-II/v002): Nesting transfer operator

\textbf{Statement as recorded:}

Define, for \(\psi\in\Phi_A\), \(d\ge1\),
\[F^{(k)}_{\psi,d}(s)=\sum_{|w|=k,\ \psi(w)=\psi,\ d\mid n(w)}E_\varphi(w)^{-s},\]
and let \(J_s(h)=\sum_{e\mid h}\mu_{\rm M}(e)(h/e)^s\) (\(\mu_{\rm M}\)
the Möbius function).

\textbf{Theorem.} For \(\operatorname{Re}s>1\), \(F^{(k+1)}=L_sF^{(k)}\)
with
\[(L_sF)_{\psi',d}=\sum_{\ell}\mu(\ell)^{-s}v^{-s}\sum_{h\mid v}J_s(h)\,F_{\psi,\,\operatorname{lcm}(h,D(d))},\]
where for each \(\ell\), \(\psi\) is the class carried to \(\psi'\) and
\(u,v,D\) belong to \((\psi,\ell)\). The initial vector is
\(F^{(0)}_{\psi,d}=a_\varphi^{-s}[\psi=\varphi][d=1]\), and
\(\mathcal G_{\varphi,k}=\sum_\psi F^{(k)}_{\psi,1}\).

\textbf{Depends on:} ANS2-T10, ANS2-T14. \textbf{Used by:} ANS2-T12,
ANS2-T13.

\hypertarget{ans2-t12-torsion-carry-operator}{%
\subsection{ANS2-T12: Torsion carry
operator}\label{ans2-t12-torsion-carry-operator}}

\textbf{Proof status:} PROVED. \textbf{Review:} UNNS Research Collective
review: valid.

\textbf{Review record:} Release review records. \textbf{Release and
computation:} Archived release.

\textbf{Source:} ANS-II v002, lines 32-34. \textbf{Scope:} Extended
grammar with torsion multipliers. \textbf{Literature priority:}
Unreviewed.

\textbf{Release table summary} (ANS-II/v002): Torsion carry operator,
non-negative exponents

\textbf{Statement as recorded:}

\textbf{The two kinds of letter.} Nesting letters keep the class: with
\(v=N\), \(\mu^{-s}v^{-s}=(N+1)^{-2s}\), giving
\[(T_NF)_{\psi,d}=(N+1)^{-2s}\sum_{h\mid N}J_s(h)F_{\psi,\operatorname{lcm}(h,D_N(d))}.\]
A torsion letter with non-negative exponents has \(v=1\) and integer
carry \(u=c\), giving the pure shift
\((S_\alpha F)_{\psi',d}=\alpha^{-s}F_{\psi,\,d/\gcd(d,c)}\).

\textbf{Depends on:} ANS2-T11. \textbf{Used by:} ANS2-T13.

\hypertarget{ans2-t13-closed-countable-operator-system}{%
\subsection{ANS2-T13: Closed countable operator
system}\label{ans2-t13-closed-countable-operator-system}}

\textbf{Proof status:} PROVED. \textbf{Review:} UNNS Research Collective
review: valid.

\textbf{Review record:} Release review records. \textbf{Release and
computation:} Archived release.

\textbf{Source:} ANS-II v002, lines 22-28, 32-34. \textbf{Scope:}
Extended grammar with torsion multipliers. \textbf{Literature priority:}
Unreviewed.

\textbf{Release table summary} (ANS-II/v002): Closed countable operator
system, Re s \textgreater{} 1

\textbf{Statement as recorded:}

Define, for \(\psi\in\Phi_A\), \(d\ge1\),
\[F^{(k)}_{\psi,d}(s)=\sum_{|w|=k,\ \psi(w)=\psi,\ d\mid n(w)}E_\varphi(w)^{-s},\]
and let \(J_s(h)=\sum_{e\mid h}\mu_{\rm M}(e)(h/e)^s\) (\(\mu_{\rm M}\)
the Möbius function).

\textbf{Theorem.} For \(\operatorname{Re}s>1\), \(F^{(k+1)}=L_sF^{(k)}\)
with
\[(L_sF)_{\psi',d}=\sum_{\ell}\mu(\ell)^{-s}v^{-s}\sum_{h\mid v}J_s(h)\,F_{\psi,\,\operatorname{lcm}(h,D(d))},\]
where for each \(\ell\), \(\psi\) is the class carried to \(\psi'\) and
\(u,v,D\) belong to \((\psi,\ell)\). The initial vector is
\(F^{(0)}_{\psi,d}=a_\varphi^{-s}[\psi=\varphi][d=1]\), and
\(\mathcal G_{\varphi,k}=\sum_\psi F^{(k)}_{\psi,1}\).

\textbf{The two kinds of letter.} Nesting letters keep the class: with
\(v=N\), \(\mu^{-s}v^{-s}=(N+1)^{-2s}\), giving
\[(T_NF)_{\psi,d}=(N+1)^{-2s}\sum_{h\mid N}J_s(h)F_{\psi,\operatorname{lcm}(h,D_N(d))}.\]
A torsion letter with non-negative exponents has \(v=1\) and integer
carry \(u=c\), giving the pure shift
\((S_\alpha F)_{\psi',d}=\alpha^{-s}F_{\psi,\,d/\gcd(d,c)}\).

\textbf{Depends on:} ANS2-T11, ANS2-T12.

\hypertarget{ans2-t15-uniform-and-weighted-l1-realizations}{%
\subsection{ANS2-T15: Uniform and weighted l\^{}1
realizations}\label{ans2-t15-uniform-and-weighted-l1-realizations}}

\textbf{Proof status:} PROVED. \textbf{Review:} UNNS Research Collective
review: valid, with a recorded qualification. Bounded only when the norm
is finite (in the text).

\textbf{Review record:} Release review records. \textbf{Release and
computation:} Archived release.

\textbf{Source:} ANS-II v002, lines 39-46. \textbf{Scope:} Extended
grammar with torsion multipliers. \textbf{Literature priority:}
Unreviewed.

\textbf{Release table summary} (ANS-II/v002): Realizations: uniform norm
B\_A(s); norm Phi\_A(s, theta) on weighted l\^{}1

\textbf{Statement as recorded:}

Let \(s>1\) be real.

\textbf{Uniform norm.} On \(\ell^\infty(\Phi_A\times\mathbb N)\) the
operator \(L_s\) is positive and \(L_s\mathbf 1=B_A(s)\mathbf 1\),
because \(\sum_{h\mid v}J_s(h)=v^s\). Hence \(\|L_s\|=B_A(s)\) and
\(B_A(s)\) is its spectral radius. This norm does not see the drifts.

\textbf{Measures on states.} Put
\(m_k(\psi,n)=\sum_{|w|=k,\,\psi(w)=\psi,\,n(w)=n}E_\varphi(w)^{-s}\).
Then \(m_{k+1}=Mm_k\), where \(M\) sends a unit mass at \((\psi,n)\) to
the masses \(\mu^{-s}(\gcd(n,v)/v)^s\) at \((\psi',un/\gcd(n,v))\), and
\(\mathcal G_{\varphi,k}(s)=\|m_k\|_1\). For
\(\theta=(\theta_q)\in[0,1]^{\text{primes}}\) let
\[V_\theta(\psi,n)=\prod_qq^{\,\theta_qs\,v_q(a_\psi n)}\ge1,\qquad\Phi_A(s,\theta)=\sum_\ell\mu(\ell)^{-s}\prod_qq^{\,\theta_qsX_q(\ell)}\in(0,\infty].\]

\textbf{Theorem.} \(\|M\|_{\ell^1(V_\theta)}=\Phi_A(s,\theta)\), in the
following sense: the supremum of the weighted column sums equals
\(\Phi_A(s,\theta)\); \(M\) is a bounded operator on
\(\ell^1(V_\theta)\) with this norm when \(\Phi_A(s,\theta)<\infty\),
and is unbounded when it is infinite.

\textbf{Depends on:} ANS2-T14. \textbf{Used by:} ANS2-T16.

\hypertarget{ans2-t16-growth-rate-abscissa-drift-dichotomy}{%
\subsection{ANS2-T16: Growth rate, abscissa, drift
dichotomy}\label{ans2-t16-growth-rate-abscissa-drift-dichotomy}}

\textbf{Proof status:} PROVED. \textbf{Review:} UNNS Research Collective
review: valid.

\textbf{Review record:} Release review records. \textbf{Release and
computation:} Archived release.

\textbf{Source:} ANS-II v002, lines 53-59, 79-79, 83-83. \textbf{Scope:}
Extended grammar with torsion multipliers. \textbf{Literature priority:}
Unreviewed.

\textbf{Release table summary} (ANS-II/v002): lim G\_k(a)\^{}(1/k) = min
over theta of Phi\_A(a, theta); abscissa is the root; sharp drift
dichotomy

\textbf{Statement as recorded:}

For \(a>1\) let
\[\rho_A(a)=\min_{\theta\in[0,1]^{\text{primes}}}\Phi_A(a,\theta).\] The
minimum exists: the cube is compact, \(\Phi_A\) is lower semicontinuous
(a sum of positive continuous functions of finitely many coordinates),
and \(\Phi_A(a,0)=B_A(a)<\infty\). \(\Phi_A(a,\cdot)\) is convex.

\textbf{Theorem.} For every finite list \(A\) as in §0 (any signs of the
drifts), every \(a>1\) and every class \(\varphi\),
\[\lim_{k\to\infty}\mathcal G_{\varphi,k}(a)^{1/k}=\rho_A(a).\]

\textbf{Corollary 2 (abscissa).} The abscissa \(\sigma^\ast\) of
\(\mathcal G_\varphi(s,1)=\sum_k\mathcal G_{\varphi,k}(s)\) is the
unique root of \(\rho_A(a)=1\); it is the same in every class and lies
in \([\sigma_c,\sigma_A]\). Consequently
\(\limsup\log A_\varphi(X)/\log X=\sigma^\ast\).

\textbf{Corollary 3 (sharp dichotomy).} \(\sigma^\ast=\sigma_A\) if
every drift \(\mu_p(\sigma_A)\), \(p\in P_0\), is \(\ge0\); and
\(\sigma^\ast<\sigma_A\) if one of them is negative.

\textbf{Depends on:} ANS2-L1, ANS2-L2, ANS2-T15, ANS2-T02. \textbf{Used
by:} ANS2-C16, ANS2-T09, ANS2-T09a, ANS2N-T01, ANS2N-T02, ANS2N-T04a,
ANS2N-T04b, ANS2N-T05. \textbf{Classical inputs:} Doob.

\textbf{Relations:} ANS2-T16 gives the exact value bounded in ANS2-T05;
ANS2-T16 treats the abscissa for any drift signs, beyond ANS2-T02.

\hypertarget{ans2-t09-nonnegative-drift-finiteness-at-the-abscissa}{%
\subsection{ANS2-T09: Nonnegative drift: finiteness at the
abscissa}\label{ans2-t09-nonnegative-drift-finiteness-at-the-abscissa}}

\textbf{Proof status:} PROVED. \textbf{Review:} UNNS Research Collective
review: valid.

\textbf{Review record:} Release review records. \textbf{Release and
computation:} Archived release.

\textbf{Source:} ANS-II v002, lines 92-92. \textbf{Scope:} Extended
grammar with torsion multipliers. \textbf{Literature priority:}
Unreviewed.

\textbf{Release table summary} (ANS-II/v002): Drifts \textgreater= 0
(zero allowed): G(sigma\_A, 1) finite and abscissa sigma\_A

\textbf{Statement as recorded:}

\textbf{Proposition (T09, the UNNS Research Collective review).} If one
or more drifts vanish at \(\sigma_A\) and the others are positive, then
\(\mathcal G_\varphi(\sigma_A,1)<\infty\) and the abscissa is
\(\sigma_A\).

\textbf{Depends on:} ANS2-T16.

\hypertarget{ans2-t09a-nonnegative-drift-depth-order}{%
\subsection{ANS2-T09a: Nonnegative drift: depth
order}\label{ans2-t09a-nonnegative-drift-depth-order}}

\textbf{Proof status:} PROVED. \textbf{Review:} UNNS Research Collective
review: valid.

\textbf{Review record:} Release review records. \textbf{Release and
computation:} Archived release.

\textbf{Source:} ANS-II v002, lines 98-99. \textbf{Scope:} Extended
grammar with torsion multipliers. \textbf{Literature priority:}
Unreviewed.

\textbf{Release table summary} (ANS-II/v002): Drifts \textgreater= 0 at
a \textgreater{} 1: -log E{[}K\^{}-b{]} between orders k\^{}beta / log k
and k\^{}beta (no sqrt k term)

\textbf{Statement as recorded:}

\textbf{Theorem (T09a).} Let \(a>1\), \(b>0\), \(\beta=1/a\), and
suppose \(\mu_p(a)\ge0\) for all \(p\in P_0\). Then for large \(k\)
\[\frac{c\,k^\beta}{\log k}\le-\log\mathbb E_{\nu_a}K_{\varphi,k}^{-b}\le C\,k^\beta.\]

\textbf{Depends on:} ANS2-T02, ANS2-T16, ANS-T23. \textbf{Classical
inputs:} Ville.

\hypertarget{ans2-c16-interval-for-2sqrt3-floating-point-certified-as-ans2n-t06}{%
\subsection{ANS2-C16: Interval for 2/sqrt3 (floating point; certified as
ANS2N-T06)}\label{ans2-c16-interval-for-2sqrt3-floating-point-certified-as-ans2n-t06}}

\textbf{Proof status:} COMPUTED. \textbf{Review:} UNNS Research
Collective review: valid.

\textbf{Review record:} Release review records. \textbf{Release and
computation:} Archived release; floating-point computation.

\textbf{Source:} ANS-II v002, lines 123-129. \textbf{Scope:} Extended
grammar with torsion multipliers. \textbf{Literature priority:}
Unreviewed.

\textbf{Release table summary} (ANS-II/v002): 2/sqrt3: 2.06055
\textless= abscissa \textless= 2.06072

\textbf{Statement as recorded:}

\begin{itemize}
\tightlist
\item
  \emph{Lower side.} By convexity, for any \(\theta^0\),
  \(\min_{[0,1]}\Phi_M\ge\Phi_M(\theta^0)+\sum_q\min\big(-g_q\theta^0_q,\;g_q(1-\theta^0_q)\big)\),
  \(g=\nabla\Phi_M(\theta^0)\). Since \(\min\Phi_M\le\rho_A\), a value
  \(\ge1\) at \(a\) gives \(\sigma^\ast\ge a\).
\item
  \emph{Upper side.} If all \(\theta_q\le t<1\), the nesting term of
  \(N\) is at most \(N^{-a(1-2t)}\), so
  \(\Phi_A(a,\theta)\le\Phi_M(a,\theta)+M^{1-c}/(c-1)\) with
  \(c=a(1-2t)>1\). A value \(<1\) gives \(\sigma^\ast<a\).
\end{itemize}

With \(M=2\cdot10^5\) for the lower side and \(10^6\) for the upper
side, \[2.06055\le\sigma^\ast\le2.06072.\] The computed values are
\(2.0605567\ldots\) and \(2.0607189\ldots\), rounded outward. (A
certified evaluation is in ANS-II-ND v001.)

\textbf{Depends on:} ANS2-T16. \textbf{Used by:} ANS2N-T06.

\textbf{Relations:} ANS2N-T06 certifies ANS2-C16.

Floating-point evaluation. The interval is certified by ANS2N-T06.

\hypertarget{k.-negative-drift-entropy-and-attainment}{%
\section{K. Negative drift: entropy and
attainment}\label{k.-negative-drift-entropy-and-attainment}}

\textbf{Setting (ANS-II-ND v001).}

As in ANS-II v002: letters \(\ell\) are the nesting labels \(N\ge1\) and
a finite list \(A\) of torsion multipliers; \(\mu(\ell)\) is the
multiplier and \(X_p(\ell)\) the exponent at \(p\); \(P_0\) is the set
of primes of \(A\); \(a>1\), \(B=B_A(a)\),
\[\Phi(\theta)=\Phi_A(a,\theta)=\sum_\ell\mu(\ell)^{-a}\prod_pp^{a\theta_pX_p(\ell)},\qquad\rho=\rho_A(a)=\min_{[0,1]^{\mathbb P}}\Phi,\]
\(\mathbb P\) the set of primes. For a probability \(\pi\) on letters
with \(\mathbb E_\pi|\log\mu|<\infty\) let \(H(\pi)\) be its entropy,
\(m_p(\pi)=\mathbb E_\pi X_p\), and
\[\Lambda(\pi)=\mathbb E_\pi\log\mu+\sum_p[-m_p(\pi)]_+\log p .\] Call
\(\pi\) \textbf{resource-neutral} if \(m_p(\pi)\ge0\) for every \(p\);
then \(\Lambda(\pi)=\mathbb E_\pi\log\mu\). Gibbs' inequality, for any
\(\theta\in[0,1]^{\mathbb P}\):
\[\log\Phi(\theta)\ge H(\pi)-a\,\mathbb E_\pi\log\mu+a\sum_p\theta_pm_p(\pi)\log p,\qquad(G)\]
with equality exactly when \(\pi(\ell)\) is proportional to the term of
\(\ell\). The exchange of sums is legitimate since
\(\sum_p|X_p(N)|\log p\le\log N+2\log(N+1)\).

\hypertarget{ans2n-t02-weights-above-one-never-help}{%
\subsection{ANS2N-T02: Weights above one never
help}\label{ans2n-t02-weights-above-one-never-help}}

\textbf{Proof status:} PROVED. \textbf{Review:} UNNS Research Collective
review: valid.

\textbf{Review record:} Release review records. \textbf{Release and
computation:} Archived release.

\textbf{Source:} ANS-II-ND v001, lines 17-17, 23-25. \textbf{Scope:}
Extended grammar with torsion multipliers. \textbf{Literature priority:}
Unreviewed.

\textbf{Release table summary} (ANS-II/ND/v001): Weights above 1 never
help; unique minimiser below 1; supremum attained in the limit by
resource-neutral laws; deficits are strictly suboptimal

\textbf{Statement as recorded:}

\textbf{Lemma 1.} If \(\theta\ge0\) and \(\Phi_M(\theta)\le B\), then
\(\theta_q\le R=2c\) for all primes \(q\le M+1\), where
\(c=1+\max(0,\log B)/(2a\log2)\).

So \(\Phi_M\) attains its minimum \(L_M\) over the whole orthant
\([0,\infty)^{\{p\le M+1\}}\) at some \(\theta^M\in[0,R]\).

\textbf{Theorem 1.} (i) \(\Phi(\theta)\ge\rho\) for every
\(\theta\in[0,\infty)^{\mathbb P}\): the upper restriction
\(\theta_p\le1\) is never active. (ii) The minimiser \(\theta^\ast\) of
\(\Phi\) on the cube is unique, and \(\theta^\ast_p<1\) for every \(p\).
(iii) \(L_M\uparrow\rho\) and \(\theta^M\to\theta^\ast\) coordinatewise.

\textbf{Depends on:} ANS2-T16. \textbf{Used by:} ANS2N-T01, ANS2N-T03,
ANS2N-T05, ANS2N-T09, ANS2N-T15.

\hypertarget{ans2n-t01-entropy-admissibility-variational-principle}{%
\subsection{ANS2N-T01: Entropy-admissibility variational
principle}\label{ans2n-t01-entropy-admissibility-variational-principle}}

\textbf{Proof status:} PROVED. \textbf{Review:} UNNS Research Collective
review: valid.

\textbf{Review record:} Release review records. \textbf{Release and
computation:} Archived release.

\textbf{Source:} ANS-II-ND v001, lines 36-46. \textbf{Scope:} Extended
grammar with torsion multipliers. \textbf{Literature priority:}
Unreviewed.

\textbf{Release table summary} (ANS-II/ND/v001): log rho = sup {[}H - a
Lambda{]}; abscissa = sup H / Lambda

\textbf{Statement as recorded:}

\textbf{Theorem 2.}
\[\log\rho=\sup_{\pi}\big[H(\pi)-a\Lambda(\pi)\big]=\sup_{\pi\ \text{resource-neutral}}\big[H(\pi)-a\,\mathbb E_\pi\log\mu\big],\]
both over finitely supported \(\pi\). For multipliers \(>1\), the
abscissa \(\sigma^\ast\) satisfies
\[\sigma^\ast=\sup_\pi\frac{H(\pi)}{\Lambda(\pi)}=\sup_{\pi\ \text{resource-neutral}}\frac{H(\pi)}{\mathbb E_\pi\log\mu}.\]
The first formula in each line is the UNNS Research Collective's
proposal.

\emph{Proof.} In \((G)\) at \(\theta^\ast\),
\(\theta^\ast_pm_p\ge-[-m_p]_+\); so \(\log\rho\ge H-a\Lambda\) for
every \(\pi\). Let \(\pi_M(\ell)\) be the term of \(\ell\) at
\(\theta^M\) divided by \(L_M\). The minimum over the orthant gives
\(a\log p\,L_M\,m_p(\pi_M)=\partial_p\Phi_M\ge0\), with equality if
\(\theta^M_p>0\). So \(\pi_M\) is resource-neutral and, by equality in
\((G)\), \(H(\pi_M)-a\,\mathbb E\log\mu=\log L_M\to\log\rho\). For
\(\sigma^\ast\): \(\rho_A(\sigma^\ast)=1\) gives
\(H\le\sigma^\ast\Lambda\) always, and neutral \(\pi_M\) at
\(a=\sigma^\ast\) have \(H-\sigma^\ast\mathbb E\log\mu\to0\) with
\(\mathbb E\log\mu\) bounded below by a positive constant. \(\square\)

\textbf{Corollary 1 (depletion is never optimal).} For every \(\pi\),
\[\log\rho-\big[H(\pi)-a\Lambda(\pi)\big]\ \ge\ a\sum_{p:\,m_p<0}(1-\theta^\ast_p)\,|m_p(\pi)|\log p,\]
which is positive as soon as some drift is negative. \emph{Proof:}
\((G)\) at \(\theta^\ast\) and \(\theta^\ast_p<1\). \(\square\)

\textbf{Depends on:} ANS2N-T02, ANS2-T16. \textbf{Used by:} ANS2N-T04a,
ANS2N-T04b. \textbf{Classical inputs:} Gibbs inequality.

\hypertarget{ans2n-t03-optimal-law-and-attainment-criterion}{%
\subsection{ANS2N-T03: Optimal law and attainment
criterion}\label{ans2n-t03-optimal-law-and-attainment-criterion}}

\textbf{Proof status:} PROVED. \textbf{Review:} UNNS Research Collective
review: valid, with a recorded qualification. Convergence valid;
attainment is open.

\textbf{Review record:} Release review records. \textbf{Release and
computation:} Archived release.

\textbf{Source:} ANS-II-ND v001, lines 49-51. \textbf{Scope:} Extended
grammar with torsion multipliers. \textbf{Literature priority:}
Unreviewed.

\textbf{Release table summary} (ANS-II/ND/v001): Optimal finite laws
converge to one probability; criterion and a sufficient condition for
attainment

\textbf{Statement as recorded:}

\textbf{Theorem 3.} Let
\(\nu^\ast(\ell)=\mu(\ell)^{-a}\prod p^{a\theta^\ast_pX_p(\ell)}/\rho\).
Then \(\nu^\ast\) is a probability, \(\pi_M\to\nu^\ast\) in total
variation, and \(\mathbb E_{\nu^\ast}\log\mu<\infty\). There is a
probability \(\pi\) on the full alphabet, with
\(\mathbb E_\pi\log\mu<\infty\), that is resource-neutral and attains
\(H(\pi)-a\,\mathbb E_\pi\log\mu=\log\rho\) if and only if
\[m_p(\nu^\ast)\ge0\ \text{ and }\ \theta^\ast_p\,m_p(\nu^\ast)=0\quad\text{for all }p;\qquad(\mathrm N)\]
in that case \(\pi=\nu^\ast\). Condition (N) holds if
\(\sup_p\theta^\ast_p<(a-1)/(2a)\).

\textbf{Depends on:} ANS2N-T02. \textbf{Used by:} ANS2N-T07, ANS2N-T08.
\textbf{Classical inputs:} Scheffe, Gibbs.

\hypertarget{ans2n-t04a-counting-limits-for-a-and-c}{%
\subsection{ANS2N-T04a: Counting limits for A and
C}\label{ans2n-t04a-counting-limits-for-a-and-c}}

\textbf{Proof status:} PROVED. \textbf{Review:} UNNS Research Collective
review: valid.

\textbf{Review record:} Release review records. \textbf{Release and
computation:} Archived release.

\textbf{Source:} ANS-II-ND v001, lines 58-60. \textbf{Scope:} Extended
grammar with torsion multipliers. \textbf{Literature priority:}
Unreviewed.

\textbf{Release table summary} (ANS-II/ND/v001): log A / log X and log C
/ log X tend to the abscissa

\textbf{Statement as recorded:}

\textbf{Theorem 4.} Let all torsion multipliers exceed \(1\). For every
state \(x=a_\varphi n\), with \(N_x\) the number of histories from
\(x\), \(A_\varphi\) the number of minimal histories and \(C_\varphi\)
the all-source count, each with endpoint at most \(X\),
\[\lim_{X\to\infty}\frac{\log N_x(X)}{\log X}=\lim\frac{\log A_\varphi(X)}{\log X}=\lim\frac{\log C_\varphi(X)}{\log X}=\sigma^\ast,\]
for any signs of the drifts.

\textbf{Depends on:} ANS2N-T01, ANS2-T16. \textbf{Classical inputs:}
Doob.

\hypertarget{ans2n-t04b-fixed-source-counting-exponent}{%
\subsection{ANS2N-T04b: Fixed-source counting
exponent}\label{ans2n-t04b-fixed-source-counting-exponent}}

\textbf{Proof status:} PROVED. \textbf{Review:} UNNS Research Collective
review: valid.

\textbf{Review record:} Release review records. \textbf{Release and
computation:} Archived release.

\textbf{Source:} ANS-II-ND v001, lines 58-60. \textbf{Scope:} Extended
grammar with torsion multipliers. \textbf{Literature priority:}
Unreviewed.

\textbf{Release table summary} (ANS-II/ND/v001): Every fixed-source
count has the same limit

\textbf{Statement as recorded:}

\textbf{Theorem 4.} Let all torsion multipliers exceed \(1\). For every
state \(x=a_\varphi n\), with \(N_x\) the number of histories from
\(x\), \(A_\varphi\) the number of minimal histories and \(C_\varphi\)
the all-source count, each with endpoint at most \(X\),
\[\lim_{X\to\infty}\frac{\log N_x(X)}{\log X}=\lim\frac{\log A_\varphi(X)}{\log X}=\lim\frac{\log C_\varphi(X)}{\log X}=\sigma^\ast,\]
for any signs of the drifts.

\textbf{Depends on:} ANS2N-T01, ANS2-T16, ANS2-L1, ANS-T24a.
\textbf{Classical inputs:} Doob.

\hypertarget{ans2n-t05-optimal-weights-zero-or-infinitely-supported}{%
\subsection{ANS2N-T05: Optimal weights zero or infinitely
supported}\label{ans2n-t05-optimal-weights-zero-or-infinitely-supported}}

\textbf{Proof status:} PROVED. \textbf{Review:} UNNS Research Collective
review: valid.

\textbf{Review record:} Release review records. \textbf{Release and
computation:} Archived release.

\textbf{Source:} ANS-II-ND v001, lines 73-73. \textbf{Scope:} Extended
grammar with torsion multipliers. \textbf{Literature priority:}
Unreviewed.

\textbf{Release table summary} (ANS-II/ND/v001): Optimal weights are
zero or supported on infinitely many primes

\textbf{Statement as recorded:}

\textbf{Theorem 5.} Let \(a>1\). If all drifts \(\mu_p(a)\) under the
unweighted law \(\nu_a\) are non-negative, then \(\theta^\ast=0\). If
one is negative, then \(\theta^\ast_p>0\) for infinitely many primes
\(p\).

\textbf{Depends on:} ANS2N-T02, ANS2-T16. \textbf{Used by:} ANS2N-T11.
\textbf{Classical inputs:} Dirichlet theorem on primes.

\hypertarget{ans2n-t06-certified-interval-for-2sqrt3}{%
\subsection{ANS2N-T06: Certified interval for
2/sqrt3}\label{ans2n-t06-certified-interval-for-2sqrt3}}

\textbf{Proof status:} CERTIFIED. \textbf{Review:} UNNS Research
Collective review: valid, with a recorded qualification. Method valid
and stored result reviewed; no independent rerun.

\textbf{Review record:} Release review records. \textbf{Release and
computation:} Archived release; certificate not independently rerun.

\textbf{Source:} ANS-II-ND v001, lines 90-91. \textbf{Scope:} Extended
grammar with torsion multipliers. \textbf{Literature priority:}
Unreviewed.

\textbf{Release table summary} (ANS-II/ND/v001): 2.06055 \textless=
abscissa of 2/sqrt3 \textless= 2.06072

\textbf{Statement as recorded:}

\[2.06055\le\sigma^\ast\le2.06072 .\] \emph{Method.} The inequalities of
ANS-II v002 §7 are valid for any weights. The weights are found in
floating point; the two certificates are then evaluated with MPFR
arithmetic at 80 bits, every operation rounded in the safe direction
(\texttt{certify.py}), at \(a=2.06055+2^{-40}\) (lower certificate
\(\ge1\), margin \(5.8\cdot10^{-6}\)) and \(a=2.06072-2^{-40}\) (upper
certificate \(<1\), margin \(1.3\cdot10^{-6}\)). The statement relies on
the correct rounding of MPFR and on the program.

\textbf{Depends on:} ANS2-C16. \textbf{Classical inputs:} MPFR directed
rounding.

\textbf{Relations:} ANS2N-T06 certifies ANS2-C16.

\hypertarget{ans2n-t07-euler-summability-criterion-hypothesis-superseded-by-t10}{%
\subsection{ANS2N-T07: Euler-summability criterion (hypothesis
superseded by
T10)}\label{ans2n-t07-euler-summability-criterion-hypothesis-superseded-by-t10}}

\textbf{Proof status:} PROVED. \textbf{Review:} UNNS Research Collective
review: valid, with a recorded qualification. Valid; its hypothesis is
superseded by ANS2N-T10.

\textbf{Review record:} Release review records. \textbf{Release and
computation:} Archived release.

\textbf{Source:} ANS-II-ND v002, lines 13-13. \textbf{Scope:} Extended
grammar with torsion multipliers. \textbf{Literature priority:}
Unreviewed.

\textbf{Release table summary} (ANS-II/ND/v002): Attainment if all
weights are below 1/2 and sum p\^{}(-a(1-2 theta\_p)) converges (and its
limsup corollary)

\textbf{Statement as recorded:}

\textbf{Theorem 7.} If \(\theta^\ast_p<\tfrac12\) for all \(p\) and
\(\sum_pp^{-c_p}<\infty\), where \(c_p=a(1-2\theta^\ast_p)\), then
\(\nu^\ast\) attains.

\textbf{Depends on:} ANS2N-T03.

\textbf{Relations:} ANS2N-T10 supersedes the hypothesis of ANS2N-T07.

\hypertarget{ans2n-t08-attainment-iff-uniform-integrability}{%
\subsection{ANS2N-T08: Attainment iff uniform
integrability}\label{ans2n-t08-attainment-iff-uniform-integrability}}

\textbf{Proof status:} PROVED. \textbf{Review:} UNNS Research Collective
review: valid.

\textbf{Review record:} Release review records. \textbf{Release and
computation:} Archived release.

\textbf{Source:} ANS-II-ND v002, lines 18-18. \textbf{Scope:} Extended
grammar with torsion multipliers. \textbf{Literature priority:}
Unreviewed.

\textbf{Release table summary} (ANS-II/ND/v002): Attainment iff log mu
is uniformly integrable under the finite optima

\textbf{Statement as recorded:}

\textbf{Theorem 8.} \(\nu^\ast\) attains if and only if \(\log\mu\) is
uniformly integrable under the family \((\pi_M)\); equivalently, if and
only if \(\mathbb E_{\pi_M}\log\mu\to E\).

\textbf{Depends on:} ANS2N-T03. \textbf{Used by:} ANS2N-T13, ANS2N-T14.

\hypertarget{ans2n-t09-variational-inequality-at-the-candidate-law}{%
\subsection{ANS2N-T09: Variational inequality at the candidate
law}\label{ans2n-t09-variational-inequality-at-the-candidate-law}}

\textbf{Proof status:} PROVED. \textbf{Review:} UNNS Research Collective
review: valid, with a recorded qualification. Valid; a proof
clarification is recorded in ANS-II-ND v003.

\textbf{Review record:} Release review records. \textbf{Release and
computation:} Archived release.

\textbf{Source:} ANS-II-ND v002, lines 25-26. \textbf{Scope:} Extended
grammar with torsion multipliers. \textbf{Literature priority:}
Unreviewed.

\textbf{Release table summary} (ANS-II/ND/v002): sum (w\_p - theta\_p)
m\_p log p \textgreater= 0 for every weighting w with finite Phi(w)

\textbf{Statement as recorded:}

\textbf{Theorem 9.} For every \(w\in[0,\infty)^{\mathbb P}\) with
\(\Phi(w)<\infty\), \[\sum_p(w_p-\theta^\ast_p)\,m_p\log p\ \ge\ 0 .\]

\textbf{Depends on:} ANS2N-T02. \textbf{Used by:} ANS2N-T10.

\hypertarget{ans2n-t10-depletion-surplus-inequality}{%
\subsection{ANS2N-T10: Depletion-surplus
inequality}\label{ans2n-t10-depletion-surplus-inequality}}

\textbf{Proof status:} PROVED. \textbf{Review:} UNNS Research Collective
review: valid.

\textbf{Review record:} Release review records. \textbf{Release and
computation:} Archived release.

\textbf{Source:} ANS-II-ND v002, lines 31-33, 38-42. \textbf{Scope:}
Extended grammar with torsion multipliers. \textbf{Literature priority:}
Unreviewed.

\textbf{Release table summary} (ANS-II/ND/v002): Depletion-surplus
inequality; neutrality of the candidate is equivalent to attainment;
attainment if weights are below 1/2 at depleting primes

\textbf{Statement as recorded:}

\textbf{Theorem 10.}
\[\sum_p\Big(\tfrac12-\theta^\ast_p\Big)m_p^-\log p\ +\ \sum_p\theta^\ast_p\,m_p^+\log p\ \le\ 0 .\qquad(\mathrm D)\]

\begin{enumerate}
\def\labelenumi{\arabic{enumi}.}
\tightlist
\item
  \emph{Neutrality is enough.} \(\nu^\ast\) attains if and only if
  \(m_p\ge0\) for all \(p\). (If every \(m_p^-=0\), (D) forces
  \(\theta^\ast_pm_p^+=0\).)
\item
  \emph{Attainment criterion.} \(\nu^\ast\) attains if
  \(\theta^\ast_p<\frac12\) at every prime with \(m_p<0\). In particular
  \textbf{it attains if \(\theta^\ast_p<\frac12\) for all \(p\)}. No
  summability, decay or uniform margin is needed; Theorem 7 and the
  criterion of v001 Theorem 3 are special cases.
\item
  \emph{What failure would require.} If \(\nu^\ast\) does not attain,
  there is a prime \(s\) with \(m_s<0\) and \(\theta^\ast_s\ge\frac12\),
  and \(\Phi(\theta^\ast+\epsilon e_s)=\infty\) for every
  \(\epsilon>0\). (The last assertion: if it were finite, Theorem 9
  would give \(\epsilon\,m_s\ge0\).) Moreover
  \(H(\nu^\ast)-aE=\log\rho-a\Gamma\) and
  \[\tfrac12\sum_pm_p^-\log p\ \le\ -\Gamma\ \le\ \sum_p\theta^\ast_pm_p^-\log p .\]
\item
  If \(\theta^\ast_p\le\frac12\) for all \(p\), depletion can occur only
  at primes with \(\theta^\ast_p=\frac12\), and \(m_p=0\) at every prime
  with \(0<\theta^\ast_p<\frac12\).
\end{enumerate}

\textbf{Depends on:} ANS2N-T09. \textbf{Used by:} ANS2N-C12, ANS2N-T11,
ANS2N-T14.

\textbf{Relations:} ANS2N-T10 supersedes the hypothesis of ANS2N-T07;
ANS2N-T14 contains, by a second argument, part of ANS2N-T10.

\hypertarget{ans2n-t11-weights-below-one-half-cannot-tend-to-zero}{%
\subsection{ANS2N-T11: Weights below one half cannot tend to
zero}\label{ans2n-t11-weights-below-one-half-cannot-tend-to-zero}}

\textbf{Proof status:} PROVED. \textbf{Review:} UNNS Research Collective
review: valid.

\textbf{Review record:} Release review records. \textbf{Release and
computation:} Archived release.

\textbf{Source:} ANS-II-ND v002, lines 45-49. \textbf{Scope:} Extended
grammar with torsion multipliers. \textbf{Literature priority:}
Unreviewed.

\textbf{Release table summary} (ANS-II/ND/v002): Under negative drift,
``all weights below 1/2'' and ``weights tend to zero'' cannot both hold;
limsup bound 2/(3L) times the supremum under T07's hypotheses

\textbf{Statement as recorded:}

Let \(L\) be an admissible Linnik exponent: the least prime
\(q\equiv b\pmod d\), \((b,d)=1\), is at most \(C_Ld^L\). Xylouris (Acta
Arith. 150 (2011)) gives \(L=5.18\).

\textbf{Theorem 11.} Suppose that some unweighted drift \(\mu_p(a)\) is
negative, that \(\theta^\ast_p<\frac12\) for all \(p\), and that
\(\sum_pp^{-c_p}<\infty\). Then
\[\limsup_{q\to\infty}\theta^\ast_q\ \ge\ \frac{2}{3L}\,\sup_p\theta^\ast_p\ >0 .\]
Consequently, when some drift is negative, \textbf{the two statements
``\(\theta^\ast_p<\frac12\) for all \(p\)'' and
``\(\theta^\ast_p\to0\)'' cannot both hold.}

\textbf{Depends on:} ANS2N-T10, ANS2N-T05. \textbf{Used by:} ANS2N-C12.
\textbf{Classical inputs:} Linnik (Xylouris exponent).

\hypertarget{ans2n-t12a-weight-above-one-half-forces-weight-on-divisors-of-c-pv---1}{%
\subsection{ANS2N-T12a: Weight above one half forces weight on divisors
of C p\^{}V -
1}\label{ans2n-t12a-weight-above-one-half-forces-weight-on-divisors-of-c-pv---1}}

\textbf{Proof status:} PROVED. \textbf{Review:} UNNS Research Collective
review: valid.

\textbf{Review record:} UNNS Research Collective review, 27 November
2026 (PROVENANCE/REVIEWS). \textbf{Release and computation:} Archived
release.

\textbf{Source:} ANS-II-ND v003, lines 11-12, 16-16. \textbf{Scope:}
Extended grammar with torsion multipliers. \textbf{Literature priority:}
Unreviewed.

\textbf{Release table summary} (ANS-II/ND/v003): Weight above 1/2 at p
forces weight on the divisors of C p\^{}V - 1

\textbf{Statement as recorded:}

\textbf{Lemma.} Let \(w\) be feasible, \(p\) a prime, \(C\ge1\) an
integer and \(N_V=Cp^V-1\). Then
\[\sum_{q\mid N_V}w_q\,v_q(N_V)\log q\;-\;(2w_p-1)\,V\log p\;\longrightarrow\;+\infty\qquad(V\to\infty).\]

So a weight above \(\frac12\) at \(p\) forces the prime divisors of
\(Cp^V-1\) to carry weight, for every \(C\) and every large \(V\).

\textbf{Depends on:} no earlier result cited. \textbf{Used by:}
ANS2N-T12.

\hypertarget{ans2n-t12-tail-weight-constraint}{%
\subsection{ANS2N-T12: Tail-weight
constraint}\label{ans2n-t12-tail-weight-constraint}}

\textbf{Proof status:} PROVED. \textbf{Review:} UNNS Research Collective
review: valid.

\textbf{Review record:} UNNS Research Collective review, 27 November
2026 (PROVENANCE/REVIEWS). \textbf{Release and computation:} Archived
release.

\textbf{Source:} ANS-II-ND v003, lines 18-18. \textbf{Scope:} Extended
grammar with torsion multipliers. \textbf{Literature priority:}
Unreviewed.

\textbf{Release table summary} (ANS-II/ND/v003): Tail-weight constraint:
limsup w\_p \textgreater= 2 sup w\_p - 1 for every feasible weighting

\textbf{Statement as recorded:}

\textbf{Theorem 12.} For feasible \(w\), with \(T=\sup_pw_p\) and
\(T_\infty=\limsup_pw_p\): \(T_\infty\ge2T-1\).

\textbf{Depends on:} ANS2N-T12a. \textbf{Used by:} ANS2N-C12.

\hypertarget{ans2n-c12-decaying-weights-have-maximum-exactly-one-half}{%
\subsection{ANS2N-C12: Decaying weights have maximum exactly one
half}\label{ans2n-c12-decaying-weights-have-maximum-exactly-one-half}}

\textbf{Proof status:} PROVED. \textbf{Review:} UNNS Research Collective
review: valid, with a recorded qualification. Valid under its
hypotheses.

\textbf{Review record:} UNNS Research Collective review, 27 November
2026 (PROVENANCE/REVIEWS). \textbf{Release and computation:} Archived
release.

\textbf{Source:} ANS-II-ND v003, lines 22-22. \textbf{Scope:} Extended
grammar with torsion multipliers. \textbf{Literature priority:}
Unreviewed.

\textbf{Release table summary} (ANS-II/ND/v003): Negative drift and
decaying weights force a maximum of exactly 1/2

\textbf{Statement as recorded:}

\textbf{Corollary (UNNS Research Collective's reduction, valid).} If
some drift is negative and \(\theta^\ast_p\to0\), then
\(\max_p\theta^\ast_p=\frac12\) exactly: \(T_\infty=0\) gives
\(T\le\frac12\), and v002 Theorem 11 excludes \(\theta^\ast_p<\frac12\)
for all \(p\). In addition, the set \(F=\{p:\theta^\ast_p=\frac12\}\) is
finite and non-empty, depletion can occur only on \(F\) (v002 Theorem
10, corollary 4), and \(\Phi(\theta^\ast+\epsilon e_s)=\infty\) for
every \(s\in F\), \(\epsilon>0\) (Theorem 12): these primes lie on the
boundary of the domain of \(\Phi\).

\textbf{Depends on:} ANS2N-T12, ANS2N-T11, ANS2N-T10.

\hypertarget{ans2n-t13-logarithmic-moment-defect-bounds}{%
\subsection{ANS2N-T13: Logarithmic moment defect
bounds}\label{ans2n-t13-logarithmic-moment-defect-bounds}}

\textbf{Proof status:} PROVED. \textbf{Review:} UNNS Research Collective
review: valid.

\textbf{Review record:} UNNS Research Collective review, 27 November
2026 (PROVENANCE/REVIEWS). \textbf{Release and computation:} Archived
release.

\textbf{Source:} ANS-II-ND v003, lines 25-25. \textbf{Scope:} Extended
grammar with torsion multipliers. \textbf{Literature priority:}
Unreviewed.

\textbf{Release table summary} (ANS-II/ND/v003): Moment defect: -Gamma
\textless= delta \textless= -a Gamma/(a-1)

\textbf{Statement as recorded:}

\textbf{Theorem 13.} If \(E_{M_j}-E\to\delta\), then
\(-\Gamma\le\delta\le-\dfrac{a}{a-1}\Gamma\).

\textbf{Depends on:} ANS2N-T08. \textbf{Used by:} ANS2N-T14.

\hypertarget{ans2n-t14-location-of-the-defect}{%
\subsection{ANS2N-T14: Location of the
defect}\label{ans2n-t14-location-of-the-defect}}

\textbf{Proof status:} PROVED. \textbf{Review:} UNNS Research Collective
review: valid, with a recorded qualification. Proof-writing
clarification: in the final Fatou step use a limit superior for the
drift difference, or a further convergent subsequence. The theorem is
unaffected; the release file is not edited.

\textbf{Review record:} UNNS Research Collective review, 27 November
2026 (PROVENANCE/REVIEWS). \textbf{Release and computation:} Archived
release.

\textbf{Source:} ANS-II-ND v003, lines 31-35, 45-48. \textbf{Scope:}
Extended grammar with torsion multipliers. \textbf{Literature priority:}
Unreviewed.

\textbf{Release table summary} (ANS-II/ND/v003): D/2 \textless= omega
\textless= -Gamma \textless= delta \textless= (-a Gamma - omega)/(a-1);
attainment iff E v\_p(N+1) converges to its candidate value for every p;
defect D/2 at the threshold

\textbf{Statement as recorded:}

\[\mathbb E_{\pi_{M_j}}v_p(N+1)-\mathbb E_{\nu^\ast}v_p(N+1)\to\kappa_p .\]
(Torsion letters count as \(v_p(N+1)=0\). Boundedness of \(E_M\) gives
such a subsequence.) By Fatou's lemma \(\kappa_p\ge0\). Put
\(\omega=\sum_p\kappa_p\log p\); it is the part of \(\delta\) carried by
divisibility of \(N+1\) at fixed primes.

\textbf{Theorem 14.}
\[\frac D2\ \le\ \omega\ \le\ -\Gamma\ \le\ \delta\ \le\ \frac{-a\Gamma-\omega}{a-1}.\]

\begin{enumerate}
\def\labelenumi{\arabic{enumi}.}
\tightlist
\item
  \emph{Prime-by-prime criterion.} \(\nu^\ast\) attains \textbf{if and
  only if} \(\mathbb E_{\pi_M}v_p(N+1)\to\mathbb E_{\nu^\ast}v_p(N+1)\)
  for every prime \(p\). (If it attains, \(\delta=0\) by v002 Theorem 8,
  so \(\omega=0\). If all \(\kappa_p=0\) then \(D=0\).) Failure of
  attainment is therefore visible at single fixed primes: at every
  depleted prime.
\item
  \(D/2\le\omega\le-\Gamma\) contains v002 Theorem 10, by a second
  argument.
\item
  \(\sum_{p\ \text{not depleted}}\kappa_p\log p\le\sum_{p\ \text{depleted}}(\theta^\ast_p-\tfrac12)\,m_p^-\log p\).
\item
  \emph{Exact defect at the threshold.} If \(\theta^\ast_p=\frac12\) at
  every depleted prime, then \(-\Gamma\le D/2\), so
  \(\omega=-\Gamma=D/2\) and \(\delta=D/2\) for every limit point:
  \(E_M\to E+D/2\). This applies to the exceptional decay case of
  section 1.
\end{enumerate}

\textbf{Depends on:} ANS2N-T13, ANS2N-T08, ANS2N-T10.

\textbf{Relations:} ANS2N-T14 contains, by a second argument, part of
ANS2N-T10.

\hypertarget{ans2n-t15-average-weight-at-most-1---1a}{%
\subsection{ANS2N-T15: Average weight at most 1 -
1/a}\label{ans2n-t15-average-weight-at-most-1---1a}}

\textbf{Proof status:} PROVED. \textbf{Review:} UNNS Research Collective
review: valid.

\textbf{Review record:} UNNS Research Collective review, 27 November
2026 (PROVENANCE/REVIEWS). \textbf{Release and computation:} Archived
release.

\textbf{Source:} ANS-II-ND v003, lines 51-55. \textbf{Scope:} Extended
grammar with torsion multipliers. \textbf{Literature priority:}
Unreviewed.

\textbf{Release table summary} (ANS-II/ND/v003): limsup eps sum w\_p log
p / p\^{}(1+eps) \textless= 1 - 1/a for every feasible weighting

\textbf{Statement as recorded:}

\textbf{Theorem 15.} (i) For \(1<a'<a\), \(s=1-a'/a\), with
\(B,E_{a'},\mu_p(a')\) the characteristic sum, mean logarithm and drifts
of the unweighted law at \(a'\), and any feasible \(w\):
\[\sum_pw_p\,\mu_p(a')\log p\ \le\ s\,E_{a'}-\frac1a+\frac{\Phi(w)}{a\,B(a')}.\]
(ii) Consequently
\[\limsup_{\epsilon\downarrow0}\ \epsilon\sum_pw_p\,\frac{\log p}{p^{1+\epsilon}}\ \le\ 1-\frac1a .\]
In particular this holds for \(w=\theta^\ast\).

\textbf{Depends on:} ANS2N-T02.

\end{document}
